% La marche vers l’indépendance — Histoire Tle Tech — Cours signé OnBuch+
% Programme officiel MINESEC. Compiler avec : tectonic N06-marche-vers-independance.tex
\newcommand{\DOCMATIERE}{Histoire}
\newcommand{\DOCNIVEAU}{Tle Tech}
\newcommand{\DOCMODULE}{Histoire – classe de Terminale technique}
\newcommand{\DOCLECON}{N06}
\newcommand{\DOCTITRE}{La marche vers l’indépendance}
% OnBuch+ — préambule « cours signés » ENRICHI (compilé avec tectonic / XeLaTeX)
% Système visuel « Pop » d'OnBuch+ : fond crème (jamais blanc), bordures encre
% épaisses, ombres « dures » décalées, titres Archivo Black en capitales.
% Version MATHÉMATIQUES : graphes (pgfplots/tikz), boîtes pédagogiques
% (définition, propriété, méthode, attention, le savais-tu, exercice résolu,
% exercice, TP, bilan), page de garde et signature OnBuch+ ancrée en bas.
%
% Macros attendues dans main.tex AVANT \input{preamble} :
%   \DOCMATIERE  \DOCNIVEAU  \DOCMODULE  \DOCLECON (numéro)  \DOCTITRE
\documentclass[11pt]{article}

\usepackage{fontspec}
\usepackage[french]{babel}
\usepackage[a4paper,margin=2cm,top=3.4cm,bottom=2.8cm,headheight=1.4cm,footskip=1.3cm]{geometry}
\usepackage{amsmath,amssymb,amsfonts}
\usepackage{mathtools} % \xmapsto, \coloneqq... souvent utilisés par les modèles
\usepackage{cancel} % simplifications barrées (bilans de forces)
% \abs{z} (module, valeur absolue) et \norm{u} (norme), utilisés par les modèles
\providecommand{\abs}{}\let\abs\relax\DeclarePairedDelimiter{\abs}{\lvert}{\rvert}
\providecommand{\norm}{}\let\norm\relax\DeclarePairedDelimiter{\norm}{\lVert}{\rVert}
\usepackage[table]{xcolor} % [table] : fournit aussi \rowcolors (alternance de lignes)
\usepackage{pagecolor}
\usepackage{tikz}
\usetikzlibrary{arrows.meta,positioning,calc,shapes.geometric,shapes.misc,shapes.arrows,decorations.pathmorphing,decorations.pathreplacing,decorations.markings,patterns,shadows,mindmap,trees,fit,backgrounds,matrix,intersections,angles,quotes}
\usepackage{pgfplots}
\usepackage[version=4]{mhchem} % équations nucléaires \ce{^{235}_{92}U}
\usepackage[european,siunitx]{circuitikz} % schémas électriques
\pgfplotsset{compat=1.18}
\usepgfplotslibrary{fillbetween,groupplots} % aires entre courbes (calcul intégral)
\usepackage{enumitem}
\usepackage{fancyhdr}
% [most] charge toutes les bibliothèques tcolorbox (listings, minted, raster,
% documentation...) et gonfle inutilement la mémoire TeX sur un document long
% (~300 boîtes) : on ne charge que ce qui est réellement utilisé.
\usepackage[skins,breakable]{tcolorbox}
\usepackage{graphicx}
\usepackage{colortbl}
\usepackage{booktabs}
\usepackage{tabularx}
\usepackage{multirow}
\usepackage{diagbox} % case d'en-tête diagonale des tableaux à double entrée
% Colonnes extensibles centrée / gauche / droite (C, L, R), souvent utilisées par les modèles
\newcolumntype{C}{>{\centering\arraybackslash}X}
\newcolumntype{L}{>{\raggedright\arraybackslash}X}
\newcolumntype{R}{>{\raggedleft\arraybackslash}X}
\usepackage{array}
\usepackage{multicol}
\usepackage{siunitx}
% parse-numbers=false : accepte aussi \SI{2,5 \times 10^{-3}}{...}
\sisetup{locale=FR,output-decimal-marker={,},group-minimum-digits=4,parse-numbers=false}
\DeclareSIUnit\ohm{\ensuremath{\symup{\Omega}}} % U+2126 absent de la police
\DeclareSIUnit\division{div} % réglages d'oscilloscope (ms/div, V/div)
\DeclareSIUnit\tr{tr} % tours (vitesse de rotation, ex. \SI{1500}{\tr\per\minute})
\DeclareSIUnit\molar{M} % concentration molaire (ex. \SI{3}{\molar}, \SI{1}{\micro\molar})
\DeclareSIUnit\an{an} % année (le modèle confond parfois avec \year, un compteur TeX) — ex. \SI{5}{\kilo\meter\per\an}
\DeclareSIUnit\FCFA{FCFA} % franc CFA (ex. \SI{500}{\FCFA\per\kilo\gram})
\DeclareSIUnit\degreeBrix{\textdegree Brix} % teneur en sucre d'un sirop/jus (réfractométrie)
\DeclareSIUnit\month{mois} % le modèle utilise parfois \month, un compteur TeX, comme unité de durée
\DeclareSIUnit\microliter{\micro\liter} % le modèle écrit parfois \microliter en un seul token
\DeclareSIUnit\million{millions} % dénombrement (ex. \SI{15}{\million\per\mL} spermatozoïdes)
\usepackage{qrcode}
% \degree : commande fréquemment utilisée par les modèles pour un angle ou une
% température en dehors de \SI{}{} (non fournie par les paquets chargés).
\providecommand{\degree}{^{\circ}}
% \qed (fin de démonstration), utilisé par les modèles sans amsthm
\providecommand{\qed}{\hfill$\square$}
% \barre : double barre des valeurs interdites dans un tableau de variations (tkz-tab)
\providecommand{\barre}{\ensuremath{\|}}
% environnement proof (amsthm non chargé) utilisé par les modèles
\ifdefined\proof\else\newenvironment{proof}[1][Démonstration]{\par\noindent\textit{#1.}\ }{\qed\par}\fi
\usepackage{lastpage}
\usepackage{hyperref}
\hypersetup{hidelinks}

% ============================================================
% Palette « Pop » OnBuch+ — mêmes hex que la classe P (plus_ui.dart)
% ============================================================
\definecolor{poporange}{HTML}{F59321}
\definecolor{popdark}{HTML}{E07A0C}
\definecolor{popcream}{HTML}{FFF6E8}
\definecolor{popink}{HTML}{1C1714}
\definecolor{poppurple}{HTML}{7A5AE0}
\definecolor{popgreen}{HTML}{2E9E6A}
\definecolor{poppink}{HTML}{E0457B}
\definecolor{popblue}{HTML}{2F7DE1}
\definecolor{popgold}{HTML}{C98A12}
\definecolor{popmuted}{HTML}{8A7F76}
\definecolor{popline}{HTML}{EADFD2}
\definecolor{popgray}{HTML}{8A7F76}
\colorlet{popgrey}{popgray}
% Alias tolérants (le modèle nomme parfois les couleurs différemment)
\colorlet{popwhite}{white}
\colorlet{popblack}{popink}
\colorlet{popred}{poppink}
\colorlet{popyellow}{popgold}
% Teintes claires (fonds de tableaux/encadrés) — le modèle nomme parfois
% les couleurs avec un suffixe L (ex. popblueL) sans les avoir définies.
\colorlet{poporangeL}{poporange!20}
\colorlet{popdarkL}{popdark!20}
\colorlet{poppurpleL}{poppurple!20}
\colorlet{popgreenL}{popgreen!20}
\colorlet{poppinkL}{poppink!20}
\colorlet{popblueL}{popblue!20}
\colorlet{popgoldL}{popgold!20}
\colorlet{popredL}{popred!20}
\colorlet{popgrayL}{popgray!20}
% Teintes claires pour fonds de tableaux / zones
\colorlet{popblueL}{popblue!10!white}
\colorlet{poporangeL}{poporange!14!white}
\colorlet{popgreenL}{popgreen!12!white}
\colorlet{poppurpleL}{poppurple!12!white}
\colorlet{poppinkL}{poppink!12!white}
\colorlet{popgoldL}{popgold!14!white}

\pagecolor{popcream}

% ============================================================
% Typographie
% ============================================================
\defaultfontfeatures{Ligatures=TeX}
\setmainfont{PlusJakartaSans-Medium.ttf}[Path=../../../fonts/,BoldFont=PlusJakartaSans-Bold.ttf]
% Maths en sans-serif (Fira Math) avec chiffres et lettres droites de Plus
% Jakarta, pour que les formules se fondent dans le texte.
\usepackage[math-style=ISO,bold-style=ISO]{unicode-math}
\setmathfont{FiraMath-Regular.otf}[Path=../../../fonts/]
\setmathfont{PlusJakartaSans-Medium.ttf}[Path=../../../fonts/,range={up/num,up/{Latin,latin},\mathup}]
\usepackage{icomma} % virgule décimale sans espace en mode math
% Caractères mathématiques Unicode absents des polices de texte (Plus Jakarta,
% Archivo Black) : rendus par leur équivalent mathématique, en texte comme en
% formule — sinon ils s'affichent en carré vide (ex. « Arithmétique dans ℤ »).
\usepackage{newunicodechar}
\newunicodechar{ℝ}{\ensuremath{\mathbb{R}}}
\newunicodechar{ℤ}{\ensuremath{\mathbb{Z}}}
\newunicodechar{ℂ}{\ensuremath{\mathbb{C}}}
\newunicodechar{ℕ}{\ensuremath{\mathbb{N}}}
\newunicodechar{ℚ}{\ensuremath{\mathbb{Q}}}
\newunicodechar{𝔻}{\ensuremath{\mathbb{D}}}
\newunicodechar{α}{\ensuremath{\alpha}}
\newunicodechar{β}{\ensuremath{\beta}}
\newunicodechar{γ}{\ensuremath{\gamma}}
\newunicodechar{δ}{\ensuremath{\delta}}
\newunicodechar{ε}{\ensuremath{\varepsilon}}
\newunicodechar{θ}{\ensuremath{\theta}}
\newunicodechar{λ}{\ensuremath{\lambda}}
\newunicodechar{μ}{\ensuremath{\mu}}
\newunicodechar{σ}{\ensuremath{\sigma}}
\newunicodechar{τ}{\ensuremath{\tau}}
\newunicodechar{φ}{\ensuremath{\varphi}}
\newunicodechar{ψ}{\ensuremath{\psi}}
\newunicodechar{ω}{\ensuremath{\omega}}
\newunicodechar{Σ}{\ensuremath{\Sigma}}
% majuscules produites par \MakeUppercase dans les titres (α-aminés -> Α)
\newunicodechar{Α}{\ensuremath{\alpha}}
\newunicodechar{Β}{\ensuremath{\beta}}
\newunicodechar{Γ}{\ensuremath{\Gamma}}
\newunicodechar{Δ}{\ensuremath{\Delta}}
\newunicodechar{Ε}{\ensuremath{\varepsilon}}
\newunicodechar{Θ}{\ensuremath{\Theta}}
\newunicodechar{Λ}{\ensuremath{\Lambda}}
\newunicodechar{Μ}{\ensuremath{\mu}}
\newunicodechar{Τ}{\ensuremath{\tau}}
\newunicodechar{Φ}{\ensuremath{\Phi}}
\newunicodechar{Ψ}{\ensuremath{\Psi}}
\newunicodechar{Ω}{\ensuremath{\Omega}}
\newunicodechar{↦}{\ensuremath{\mapsto}}
\newunicodechar{∈}{\ensuremath{\in}}
\newunicodechar{∉}{\ensuremath{\notin}}
\newunicodechar{∧}{\ensuremath{\wedge}}
\newunicodechar{⇔}{\ensuremath{\Leftrightarrow}}
\newunicodechar{⟺}{\ensuremath{\Longleftrightarrow}}
\newunicodechar{⇒}{\ensuremath{\Rightarrow}}
\newunicodechar{⟹}{\ensuremath{\Longrightarrow}}
\newunicodechar{∘}{\ensuremath{\circ}}
\newunicodechar{∪}{\ensuremath{\cup}}
\newunicodechar{∩}{\ensuremath{\cap}}
\newunicodechar{≡}{\ensuremath{\equiv}}
\newunicodechar{⊂}{\ensuremath{\subset}}
\newunicodechar{⊥}{\ensuremath{\perp}}
\newunicodechar{∃}{\ensuremath{\exists}}
\newunicodechar{∀}{\ensuremath{\forall}}
\newunicodechar{∛}{\ensuremath{\sqrt[3]{\,}}}
\newunicodechar{∜}{\ensuremath{\sqrt[4]{\,}}}
\newunicodechar{⃗}{\ensuremath{{}^{\to}}}
\newunicodechar{ⁿ}{\ifmmode^{n}\else\textsuperscript{\ensuremath{n}}\fi}
\newunicodechar{ˣ}{\ifmmode^{x}\else\textsuperscript{\ensuremath{x}}\fi}
\newunicodechar{ᵇ}{\ifmmode^{b}\else\textsuperscript{\ensuremath{b}}\fi}
\newunicodechar{ʳ}{\ifmmode^{r}\else\textsuperscript{\ensuremath{r}}\fi}
\newunicodechar{ᵘ}{\ifmmode^{u}\else\textsuperscript{\ensuremath{u}}\fi}
\newunicodechar{ʸ}{\ifmmode^{y}\else\textsuperscript{\ensuremath{y}}\fi}
\newunicodechar{ᵃ}{\ifmmode^{a}\else\textsuperscript{\ensuremath{a}}\fi}
\newunicodechar{ᵉ}{\ifmmode^{e}\else\textsuperscript{\ensuremath{e}}\fi}
\newunicodechar{ᵏ}{\ifmmode^{k}\else\textsuperscript{\ensuremath{k}}\fi}
\newunicodechar{ᵖ}{\ifmmode^{p}\else\textsuperscript{\ensuremath{p}}\fi}
\newunicodechar{ᴺ}{\ifmmode^{N}\else\textsuperscript{\ensuremath{N}}\fi}
\newunicodechar{ᶜ}{\ifmmode^{c}\else\textsuperscript{\ensuremath{c}}\fi}
\newunicodechar{ʰ}{\ifmmode^{h}\else\textsuperscript{\ensuremath{h}}\fi}
\newunicodechar{ᵠ}{\ifmmode^{q}\else\textsuperscript{\ensuremath{q}}\fi}
\newunicodechar{ₙ}{\ifmmode_{n}\else\textsubscript{\ensuremath{n}}\fi}
\newunicodechar{ₐ}{\ifmmode_{a}\else\textsubscript{\ensuremath{a}}\fi}
\newunicodechar{ᵢ}{\ifmmode_{i}\else\textsubscript{\ensuremath{i}}\fi}
\newunicodechar{ᵣ}{\ifmmode_{r}\else\textsubscript{\ensuremath{r}}\fi}
\newunicodechar{ₖ}{\ifmmode_{k}\else\textsubscript{\ensuremath{k}}\fi}
\newunicodechar{ᵦ}{\ifmmode_{\beta}\else\textsubscript{\ensuremath{\beta}}\fi}
\newunicodechar{ₓ}{\ifmmode_{x}\else\textsubscript{\ensuremath{x}}\fi}
\newfontfamily\popdisplay{ArchivoBlack-Regular.ttf}[Path=../../../fonts/]
\newfontfamily\poplogo{SpaceGrotesk-Bold.ttf}[Path=../../../fonts/]
\newfontfamily\popsemi{PlusJakartaSans-SemiBold.ttf}[Path=../../../fonts/]
\newfontfamily\popxbold{PlusJakartaSans-ExtraBold.ttf}[Path=../../../fonts/]
\newfontfamily\popxboldls{PlusJakartaSans-ExtraBold.ttf}[Path=../../../fonts/,LetterSpace=3.0]

\setlist[enumerate]{leftmargin=1.7em,itemsep=5pt,topsep=4pt}
\setlist[itemize]{leftmargin=1.7em,itemsep=4pt,topsep=4pt,label={\color{poporange}\textbullet}}
\setlength{\parindent}{0pt}
\setlength{\parskip}{5pt}
\renewcommand{\arraystretch}{1.25}

% pgfplots : style Pop par défaut
\pgfplotsset{
  every axis/.append style={axis line style={popink,line width=1pt},tick style={popink},
    grid style={popline,line width=0.5pt},label style={font=\small},tick label style={font=\footnotesize}},
  popcurve/.style={poporange,line width=1.8pt,no markers,smooth},
}
% Même style disponible dans un \draw TikZ simple (hors pgfplots)
\tikzset{popcurve/.style={poporange,line width=1.8pt}}
\tikzset{popgrid/.style={popline,very thin}}
\colorlet{popcurve}{poporange} % le nom du style est parfois utilisé comme couleur

% ============================================================
% Pill / squiggle
% ============================================================
\newcommand{\poppill}[3]{%
  \tikz[baseline=-0.55ex]\node[draw=popink,line width=1.2pt,rounded corners=2.6pt,%
    fill=#2,inner xsep=6.5pt,inner ysep=3pt]{%
    \popxboldls\fontsize{7}{7}\selectfont\color{#3}\MakeUppercase{#1}};%
}
\newcommand{\popsquiggle}[1]{%
  \tikz[baseline=-0.4ex]{
    \draw[#1,line width=2.6pt,line cap=round]
      plot [smooth] coordinates {(0,0.09) (0.34,0.01) (0.68,0.21) (1.02,0.01) (1.36,0.10)};
  }%
}
% Mot-clé surligné (marqueur orange sous le texte)
\newcommand{\cle}[1]{{\popxbold\color{popdark}#1}}

% ============================================================
% En-tête / pied
% ============================================================
\pagestyle{fancy}
\fancyhf{}
\renewcommand{\headrulewidth}{0pt}
\renewcommand{\footrulewidth}{0pt}
\lhead{\raisebox{-0.15ex}{\poplogo\fontsize{13}{13}\selectfont\color{popink}OnBuch\color{poporange}+}}
\rhead{\poppill{\DOCMATIERE\ \textbullet\ \DOCNIVEAU\ \textbullet\ Leçon \DOCLECON}{white}{popink}}
\lfoot{\popxbold\fontsize{7.6}{7.6}\selectfont\color{popmuted}\MakeUppercase{Cours signé OnBuch+}\ \textbullet\ {\popsemi\DOCTITRE}}
\rfoot{\tikz[baseline=-0.6ex]\node[rounded corners=8.5pt,fill=popink,minimum height=17pt,minimum width=17pt,inner xsep=5pt,inner sep=0]{\popxbold\fontsize{8}{8}\selectfont\color{popcream}\thepage\,/\,\pageref*{LastPage}};}

% ============================================================
% Titres
% ============================================================
\newcommand{\chaptitre}[2]{%
  \begin{tcolorbox}[enhanced,colback=poporange,colframe=popink,boxrule=2.6pt,arc=11pt,
      drop shadow=popink,left=15pt,right=15pt,top=12pt,bottom=11pt,before skip=3pt,after skip=18pt]
    \poppill{#2}{popink}{popcream}\par\vspace{9pt}
    {\raggedright\hyphenpenalty=10000\exhyphenpenalty=10000\popdisplay\fontsize{22}{24}\selectfont\color{popcream}\MakeUppercase{#1}\par}
  \end{tcolorbox}
}
% Section numérotée : pastille orange avec le numéro + titre Archivo
\newcounter{popsec}
\newcommand{\coursec}[1]{%
  \stepcounter{popsec}\setcounter{popsub}{0}%
  \par\vspace{16pt}\noindent
  \begin{minipage}{\linewidth}
  \tikz[baseline=-0.7ex]\node[draw=popink,line width=1.6pt,fill=poporange,rounded corners=4pt,minimum size=22pt,inner sep=0]{\popdisplay\fontsize{12}{12}\selectfont\color{popcream}\thepopsec};%
  \hspace{8pt}{\popdisplay\fontsize{15}{17}\selectfont\color{popink}\MakeUppercase{#1}}\par
  \vspace{4pt}\noindent{\color{poporange}\rule{36pt}{3.6pt}}
  \end{minipage}\par\nopagebreak\vspace{8pt}
}
\newcounter{popsub}
\newcommand{\courssub}[1]{%
  \stepcounter{popsub}%
  \par\vspace{9pt}\noindent{\popxbold\fontsize{11.5}{13}\selectfont\color{popink}{\color{poporange}\thepopsec.\thepopsub}\hspace{6pt}#1}\par\nopagebreak\vspace{4pt}
}

% ============================================================
% Boîtes pédagogiques — même recette : fond blanc, bordure épaisse colorée,
% ombre dure encre, badge plein. Titre optionnel [#1].
% ============================================================
\tcbset{popbase/.style={breakable,enhanced,colback=white,boxrule=2.3pt,arc=9pt,
  shadow={3pt}{-3pt}{0pt}{popink},left=14pt,right=14pt,top=11pt,bottom=11pt,before skip=11pt,after skip=15pt,
  fontupper=\popsemi\fontsize{10.6}{14.5}\selectfont\color{popink},
  fontlower=\popsemi\fontsize{10.6}{14.5}\selectfont\color{popink}}}
% Coche dessinée (le glyphe ✓ n'existe pas dans les polices du document)
\newcommand{\popcheck}[1]{\tikz[baseline=-0.5ex]\draw[#1,line width=1.6pt,line cap=round,line join=round](-0.13,0.0)--(-0.04,-0.09)--(0.13,0.1);}
\providecommand{\coche}{\popcheck{popgreen}}
\providecommand{\resultat}[1]{\par\smallskip\noindent\fcolorbox{popgreen}{popgreenL}{\parbox{\dimexpr\linewidth-2\fboxsep-2\fboxrule\relax}{\textbf{Résultat.} #1}}\par\smallskip}
\newcommand{\popboxhead}[3]{\poppill{#1}{#2}{white}\ifx\relax#3\relax\else\hspace{6pt}{\popxbold\fontsize{10.6}{12}\selectfont\color{popink}#3}\fi\par\vspace{7pt}}

\newtcolorbox{aretenir}[1][]{popbase,colframe=popgreen,before upper={\popboxhead{À retenir}{popgreen}{#1}}}
\newtcolorbox{exemplebox}[1][]{popbase,colframe=poporange,before upper={\popboxhead{Exemple}{poporange}{#1}}}
\newtcolorbox{definition}[1][]{popbase,colframe=popblue,before upper={\popboxhead{Définition}{popblue}{#1}}}
\newtcolorbox{propriete}[1][]{popbase,colframe=poppurple,before upper={\popboxhead{Propriété}{poppurple}{#1}}}
\newtcolorbox{methode}[1][]{popbase,colframe=popgold,before upper={\popboxhead{Méthode}{popgold}{#1}}}
\newtcolorbox{attention}[1][]{popbase,colframe=poppink,before upper={\popboxhead{Attention}{poppink}{#1}}}
\newtcolorbox{experience}[1][]{popbase,colframe=popblue,colback=popblueL,before upper={\popboxhead{Expérience}{popblue}{#1}}}
% « Le savais-tu ? » : carte violette pleine (contexte, culture, Cameroun)
\newtcolorbox{savaistu}[1][]{popbase,colback=poppurple,colframe=popink,
  fontupper=\popsemi\fontsize{10.4}{14.2}\selectfont\color{white},
  before upper={\poppill{Le savais-tu\,?}{white}{poppurple}\ifx\relax#1\relax\else\hspace{6pt}{\popxbold\color{white}#1}\fi\par\vspace{7pt}}}
% Exercice résolu : énoncé en haut, \tcblower puis la solution sur fond crème
\newcounter{popexo}\newcounter{popexr}
\newtcolorbox{exoresolu}[1][]{popbase,colframe=poporange,colbacklower=popcream,
  segmentation style={popink,line width=1.2pt,dashed},
  before upper={\stepcounter{popexr}\popboxhead{Exercice résolu \thepopexr}{poporange}{#1}},
  before lower={\poppill{Solution}{popink}{popcream}\par\vspace{6pt}}}
% Exercice (non résolu en place) — la correction est dans la partie Corrigés
\newtcolorbox{exercice}[1][]{popbase,colframe=popink,
  before upper={\stepcounter{popexo}\popboxhead{Exercice \thepopexo}{popink}{#1}}}
% Corrigé
\newtcolorbox{corrige}[1][]{popbase,colframe=popgreen,colback=popgreenL,
  before upper={\popboxhead{Corrigé}{popgreen}{#1}}}
% Objectifs / prérequis / bilan
\newtcolorbox{objectifs}{popbase,colframe=popink,colback=poporangeL,
  before upper={\popboxhead{Objectifs de la leçon}{popink}{}}}
\newtcolorbox{prerequis}{popbase,colframe=popink,colback=popblueL,
  before upper={\popboxhead{Prérequis}{popblue}{}}}
\newtcolorbox{bilan}{popbase,colframe=popink,colback=white,boxrule=2.8pt,
  before upper={\popboxhead{Fiche bilan}{poporange}{L'essentiel à connaître pour l'examen}}}
% Figure encadrée avec légende
\newtcolorbox{popfigure}[1][]{enhanced,breakable=false,colback=white,colframe=popline,boxrule=1.4pt,arc=8pt,
  left=10pt,right=10pt,top=10pt,bottom=8pt,before skip=10pt,after skip=12pt,halign=center,
  lower separated=false,halign lower=center,
  fontlower=\popsemi\fontsize{9}{11.5}\selectfont\color{popmuted},
  #1}
\newcommand{\legende}[1]{\par\vspace{6pt}{\popsemi\fontsize{9}{11.5}\selectfont\color{popmuted}#1}}

% ============================================================
% Signature OnBuch+
% ============================================================
\newcommand{\poplogoinline}[1]{{\poplogo\fontsize{#1}{#1}\selectfont\color{popink}OnBuch\color{poporange}+}}
\newcommand{\signatureonbuch}{%
  \begin{tcolorbox}[enhanced,colback=popink,colframe=popink,boxrule=2.2pt,arc=10pt,
      drop shadow=poporange,left=14pt,right=14pt,top=10pt,bottom=10pt,before skip=0pt,after skip=0pt]
    \tikz[baseline=-0.7ex]\node[circle,fill=popgreen,draw=popcream,line width=1.3pt,minimum size=22pt,inner sep=0]{\popcheck{white}};%
    \hspace{9pt}\begin{minipage}[c]{0.62\linewidth}
      {\popxbold\fontsize{12}{13}\selectfont\color{popcream}Signé\ \textbullet\ {\poplogo OnBuch\color{poporange}+}}\par\vspace{2pt}
      {\raggedright\popsemi\fontsize{8}{10.5}\selectfont\color{popline}\MakeUppercase{Équipe pédagogique OnBuch+}\\\MakeUppercase{Conforme au programme officiel MINESEC}\par}
    \end{minipage}\hfill
    {\poplogo\fontsize{22}{22}\selectfont\color{popcream}OnBuch\color{poporange}+}
  \end{tcolorbox}%
}

% ============================================================
% Page de garde
% #1 titre · #2 matière · #3 niveau · #4 module · #5 n° leçon · #6 durée ·
% #7 description / objectifs (texte ou itemize)
% ============================================================
\newcommand{\pagegarde}[7]{%
  \thispagestyle{empty}
  \newgeometry{a4paper,margin=1.7cm,top=1.6cm,bottom=1.4cm}%
  \begin{tikzpicture}[remember picture,overlay]
    \fill[poporange!18] ([xshift=-3.2cm,yshift=-3.6cm]current page.north east) circle (4.6cm);
    \fill[poppurple!14] ([xshift=2.4cm,yshift=7.6cm]current page.south west) circle (3.8cm);
  \end{tikzpicture}%
  {\poplogo\fontsize{30}{30}\selectfont\color{popink}OnBuch\color{poporange}+}\hfill
  \raisebox{0.6ex}{\poppill{Cours signé \textbullet\ Premium}{popink}{popcream}}\par
  \vspace{1pt}\popsquiggle{poppurple}\par
  \vspace{8pt}
  \begin{tcolorbox}[enhanced,colback=poporange,colframe=popink,boxrule=2.8pt,arc=13pt,
      drop shadow=popink,left=18pt,right=18pt,top=12pt,bottom=13pt,after skip=10pt]
    \poppill{#2 \textbullet\ #3}{popink}{popcream}\hspace{5pt}\poppill{#4}{white}{popink}\par\vspace{8pt}
    {\popdisplay\fontsize{14}{14}\selectfont\color{popink}LEÇON #5}\par\vspace{4pt}
    {\raggedright\hyphenpenalty=10000\exhyphenpenalty=10000\popdisplay\fontsize{25}{27}\selectfont\color{popcream}\MakeUppercase{#1}\par}
    \vspace{8pt}
    {\popxbold\fontsize{9.5}{9.5}\selectfont\color{popink}\MakeUppercase{Durée conseillée : #6}}
  \end{tcolorbox}
  \begin{tcolorbox}[enhanced,colback=white,colframe=popink,boxrule=2.2pt,arc=10pt,
      drop shadow=popink,left=16pt,right=16pt,top=10pt,bottom=10pt,after skip=8pt]
    \poppill{Dans cette leçon}{poppurple}{white}\par\vspace{5pt}
    \popsemi\fontsize{9.3}{12.6}\selectfont\color{popink}#7
  \end{tcolorbox}
  \noindent\begin{minipage}[t]{0.487\linewidth}
    \begin{tcolorbox}[enhanced,colback=popgreen,colframe=popink,boxrule=2.2pt,arc=10pt,
        drop shadow=popink,left=11pt,right=11pt,top=10pt,bottom=10pt,height=3.1cm,valign=center]
      \tcbox[colback=white,colframe=popink,boxrule=1.3pt,arc=5pt,boxsep=0pt,nobeforeafter,
        left=3pt,right=3pt,top=3pt,bottom=3pt,valign=center]{\qrcode[height=1.9cm]{https://play.google.com/store/apps/details?id=com.luvvix.onbuch&hl=fr}}%
      \hspace{8pt}\begin{minipage}[c]{0.5\linewidth}
        {\popdisplay\fontsize{11}{12.5}\selectfont\color{white}\MakeUppercase{Télécharge OnBuch}}\par\vspace{4pt}
        {\raggedright\popsemi\fontsize{8.4}{11}\selectfont\color{white}Gratuit \textbullet\ Google Play\\Tuteur IA, annales, concours, résultats\par}
      \end{minipage}
    \end{tcolorbox}
  \end{minipage}\hfill
  \begin{minipage}[t]{0.487\linewidth}
    \begin{tcolorbox}[enhanced,colback=poppurple,colframe=popink,boxrule=2.2pt,arc=10pt,
        drop shadow=popink,left=11pt,right=11pt,top=10pt,bottom=10pt,height=3.1cm,valign=center]
      \tikz[baseline=-0.7ex]\node[circle,fill=white,draw=popink,line width=1.3pt,minimum size=1.6cm,inner sep=0]{\popdisplay\fontsize{24}{24}\selectfont\color{poppurple}+};%
      \hspace{8pt}\begin{minipage}[c]{0.6\linewidth}
        {\popdisplay\fontsize{11}{12.5}\selectfont\color{white}\MakeUppercase{Passe à OnBuch+}}\par\vspace{4pt}
        {\raggedright\popsemi\fontsize{8.4}{11}\selectfont\color{white}Tous les cours signés, Tuteur IA illimité, fiches PDF — onglet OnBuch+ de l'app\par}
      \end{minipage}
    \end{tcolorbox}
  \end{minipage}
  \vspace{8pt}
  \popcontenu
  \vfill
  \signatureonbuch
  \restoregeometry
  \newpage
}

% Rangée « Ce cours contient » (page de garde)
\newcommand{\poptile}[3]{% #1 couleur #2 gros texte #3 légende
  \begin{tcolorbox}[enhanced,colback=#1,colframe=popink,boxrule=2pt,arc=9pt,shadow={2.5pt}{-2.5pt}{0pt}{popink},
      left=6pt,right=6pt,top=9pt,bottom=9pt,halign=center,valign=center,height=1.75cm,width=\linewidth,nobeforeafter]
    {\popdisplay\fontsize{12.5}{13}\selectfont\color{white}\MakeUppercase{#2}}\par\vspace{3pt}
    {\centering\popsemi\fontsize{7.8}{9.5}\selectfont\color{white}#3\par}
  \end{tcolorbox}}
\newcommand{\popcontenu}{%
  \par\noindent{\popxboldls\fontsize{7.5}{7.5}\selectfont\color{popmuted}CE COURS SIGNÉ CONTIENT}\par\vspace{7pt}
  \noindent\begin{minipage}[t]{0.235\linewidth}\poptile{popblue}{Cours}{complet, illustré, pas à pas}\end{minipage}\hfill
  \begin{minipage}[t]{0.235\linewidth}\poptile{poporange}{Méthodes}{et exercices résolus}\end{minipage}\hfill
  \begin{minipage}[t]{0.235\linewidth}\poptile{poppink}{Exercices}{type examen + corrigés}\end{minipage}\hfill
  \begin{minipage}[t]{0.235\linewidth}\poptile{popgreen}{Fiche bilan}{l'essentiel à retenir}\end{minipage}\par}

% Sommaire
\newtcolorbox{sommaire}{enhanced,colback=white,colframe=popink,boxrule=2.2pt,arc=10pt,
  drop shadow=popink,left=18pt,right=18pt,top=13pt,bottom=13pt,before skip=6pt,after skip=20pt,
  before upper={\poppill{Sommaire}{poppurple}{white}\par\vspace{9pt}\popsemi\fontsize{10.6}{14.5}\selectfont\color{popink}}}

% Signature de fin de cours — poussée tout en bas de la dernière page
\newcommand{\findecours}{%
  \par\vfill\nopagebreak
  \begin{center}\popsquiggle{poporange}\end{center}\vspace{4pt}
  \signatureonbuch
}
\AtBeginDocument{\def\llbracket{\mathopen{[\mkern-3.5mu[}}\def\rrbracket{\mathclose{]\mkern-3.5mu]}}} % intervalles d'entiers (glyphe ⟦ absent de la police)
% Glyphes absents de Fira Math : redéfinis à partir de glyphes disponibles
\AtBeginDocument{%
  \def\setminus{\mathbin{\backslash}}\let\smallsetminus\setminus
  \def\vdots{\mathord{\vbox{\baselineskip3.2pt\lineskiplimit0pt\kern4pt\hbox{.}\hbox{.}\hbox{.}}}}%
  \def\ddots{\mathinner{\mkern1mu\raise6.5pt\hbox{.}\mkern2mu\raise3.6pt\hbox{.}\mkern2mu\raise0.7pt\hbox{.}\mkern1mu}}%
  \def\triangle{\mathord{\scalebox{0.9}{$\symup{\Delta}$}}}%
  \def\nabla{\mathord{\raisebox{\depth}{\scalebox{1}[-1]{$\symup{\Delta}$}}}}%
  \def\checkmark{\popcheck{popdark}}%
  \def\star{\mathbin{\ast}}\def\bigstar{\mathord{\ast}}%
  \def\diamond{\mathbin{\scalebox{0.6}{$\Diamond$}}}%
  \def\bigcup{\mathop{\mathchoice{\raisebox{-0.3em}{\scalebox{1.7}{$\cup$}}}{\scalebox{1.25}{$\cup$}}{\cup}{\cup}}\displaylimits}%
  \def\bigcap{\mathop{\mathchoice{\raisebox{-0.3em}{\scalebox{1.7}{$\cap$}}}{\scalebox{1.25}{$\cap$}}{\cap}{\cap}}\displaylimits}%
  \def\lesseqqgtr{\mathrel{\lessgtr}}\def\bigoplus{\mathop{\oplus}\displaylimits}%
}
\providecommand{\ssi}{si et seulement si} % abréviation fréquente
\AtBeginDocument{\providecommand{\wideparen}{\overparen}} % arc de cercle

%% FIGFIX-BEGIN v4 — correctifs globaux des figures (docs/onbuch-generation/tools/figfix)
\usepackage{adjustbox}\usepackage{etoolbox}
% toute figure TikZ/pgfplots plus large que la ligne est réduite à la largeur disponible
\BeforeBeginEnvironment{tikzpicture}{\begin{adjustbox}{max width=\linewidth}}
\AfterEndEnvironment{tikzpicture}{\end{adjustbox}}
% légendes pgfplots lisibles par-dessus les courbes
\pgfplotsset{every axis/.append style={legend style={fill=white,fill opacity=0.92,text opacity=1,draw=black!15,inner sep=3pt}}}
% étiquettes d'arêtes (syntaxe edge["..."]) avec fond blanc
\tikzset{every edge quotes/.append style={fill=white,inner sep=1.2pt}}
%% FIGFIX-END
\begin{document}

\pagegarde{La marche vers l’indépendance}{Histoire}{Tle Tech}{Histoire – classe de Terminale technique}{N06}{3 séances (fiche de progression harmonisée)}{Cette leçon analyse le processus historique qui a conduit le Cameroun de la tutelle des Nations Unies à l'indépendance et à la réunification (1945-1961), en mettant en lumière le rôle décisif des figures nationalistes et les mécanismes juridiques et politiques mis en œuvre.
\vspace{4pt}
\begin{multicols}{2}\small
\begin{itemize}
\item À la fin de cette leçon, je sais expliquer le système de tutelle de l'ONU et le distinguer de l'ancien mandat de la SDN.
\item Je sais analyser les causes profondes (internes et externes) du réveil nationaliste camerounais après 1945.
\item Je sais identifier les grandes figures du nationalisme (Um Nyobé, Félix-Roland Moumié, Ahmadou Ahidjo, John Ngu Foncha, etc.) et comparer leurs stratégies politiques.
\item Je sais ordonner chronologiquement les étapes majeures : insurrection de 1955, loi-cadre Defferre, référendum de 1961, réunification.
\item Je sais construire et commenter un croquis cartographique montrant l'évolution des frontières et le scrutin de 1961.
\item Je sais rédiger une argumentation structurée justifiant le choix de la réunification par les populations du Southern Cameroons.
\end{itemize}\end{multicols}}

\begin{sommaire}
\begin{multicols}{2}
\begin{enumerate}
\item 1. Le Cameroun sous la tutelle de l'ONU (1946-1960) : un double héritage
\item 2. L'éveil et la structuration du nationalisme camerounais (1945-1955)
\item 3. Dossier 3 : Les grandes figures du nationalisme camerounais - Portraits croisés
\item 4. De la revendication à la souveraineté : Étapes clés (1955-1961)
\item 5. Construire l'État-nation : Défis de l'unité, de la réconciliation et de la mémoire (1961-1972 / Aujourd'hui)
\item Activité d'intégration
\item Méthodes et exercices résolus
\item Exercices
\item Corrigés des exercices
\item Fiche bilan
\end{enumerate}
\end{multicols}
\end{sommaire}

\begin{objectifs}
\begin{itemize}
\item À la fin de cette leçon, je sais expliquer le système de tutelle de l'ONU et le distinguer de l'ancien mandat de la SDN.
\item Je sais analyser les causes profondes (internes et externes) du réveil nationaliste camerounais après 1945.
\item Je sais identifier les grandes figures du nationalisme (Um Nyobé, Félix-Roland Moumié, Ahmadou Ahidjo, John Ngu Foncha, etc.) et comparer leurs stratégies politiques.
\item Je sais ordonner chronologiquement les étapes majeures : insurrection de 1955, loi-cadre Defferre, référendum de 1961, réunification.
\item Je sais construire et commenter un croquis cartographique montrant l'évolution des frontières et le scrutin de 1961.
\item Je sais rédiger une argumentation structurée justifiant le choix de la réunification par les populations du Southern Cameroons.
\item J'applique la notion d'unité nationale à une situation concrète actuelle (ex: décentralisation, bilinguisme, vivre-ensemble).
\item Je sais exploiter un corpus documentaire (discours, affiches, statistiques, cartes) pour produire une synthèse historique rigoureuse.
\end{itemize}
\end{objectifs}

\begin{prerequis}
\begin{itemize}
\item La conférence de Brazzaville (1944) et la fin du code de l'indigénat (notion de 4ème / 1ère).
\item La géographie physique et humaine du Cameroun (régions, groupes ethniques, ressources) - début d'année Tle.
\item Les concepts de mandat (SDN) et de tutelle (ONU) vus en classe de 1ère (Histoire/Géographie).
\item Les grandes lignes de la Seconde Guerre mondiale et l'affaiblissement des puissances coloniales.
\item La structure de base d'une dissertation historique et la méthode d'analyse de document.
\end{itemize}
\end{prerequis}

\newpage

\coursec{Le Cameroun sous la tutelle de l'ONU (1946-1960) : un double héritage}

\textbf{Situation d'accroche.} En 1952, à Édéa, un jeune ouvrier de l'Alucam écoute, caché derrière une porte, les débats houleux de l'\cle{ALCAM} (Assemblée législative du Cameroun) sur le sort du territoire : faut-il demander l'indépendance immédiate ou une autonomie progressive sous tutelle ? Les voix s'affrontent, les positions se durcissent. Soixante-dix ans plus tard, face aux défis de l'intégration régionale et de la décentralisation, la question de la souveraineté réelle et de l'unité nationale reste au cœur de l'actualité camerounaise.

\textbf{Problématique.} Comment les pères de la nation ont-ils transformé un territoire sous tutelle, divisé entre deux puissances administratrices, en un État souverain uni ? Pour répondre à cette question, il faut d'abord comprendre ce qu'était la tutelle de l'ONU, comment elle fonctionnait, et pourquoi le Cameroun y occupait une place singulière : celle d'un territoire unique administré selon deux logiques coloniales opposées.

\courssub{Du mandat de la SDN à la tutelle de l'ONU : un changement de statut}

\textbf{Intuition.} Imagine un élève placé sous la surveillance d'un tuteur. Sous le mandat de la SDN, le tuteur rendait peu de comptes et pouvait garder l'élève presque indéfiniment. Sous la tutelle de l'ONU, le tuteur doit rendre un bulletin chaque année, recevoir des inspecteurs, et surtout préparer l'élève à devenir autonome : la tutelle a une fin prévue. C'est exactement la différence entre le mandat et la tutelle.

\textbf{Rappel chronologique.} Après la défaite allemande de 1918, le Cameroun, ancienne colonie allemande (\cle{Kamerun}, 1884-1916), est partagé entre la France et le Royaume-Uni. En 1919, la Société des Nations (SDN) place ces deux parties sous \cle{mandat} : la France administre environ les quatre cinquièmes du territoire (le Cameroun oriental, capitale Yaoundé), le Royaume-Uni administre deux étroites bandes occidentales le long de la frontière nigériane (le Cameroun britannique, capitale Buéa). Le mandat de la SDN contrôle mal les puissances mandataires : la SDN est faible et disparaît dans la Deuxième Guerre mondiale.

En 1945, l'Organisation des Nations Unies (ONU) remplace la SDN. Le \cle{13 décembre 1946}, l'Assemblée générale de l'ONU approuve les \cle{accords de tutelle} pour le Cameroun français et le Cameroun britannique. Le Cameroun devient ainsi un \emph{territoire sous tutelle}, régi par le Chapitre XII de la Charte de l'ONU.

\begin{definition}[Tutelle de l'ONU]
La \cle{tutelle} est un système de contrôle international institué par le Chapitre XII de la Charte de l'ONU (1945), succédant au régime des mandats de la SDN. Une puissance administrante y gère un territoire non autonome sous la surveillance de l'ONU, avec l'obligation explicite de conduire ce territoire vers l'\cle{autonomie} ou l'\cle{indépendance}. Contrairement à la colonisation classique, la tutelle fixe donc un objectif final : la sortie de la dépendance.
\end{definition}

\textbf{Les obligations de la puissance administrante.} L'accord de tutelle n'est pas un blanc-seing. La France et le Royaume-Uni s'engagent à :
\begin{itemize}
\item assurer le \cle{développement politique} du territoire : création d'assemblées représentatives, participation croissante des populations à la gestion de leurs affaires ;
\item promouvoir le \cle{développement économique et social} : infrastructures, santé, travail ;
\item développer l'\cle{éducation} : scolarisation, formation des cadres locaux ;
\item adresser chaque année un \cle{rapport annuel} au Conseil de tutelle de l'ONU ;
\item accepter les \cle{missions de visite} de l'ONU (par exemple les missions de 1952 et de 1958) et recevoir les \cle{pétitions} des habitants.
\end{itemize}

\begin{aretenir}[L'essentiel]
Le 13 décembre 1946, le Cameroun passe du mandat SDN à la tutelle ONU. Nouveauté décisive : la tutelle impose un contrôle international effectif (rapports annuels, missions de visite, pétitions) et fixe comme but final l'autonomie puis l'indépendance. C'est une brèche juridique que les nationalistes camerounais vont savoir exploiter.
\end{aretenir}

\courssub{Un territoire, deux puissances administratrices : la spécificité camerounaise}

\textbf{Intuition.} Deux frères grandissent dans deux familles d'accueil différentes : l'une parle français, applique des règles centralisées et identiques partout ; l'autre parle anglais, laisse les grands-parents gérer les affaires quotidiennes. Quand les deux frères se retrouvent adultes, ils partagent le même sang mais pas les mêmes habitudes. C'est le drame et la richesse du Cameroun sous tutelle.

Le Cameroun sous tutelle est un cas unique en Afrique : un même territoire, un même peuple aux racines communes, mais administré par \cle{deux puissances} aux méthodes opposées, la \cle{France} et le \cle{Royaume-Uni}. Il en résulte deux systèmes \cle{juridiques} (droit civil français / \emph{common law} britannique), deux systèmes \cle{éducatifs} (francophone / anglophone), deux systèmes \cle{politiques} (centralisation / autonomie locale) distincts. Ce double héritage pèse encore sur le Cameroun d'aujourd'hui.

\begin{popfigure}
\begin{center}
\begin{tikzpicture}[scale=0.95, every node/.style={font=\small}]
% Contour simplifié du Cameroun (polygone approximatif)
\draw[thick, popdark, fill=popcream] 
(0.5,0.5) -- (3.5,0.2) -- (5.0,1.5) -- (5.2,4.0) -- (4.5,6.0) -- (3.5,7.8) -- (2.5,9.2) -- (1.5,8.0) -- (0.8,6.0) -- (0.2,3.5) -- (0.3,1.5) -- cycle;
% Zone britannique (bande occidentale)
\draw[thick, poppurple, fill=poppurpleL] 
(0.2,3.5) -- (0.8,6.0) -- (1.5,8.0) -- (1.2,8.2) -- (0.6,6.2) -- (0.0,4.0) -- (-0.2,2.0) -- cycle;
% Ligne de partage
\draw[dashed, poppurple, thick] (0.2,3.5) -- (0.8,6.0) -- (1.5,8.0);
% Villes
\fill[popink] (2.2,2.0) circle (0.09) node[right, popdark]{\textbf{Yaoundé} (cap. française)};
\fill[popink] (3.2,1.0) circle (0.07) node[right, popdark]{Douala};
\fill[popink] (0.5,4.0) circle (0.07) node[left, poppurple]{\textbf{Buéa} (cap. brit.)};
\fill[popink] (0.8,5.5) circle (0.06) node[left, poppurple]{Bamenda};
\fill[popink] (2.0,7.0) circle (0.06) node[right, popdark]{Garoua};
\fill[popink] (2.8,5.0) circle (0.06) node[right, popdark]{Ngaoundéré};
% Rattachement au Nigéria
\draw[-{Stealth[length=3mm]}, very thick, poporange] (0.6,4.0) -- (-1.8,4.0);
\node[poporange, align=center, font=\footnotesize] at (-2.4,4.7) {Rattachement\\ administratif\\ au Nigéria\\ (Enugu puis Lagos)};
% Nord britannique
\node[poppurple, font=\footnotesize, align=center] at (1.0,7.2) {Nord\\ britannique};
\node[poppurple, font=\footnotesize, align=center] at (0.3,3.0) {Sud\\ brit.};
% Flèche Nord
\draw[-{Stealth[length=2.5mm]}, thick, popdark] (5.5,8.5) -- (5.5,9.3) node[above]{N};
% Légende zones
\node[fill=popcream, draw=popdark, font=\footnotesize, anchor=west] at (5.8,7.5) {Cameroun français (4/5 du territoire)};
\node[fill=poppurpleL, draw=poppurple, font=\footnotesize, anchor=west] at (5.8,6.8) {Cameroun britannique (bande occidentale)};
\end{tikzpicture}
\end{center}
\legende{Fig. 1 -- Croquis du Cameroun sous tutelle (1946-1960). 1 : zone française administrée depuis Yaoundé ; 2 : zone britannique (Buéa, Bamenda) rattachée administrativement au Nigéria voisin. Échelle indicative, contours simplifiés.}
\end{popfigure}

\courssub{L'administration française : assimilation, centralisation et répression}

\textbf{Intuition.} L'administration française agit comme un chef d'orchestre qui veut que tous les musiciens jouent la même partition, écrite à Paris. Elle croit à l'\cle{assimilation} : faire des administrés des Français à part entière, selon le modèle républicain unique.

L'administration française au Cameroun est \cle{assimilationniste} et \cle{centralisée} :
\begin{itemize}
\item le territoire est dirigé par un \cle{Haut-Commissaire} de la République, nommé à Paris, assisté d'une administration directe qui contrôle les circonscriptions ;
\item les chefferies traditionnelles sont marginalisées ou instrumentalisées, contrairement à la pratique britannique ;
\item une représentation locale est créée : l'\cle{ATCAM} (Assemblée représentative du Cameroun) en 1946, devenue \cle{ALCAM} (Assemblée législative du Cameroun) en 1952, puis Assemblée territoriale en 1956. Des Camerounais siègent aussi à l'Assemblée nationale française et à l'Assemblée de l'Union française ;
\item le développement économique est réel mais inégal : port de Douala, plantations, complexe Alucam à Édéa (inauguré en 1957), réseau routier, scolarisation en progression.
\end{itemize}

Mais cette administration se durcit face au nationalisme radical. L'\cle{UPC} (Union des populations du Cameroun), fondée en 1948, réclame l'indépendance immédiate et la réunification des deux Cameroun. Après les émeutes de mai 1955, le gouvernement français prend la \cle{loi du 13 juillet 1955} qui \emph{interdit l'UPC} et dissout ses organisations affiliées. Les dirigeants upécistes sont arrêtés, exilés ou passent dans la clandestinité : c'est le début d'une répression qui dégénérera en guerre civile larvée.

\begin{attention}[Erreur fréquente]
Ne confonds pas \cle{indépendance} et \cle{réunification}. L'\textbf{indépendance} du Cameroun français est proclamée le \textbf{1\textsuperscript{er} janvier 1960}. La \textbf{réunification} (rattachement du Cameroun méridional britannique au Cameroun indépendant) n'intervient que le \textbf{1\textsuperscript{er} octobre 1961}, après le référendum de l'ONU du 11 février 1961. Ce sont deux temps juridiques distincts, séparés par près de deux ans.
\end{attention}

\courssub{L'administration britannique : l'indirect rule et le sentiment d'abandon}

\begin{definition}[Indirect Rule]
L'\cle{Indirect Rule} (« administration indirecte ») est le système administratif britannique qui gouverne les populations par l'intermédiaire des \cle{chefferies traditionnelles} reconnues et encadrées (\emph{Native Authorities}). Il s'oppose à l'\cle{administration directe} française, qui gouverne par ses propres agents et uniformise les règles. L'indirect rule coûte moins cher à la métropole mais fige souvent les structures locales et néglige le développement.
\end{definition}

\textbf{Intuition.} Le Cameroun britannique est comme un enfant confié à une famille déjà nombreuse : on s'occupe d'abord des autres enfants. Le Royaume-Uni administre sa zone du Cameroun non pas comme une entité propre, mais comme une simple annexe du \cle{Nigéria}, sa grande colonie voisine.

Concrètement :
\begin{itemize}
\item le Cameroun britannique est \cle{rattaché administrativement au Nigéria} : il dépend d'abord de la province d'Enugu (région de l'Est nigérian), puis de Lagos ; il n'a ni budget propre significatif, ni administration autonome ;
\item l'administration britannique applique l'\cle{indirect rule} : les chefs traditionnels (fons, lamibe) collectent les impôts et rendent la justice sous supervision ;
\item la \cle{négligence économique} est patente : peu de routes, peu d'écoles, peu d'investissements. La Compagnie des plantations (CDC) domine l'économie sans retombées suffisantes pour la population ;
\item il en naît un profond \cle{sentiment d'abandon} : les populations se sentent « oubliées » entre Londres et Lagos. Ce sentiment nourrira plus tard le débat de 1961 : rejoindre le Nigéria indépendant ou réunifier avec le Cameroun français ?
\end{itemize}

\begin{exemplebox}[Exemple chiffré : l'écart de développement]
En 1955, le taux de scolarisation primaire atteint environ \cle{45 \%} au Cameroun français contre seulement environ \cle{15 \%} au Cameroun britannique (source : rapports annuels des puissances administrantes à l'ONU). L'écart se retrouve dans la santé (dispensaires, médecins) et les infrastructures routières. Cette inégalité de traitement entre les deux zones est une conséquence directe des deux philosophies administratives : investissement assimilateur d'un côté, gestion à moindre coût de l'autre.
\end{exemplebox}

\begin{popfigure}
\begin{center}
\begin{tikzpicture}
\begin{axis}[
ybar, bar width=14pt,
width=0.95\linewidth, height=6.5cm,
symbolic x coords={Éducation, Santé, Routes},
xtick=data,
ylabel={Dépenses publiques (indice, base 100)},
ymin=0, ymax=120,
legend style={at={(0.5,-0.2)}, anchor=north, legend columns=2, font=\small},
nodes near coords, every node near coord/.append style={font=\footnotesize},
ymajorgrids, grid style={popline!40},
]
\addplot[fill=popblue, draw=popdark] coordinates {(Éducation,100) (Santé,100) (Routes,100)};
\addplot[fill=poppurple, draw=popdark] coordinates {(Éducation,33) (Santé,28) (Routes,22)};
\legend{Cameroun français, Cameroun britannique}
\end{axis}
\end{tikzpicture}
\end{center}
\legende{Fig. 2 -- Investissement public comparé (éducation, santé, routes), moyennes 1950-1959, en indice base 100 pour le Cameroun français. Valeurs indicatives reconstituées d'après les rapports annuels à l'ONU : l'écart illustre la négligence de la zone britannique rattachée au Nigéria.}
\end{popfigure}

\courssub{Les instances de contrôle : le Conseil de tutelle, les missions, les pétitions}

La tutelle n'est pas qu'un principe : c'est une \cle{mécanique de contrôle} organisée en étages, que la figure 3 schématise.

\begin{popfigure}
\begin{center}
\begin{tikzpicture}[
node distance=1.0cm,
box/.style={draw=popdark, fill=popblueL, rounded corners, text width=4.8cm, align=center, font=\small, minimum height=1.1cm},
ctrl/.style={draw=poppurple, fill=poppurpleL, rounded corners, text width=3.6cm, align=center, font=\footnotesize, minimum height=0.9cm},
arr/.style={-{Stealth[length=2.5mm]}, thick, popdark}
]
\node[box, fill=popgoldL, align=center] (onu){\textbf{Assemblée générale de l'ONU}\\ (approuve les accords de tutelle, vote les résolutions)};
\node[box, below=of onu, align=center] (ct){\textbf{Conseil de tutelle}\\ (examine rapports annuels, reçoit pétitions, envoie des missions)};
\node[box, below=of ct, align=center] (pa){\textbf{Puissances administrantes}\\ France (Yaoundé) / Royaume-Uni (Buéa via Lagos)};
\node[box, below=of pa, align=center] (al){\textbf{Assemblées locales}\\ ATCAM (1946), ALCAM (1952), Assemblée territoriale (1956)};
\draw[arr] (onu) -- (ct);
\draw[arr] (ct) -- (pa);
\draw[arr] (pa) -- (al);
\node[ctrl, right=1.3cm of ct, align=center] (pet){\textbf{Pétitions}\\ Um Nyobé à l'ONU : 1952, 1953, 1954};
\node[ctrl, right=1.3cm of pa, align=center] (mis){\textbf{Missions de visite}\\ 1952, 1958 : enquêtes sur place};
\draw[-{Stealth[length=2mm]}, dashed, thick, poppurple] (pet) -- (ct);
\draw[-{Stealth[length=2mm]}, dashed, thick, poppurple] (mis) -- (pa);
\end{tikzpicture}
\end{center}
\legende{Fig. 3 -- Structure de la tutelle : la chaîne de contrôle descend de l'ONU vers les assemblées locales ; les pétitions des nationalistes et les missions de visite remontent l'information vers le Conseil de tutelle.}
\end{popfigure}

\textbf{Le Conseil de tutelle} examine chaque année le rapport de la puissance administrante et l'interroge publiquement. Les \cle{missions de visite} (1952, 1958) se déplacent sur le terrain, rencontrent les populations, les chefs et les partis politiques, puis publient des rapports parfois très critiques.

Surtout, la tutelle ouvre un droit inédit : le \cle{droit de pétition}. Les nationalistes camerounais l'utilisent massivement. \cle{Ruben Um Nyobé}, secrétaire général de l'UPC, se rend trois fois à New York pour plaider devant l'ONU (1952, 1953, 1954) : il y dénonce la répression française et réclame l'indépendance et la réunification. La tribune de l'ONU devient ainsi une arme politique : ce qui ne peut pas être dit à Yaoundé peut être dit à New York.

\begin{methode}[Analyser un rapport de mission de l'ONU]
Face à un extrait de rapport de mission (ex. rapport de 1958), procède en trois étapes :
\begin{enumerate}
\item \textbf{Identifier le ton} : le rapport est-il critique ou élogieux envers la puissance administrante ? Repère les adjectifs et les jugements de valeur.
\item \textbf{Extraire les chiffres clés} : taux de scolarisation, nombre de dispensaires, kilomètres de routes, évolution des assemblées. Compare-les entre les deux zones.
\item \textbf{Dégager les recommandations politiques} : le rapport propose-t-il une date d'indépendance, des élections, une amnistie ? Relie chaque recommandation à la marche vers la souveraineté.
\end{enumerate}
Conclus toujours en situant le document dans la chronologie : un rapport critique accélère la pression internationale sur la puissance administrante.
\end{methode}

\courssub{Complément Bac : la résolution 1476 (XIV) et la fixation de l'indépendance}

Sous la pression des pétitions, des rapports de mission et de l'évolution politique interne (loi-cadre de 1956, gouvernement autonome d'Ahmadou Ahidjo en 1958), la France accepte le principe de l'indépendance. Le \cle{12 décembre 1959}, l'Assemblée générale de l'ONU adopte la \cle{résolution 1476 (XIV)} : elle décide que les accords de tutelle concernant le Cameroun français prendront fin le \cle{1\textsuperscript{er} janvier 1960}, date à laquelle le territoire accède à l'indépendance sous le nom de République du Cameroun. La tutelle a ainsi tenu sa promesse juridique : conduire le territoire à la souveraineté.

\begin{exoresolu}[Exploiter un document : la résolution 1476 (XIV)]
\textbf{Énoncé.} Document : « L'Assemblée générale [...] décide que les accords de tutelle concernant le Cameroun sous administration française [...] cesseront d'être en vigueur le 1\textsuperscript{er} janvier 1960, date à laquelle le Cameroun accédera à l'indépendance. » (Résolution 1476 (XIV), 12 décembre 1959). Question : montrez que ce texte illustre le rôle moteur de l'ONU dans la marche vers l'indépendance du Cameroun.
\tcblower
\textbf{Solution détaillée.} 1) \emph{Nature du document} : il s'agit d'une résolution de l'Assemblée générale de l'ONU, c'est-à-dire une décision de l'organe supérieur de contrôle de la tutelle. 2) \emph{Contenu} : l'ONU ne se contente pas de constater, elle \emph{décide} : elle fixe elle-même la date de fin de la tutelle et la date d'indépendance (1\textsuperscript{er} janvier 1960). 3) \emph{Interprétation} : ce texte prouve que la tutelle n'était pas une colonisation déguisée mais un régime transitoire à terme fixé. L'ONU, saisie par les rapports annuels, les missions de visite et les pétitions des nationalistes (Um Nyobé dès 1952), a transformé la revendication nationaliste en échéance juridique internationale. 4) \emph{Conclusion} : la résolution 1476 (XIV) couronne la stratégie des nationalistes qui ont fait de la tribune onusienne un levier ; elle montre que l'indépendance camerounaise est à la fois le fruit des luttes internes et du mécanisme international de la tutelle.
\end{exoresolu}

\begin{savaistu}[Le savais-tu ?]
Le Cameroun est le \textbf{seul territoire d'Afrique} à avoir accédé à l'indépendance par la \cle{réunification} de deux zones coloniales différentes -- l'une française, l'autre britannique -- validée par un \cle{référendum organisé par l'ONU} (11 février 1961). Le Cameroun septentrional britannique a choisi le Nigéria, le Cameroun méridional britannique a choisi la réunification, effective le 1\textsuperscript{er} octobre 1961. Cette singularité explique le bilinguisme et le bijuridisme du Cameroun actuel.
\end{savaistu}

\begin{exercice}[À toi de jouer]
1. Définis la tutelle de l'ONU et cite deux obligations de la puissance administrante. 2. Construis un tableau comparatif de l'administration française et de l'administration britannique au Cameroun (méthode, capitale, développement, rapport aux chefferies). 3. Explique pourquoi le droit de pétition à l'ONU a constitué une arme pour les nationalistes camerounais.
\end{exercice}

\begin{corrige}[Exercice]
1. La tutelle est le système de contrôle international (Chapitre XII de la Charte de l'ONU) succédant au mandat, visant l'autonomie puis l'indépendance des territoires non autonomes. Obligations : développement politique, économique, social et éducatif ; rapport annuel à l'ONU ; acceptation des missions de visite et des pétitions. 2. Tableau attendu : France = administration directe, assimilation, centralisation, capitale Yaoundé, investissements réels mais répression de l'UPC (loi du 13 juillet 1955) ; Royaume-Uni = indirect rule via les chefferies, rattachement au Nigéria (Enugu puis Lagos), capitale Buéa, négligence économique et sentiment d'abandon. 3. Le droit de pétition permet de contourner la censure coloniale : Um Nyobé porte la revendication d'indépendance et de réunification devant l'Assemblée générale (1952, 1953, 1954), internationalise la question camerounaise et met la France sous pression de l'opinion mondiale, ce qui aboutit à la résolution 1476 (XIV) de 1959.
\end{corrige}

\begin{aretenir}[Bilan de la section]
De 1946 à 1960, le Cameroun vit sous un régime juridique original : la tutelle de l'ONU, qui oblige les puissances administrantes à préparer l'indépendance. Mais ce territoire unique reçoit un \cle{double héritage} : une zone française centralisée, assimilatrice et relativement développée, marquée par la répression de l'UPC ; une zone britannique gouvernée par indirect rule, rattachée au Nigéria et négligée. Le contrôle de l'ONU (rapports, missions, pétitions) offre aux nationalistes une tribune décisive et conduit à la résolution 1476 (XIV) fixant l'indépendance au 1\textsuperscript{er} janvier 1960. Restera alors à réunifier les deux héritages : ce sera l'enjeu de 1961.
\end{aretenir}

\coursec{L'éveil et la structuration du nationalisme camerounais (1945-1955)}

\textbf{Problématique.} Comment, au sortir de la Deuxième Guerre mondiale, les Camerounais passent-ils de la revendication syndicale et culturelle à la revendication politique nationale d'indépendance et de réunification ?

\courssub{Le contexte international et local : le souffle de la liberté}

La fin de la Deuxième Guerre mondiale (1945) bouleverse les consciences. Les \cle{tirailleurs camerounais} rentrent du front avec l'expérience du combat pour la liberté et la découverte d'autres sociétés. La Charte de l'ONU (1945) proclame le \cle{droit des peuples à disposer d'eux-mêmes}. La Conférence de Brazzaville (1944) a refusé l'indépendance mais promis des réformes (citoyenneté, fin du code de l'indigénat). Au Cameroun, la déception face à la lenteur des réformes et la persistance des discriminations (travail forcé, code du travail inégalitaire) radicalise les esprits.

\courssub{Les premières organisations : syndicalisme et culturel (1945-1948)}

Avant les partis politiques, ce sont les syndicats et associations culturelles qui structurent la contestation.
\begin{itemize}
\item \cle{USEC} (Union des Syndicats Confédérés du Cameroun, 1945) : affiliée à la CGT française, elle mène des grèves pour les salaires et les conditions de travail (grève des dockers de Douala 1945, grève générale 1947). Elle forme une génération de militants (Um Nyobé, Moumié).
\item \cle{JEUCAFRA} (Jeunesse Camerounaise Française, 1946) : association culturelle et sportive qui devient un vivier politique. Elle réclame l'abolition du code de l'indigénat et l'accès à la citoyenneté.
\item Côté britannique : \cle{KNC} (Kamerun National Congress, 1945) et \cle{KPP} (Kamerun People's Party, 1948) militent pour la réunification et l'autonomie vis-à-vis du Nigéria.
\end{itemize}

\begin{aretenir}[Continuité]
Le nationalisme camerounais naît d'abord \emph{par le bas} : luttes sociales (syndicats), revendications culturelles (JEUCAFRA), avant de se structurer en partis politiques revendiquant la souveraineté nationale.
\end{aretenir}

\courssub{La naissance de l'UPC (1948) : le parti de l'indépendance immédiate}

Le \cle{10 avril 1948}, à la salle des fêtes de Douala (Bassa), est fondée l'\cle{Union des Populations du Cameroun (UPC)}. Son mot d'ordre : \textbf{« Indépendance immédiate et réunification »}.

\begin{definition}[L'UPC : caractéristiques]
\begin{itemize}
\item \textbf{Parti de masse} : recrutement transethnique, implanté dans les villes (Douala, Yaoundé, Édéa, Nkongsamba) et les campagnes (Sanaga-Maritime, Bamiléké).
\item \textbf{Idéologie} : nationalisme radical, anti-impérialisme, panafricanisme, socialisme non-aligné.
\item \textbf{Direction collégiale} : Secrétaire général \cle{Ruben Um Nyobé} (charisme, rigueur intellectuelle), Président \cle{Félix-Roland Moumié} (médecin, diplomatie internationale), Vice-président \cle{Ernest Ouandié} (organisation, maquis).
\item \textbf{Méthodes} : action légale (pétitions ONU, élections), mobilisation populaire (meetings, grèves), presse (\emph{La Voix du Cameroun}, \emph{Étoile}).
\end{itemize}
\end{definition}

L'UPC remporte les premières élections à suffrage universel (municipales 1951, législatives ALCAM 1952) dans les grandes villes. Um Nyobé dépose des pétitions à l'ONU (1952, 1953, 1954) dénonçant la répression et réclamant un calendrier d'indépendance.

\courssub{La riposte coloniale et l'entrée dans la clandestinité (1955-1956)}

Face à la montée de l'UPC, l'administration française durcit le ton.
\begin{itemize}
\item \textbf{Mai 1955} : Émeutes à Douala, Yaoundé, Édéa, Nkongsamba. Répression sanglante (dizaines de morts officiels, centaines selon l'UPC).
\item \textbf{13 juillet 1955} : Loi interdisant l'UPC et ses organisations satellites (JEUCAFRA, USEC). Dirigeants arrêtés (Um Nyobé assigné à résidence puis clandestin), exilés (Moumié au Ghana, puis Guinée, Égypte) ou passés en maquis (Ouandié, Woungly-Massaga).
\item \textbf{1956} : Promulgation de la \cle{Loi-cadre Defferre} (23 juin) qui accorde l'autonomie interne. L'UPC, interdite, boycotte les élections de l'Assemblée territoriale (décembre 1956) remportées par les modérés (André-Marie Mbida).
\end{itemize}

\begin{attention}[Piège historiographique]
Ne pas réduire le nationalisme à l'UPC. Des figures \cle{modérées} (Mbida, Ahidjo, Foncha) participent aux institutions de la tutelle pour arracher l'indépendance « par étapes ». Le nationalisme camerounais est \textbf{pluriel} : radical (UPC) vs modéré (partis dits « légalistes »), mais tous partagent l'objectif de la réunification.
\end{attention}

\begin{popfigure}
\begin{center}
\begin{tikzpicture}[
timeline/.style={thick, popdark},
event/.style={circle, fill=popblue, draw=popdark, minimum size=3pt, inner sep=0},
eventlabel/.style={font=\footnotesize, align=center, text width=3.5cm},
datelabel/.style={font=\footnotesize\bfseries, popdark, anchor=east, xshift=-0.3cm}
]
% Ligne de temps
\draw[timeline] (0,0) -- (16,0);
% Années
\foreach \x/\year in {0/1945, 2/1947, 4/1948, 6/1950, 8/1952, 10/1954, 12/1955, 14/1956, 16/1957} {
    \draw (\x,0.1) -- (\x,-0.1) node[below, datelabel] {\year};
}
% Événements
\draw (0,0) node[event] {} node[eventlabel, above=0.3cm, align=center]{Fin 2nde GM\\ Retour tirailleurs};
\draw (2,0) node[event] {} node[eventlabel, above=0.3cm, align=center]{Grèves USEC\\ (Douala)};
\draw (4,0) node[event] {} node[eventlabel, above=0.3cm, fill=popgoldL, rounded corners, align=center]{\textbf{Création UPC}\\(10 avril 1948)};
\draw (6,0) node[event] {} node[eventlabel, below=0.3cm, align=center]{Pétitions Um Nyobé\\ à l'ONU (52, 53, 54)};
\draw (8,0) node[event] {} node[eventlabel, above=0.3cm, align=center]{Élections ALCAM\\ Victoire UPC villes};
\draw (12,0) node[event] {} node[eventlabel, below=0.3cm, fill=poppinkL, rounded corners, align=center]{\textbf{Émeutes mai 1955}\\\textbf{Interdiction UPC (juil)}};
\draw (14,0) node[event] {} node[eventlabel, above=0.3cm, align=center]{Loi-cadre Defferre\\ Autonomie interne};
\draw (16,0) node[event] {} node[eventlabel, below=0.3cm, align=center]{Gouvernement Mbida\\ (1er africain)};
\end{tikzpicture}
\end{center}
\legende{Fig. 4 -- Frise chronologique : l'éveil du nationalisme camerounais (1945-1957). En jaune : création UPC ; en rose : basculement dans la répression/clandestinité.}
\end{popfigure}

\begin{exercice}[À toi de jouer]
1. Cite trois facteurs internationaux ayant favorisé l'éveil nationaliste après 1945. 2. Présente l'UPC (date, mot d'ordre, trois dirigeants). 3. Explique pourquoi l'année 1955 constitue une rupture majeure.
\end{exercice}

\begin{corrige}[Exercice]
1. Charte ONU (droit des peuples), Guerre froide (bloc afro-asiatique Bandung 1955), décolonisation asiatique (Inde 1947). 2. 10 avril 1948 ; « Indépendance immédiate et réunification » ; Um Nyobé (SG), Moumié (Président), Ouandié (VP). 3. Émeutes de mai -> Répression -> Loi du 13 juillet interdisant l'UPC -> Clandestinité des dirigeants -> Début de la guerre de libération (maquis) -> Fin de la légalité pour le nationalisme radical.
\end{corrige}

\coursec{Dossier 3 : Les grandes figures du nationalisme camerounais - Portraits croisés}

\textbf{Consigne.} Ce dossier est au programme du Probatoire/Baccalauréat. Tu dois connaître le parcours, le rôle et le sort de chaque figure. Utilise le tableau comparatif pour réviser.

\courssub{Les figures de l'UPC : le sacrifice radical}

\begin{center}
\begin{tabularx}{\linewidth}{|l|X|X|X|}\hline
\rowcolor{popblueL}
\textbf{Figure} & \textbf{Parcours \& Rôle} & \textbf{Action clé} & \textbf{Sort} \\ \hline
\textbf{Ruben Um Nyobé} (1913-1958) & Instituteur, syndicaliste (USEC). Secrétaire général UPC (1948-1958). « Mpodol » (porte-parole). & Stratégie \emph{légale} : 3 pétitions à l'ONU (52, 53, 54). Refus de la violence aveugle. Recherche dialogue avec France. & Tué par l'armée française le \textbf{13 sept. 1958} (maquis Sanaga-Maritime). Corps exposé à Boumnyébel. \\ \hline
\textbf{Félix-Roland Moumié} (1925-1960) & Médecin. Président UPC (1952). Diplomate de l'ombre. Exil : Ghana, Guinée, Égypte, Maroc. & \emph{Diplomatie internationale} : Rallye pays non-alignés, Chine, URSS. Obtient reconnaissance UPC comme mouvement de libération. & Assassiné par empoisonnement (thallium) à \textbf{Genève (Suisse)} le \textbf{3 nov. 1960} (main rouge/SDECE). \\ \hline
\textbf{Ernest Ouandié} (1924-1971) & Instituteur. Vice-président UPC. Chef militaire. Organisateur maquis (Bamiléké, Sanaga). & \emph{Lutte armée} : Structure maquis, code de conduite. Continue lutte après 1960 contre « néocolonialisme ». & Arrêté \textbf{août 1970}. Jugé, condamné à mort. Exécuté publiquement \textbf{15 janv. 1971} (Bafoussam). \\ \hline
\end{tabularx}
\end{center}

\courssub{Les figures de la voie légale/institutionnelle : les bâtisseurs de l'État}

\begin{center}
\begin{tabularx}{\linewidth}{|l|X|X|X|}\hline
\rowcolor{popgreenL}
\textbf{Figure} & \textbf{Parcours \& Rôle} & \textbf{Action clé} & \textbf{Sort / Héritage} \\ \hline
\textbf{Ahmadou Ahidjo} (1924-1989) & Fils de chef, fonctionnaire. Député ALCAM. Vice-PM (1957), PM (1958). Président (1960-1982). & \cle{« Indépendance dans l'interdépendance »}. Négocie fin tutelle (Rés. 1476). Gère réunification 1961. Instauré parti unique (UNC 1966). & Démission 1982. Exil France. Réhabilité 1991. « Père de la Nation » officiel. \\ \hline
\textbf{André-Marie Mbida} (1917-1980) & Premier africain diplômé droit (Paris). Député français. 1er Chef de gouvernement (1957-58). & Refus ingérence France. Défense chefferies. Chute : motion censure Assemblée (soutien France/Ahidjo). & Exil (Conakry, Paris). Retour 1990. Figure démocratie chrétienne. \\ \hline
\textbf{John Ngu Foncha} (1916-1999) & Instituteur (Sud britannique). Fondateur KNDP (1955). PM Cameroun méridional (1959). & Artisan \cle{réunification}. Campagne « Votez Vert » (référendum 1961). VP République Fédérale (1961-70). & Démission 1970 (opposition unitarisme). Retraite. Symbole anglophone. \\ \hline
\end{tabularx}
\end{center}

\courssub{Synthèse comparative : deux stratégies, une même fin ?}

\begin{popfigure}
\begin{center}
\begin{tikzpicture}[
mindmap, grow cyclic, every node/.style=concept, concept color=popblue!60,
level 1/.append style={level distance=4.5cm, sibling angle=120},
level 2/.append style={level distance=3.2cm, sibling angle=60},
]
\node[concept color=popdark] {Nationalisme camerounais}
child [concept color=popblue] {node {Voie Radicale (UPC)}
  child {node {Um Nyobé : Légalité ONU}}
  child {node {Moumié : Diplomatie}}
  child {node {Ouandié : Maquis}}}
child [concept color=popgreen] {node {Voie Institutionnelle}
  child {node {Ahidjo : Négociation Etapiste}}
  child {node {Mbida : Légalisme Parlementaire}}
  child {node {Foncha : Fédéralisme Réunification}}}
child [concept color=poporange] {node {Convergences}
  child {node {Réunification}}
  child {node {Indépendance}}
  child {node {Fin travail forcé}}};
\end{tikzpicture}
\end{center}
\legende{Fig. 5 -- Carte mentale : complémentarité et tensions entre les figures du nationalisme.}
\end{popfigure}

\begin{aretenir}[Bilan Dossier 3]
L'UPC (Um Nyobé, Moumié, Ouandié) incarne la \cle{radicalité} : indépendance immédiate, réunification non négociable, recours à l'ONU puis à la lutte armée. Sacrifice ultime (mort violente). La voie institutionnelle (Ahidjo, Mbida, Foncha) incarne le \cle{réalisme politique} : autonomie progressive, négociation avec la France, construction de l'État (fédéral puis unitaire). Les deux voies ont été \textbf{nécessaires} : la radicalité a imposé l'échéance (ONU, opinion mondiale), le réalisme a géré la transition et la souveraineté.
\end{aretenir}

\begin{exoresolu}[Sujet type Bac : « Les figures du nationalisme camerounais illustrent la diversité des chemins vers l'indépendance. » Discutez.]
\textbf{Énoncé.} À partir de vos connaissances et du dossier, montrez que l'indépendance du Cameroun résulte de la convergence de stratégies distinctes.
\tcblower
\textbf{Plan détaillé.}
\textbf{Intro.} Accroche (double héritage, double stratégie). Définition nationalisme. Problématique. Annonce plan.
\textbf{I. La voie radicale (UPC) : la pression par le bas et l'international.}
A. Um Nyobé : légalité onusienne (pétitions 52-54), internationalisation conflit.
B. Moumié : diplomatie des non-alignés, armement maquis.
C. Ouandié : maquis intérieur, coût humain élevé. \emph{Apport : calendrier imposé à la France.}
\textbf{II. La voie institutionnelle : la gestion de la transition.}
A. Ahidjo : loi-cadre 1956, autonomie 1957, négociation résolution 1476 (1959).
B. Mbida : 1er gouvernement africain, tentative d'émancipation réelle.
C. Foncha : KNDP, campagne réunification, victoire référendum 1961. \emph{Apport : structures étatiques, légitimité internationale.}
\textbf{III. Convergence et tensions (1958-1961).}
A. Convergence : Réunification (objectif commun), Indépendance (1er janv 1960).
B. Tensions : Répression UPC par Ahidjo (accords défense 1960), élimination rivaux (Ouandié 1971).
\textbf{Conclusion.} Indépendance = fruit de la dialectique pression radicale / gestion institutionnelle. Héritage complexe : mémoire blessée (UPC) vs État construit (Ahidjo).
\end{exoresolu}

\coursec{De la revendication à la souveraineté : Étapes clés (1955-1961)}

\courssub{1955-1958 : La guerre cachée et l'autonomie interne}

\begin{itemize}
\item \textbf{1955-1957} : Maquis UPC (Sanaga-Maritime, Bamiléké). Répression française (corps expéditionnaire, regroupements de villages « villages de regroupement »). Parallèlement, vie parlementaire : ALCAM (majorité modérée), gouvernement Mbida (mai 1957-fév 1958).
\item \textbf{Février 1958} : Chute de Mbida (motion de censure). \cle{Ahmadou Ahidjo} devient Vice-Président du Conseil (Chef de gouvernement). Il négocie avec Paris : \cle{« Indépendance dans l'interdépendance »} (coopération militaire, monétaire, technique).
\item \textbf{1958} : Mission de visite ONU (rapport favorable à l'indépendance). Ahidjo demande la fixation de la date.
\end{itemize}

\courssub{1959-1960 : L'indépendance du Cameroun français}

\begin{center}
\begin{tabularx}{\linewidth}{|c|X|X|}\hline
\rowcolor{popgoldL}
\textbf{Date} & \textbf{Événement} & \textbf{Signification} \\ \hline
\textbf{13 mars 1959} & Accords de coopération Franco-Camerounais (Yaoundé) & Cadre post-indépendance : défense, monnaie (FCFA), aide technique. \\ \hline
\textbf{12 déc. 1959} & \textbf{Résolution 1476 (XIV) ONU} & Fin tutelle fixée au 1er janv 1960. Victoire diplomatique. \\ \hline
\textbf{1er janv 1960} & \textbf{Proclamation Indépendance} (Yaoundé) & Naissance \cle{République du Cameroun}. Ahidjo Président. Adhésion ONU (20 sept 1960). \\ \hline
\end{tabularx}
\end{center}

\courssub{Le référendum et la réunification}

Le sort du Cameroun britannique (toujours sous tutelle) doit être réglé. L'ONU organise un \cle{référendum} le \cle{11 février 1961} : deux questions, deux réponses.

\begin{popfigure}
\begin{center}
\begin{tikzpicture}[scale=0.9]
% Carte simplifiée zone britannique
\draw[thick, poppurple, fill=poppurpleL] (0,0) -- (4,0) -- (4,5) -- (3.5,4) -- (1,4.5) -- (0,3) -- cycle;
\draw[dashed, thick, popdark] (1.5,0.5) -- (1.5,3.5); % Ligne de séparation Nord/Sud
% Nord
\node[align=center, font=\small] at (0.7,1.5) {\textbf{Nord britannique}\\\textbf{Vote :} Rejoindre Nigéria\\\textbf{Résultat :} 60\% vs 40\%};
% Sud
\node[align=center, font=\small] at (3.0,2.5) {\textbf{Sud britannique}\\\textbf{Vote :} Réunification\\\textbf{Résultat :} 70\% vs 30\%};
% Flèches résultats
\draw[-{Stealth[length=3mm]}, thick, poporange] (0.7,0.5) -- (-1.0,0.5) node[left, font=\footnotesize, align=center] {Intégration\\Nigéria (mai 1961)};
\draw[-{Stealth[length=3mm]}, thick, popgreen] (3.5,3.5) -- (5.5,3.5) node[right, font=\footnotesize, align=center] {Réunification\\Rép. du Cameroun\\ (1er oct 1961)};
% Légende
\node[fill=popcream, draw=popdark, font=\footnotesize, anchor=south west] at (0,4.5) {Référendum ONU 11 fév 1961 : Choix par plébiscite};
\end{tikzpicture}
\end{center}
\legende{Fig. 6 -- Résultats du référendum de février 1961 au Cameroun britannique. Le Nord choisit le Nigéria, le Sud choisit la réunification.}
\end{popfigure}

\begin{itemize}
\item \textbf{Nord britannique} (majorité musulmane, liens économiques avec Nord Nigéria) : vote pour l'intégration au \cle{Nigéria} (indépendant depuis oct. 1960). Effectif : 1er juin 1961.
\item \textbf{Sud britannique} (majorité chrétienne, élites formées au contact UPC/Français) : vote pour la \cle{réunification} avec la République du Cameroun. Effectif : \cle{1er octobre 1961}.
\end{itemize}

\courssub{La Fédération (1961-1972) : un compromis constitutionnel}

La \cle{Conférence de Foumban} (juillet 1961) rédige la Constitution fédérale.
\begin{itemize}
\item \cle{République Fédérale du Cameroun} (1er oct 1961 - 20 mai 1972).
\item Deux États fédérés : \cle{État de l'Est} (ex-français, cap. Yaoundé) et \cle{État de l'Ouest} (ex-britannique, cap. Buéa).
\item Président fédéral : Ahidjo. VP : Foncha.
\item \cle{Bilinguisme} officiel (français/anglais), \cle{Bijuridisme} (droit civil / common law).
\item Assemblée fédérale + Assemblées fédérées.
\end{itemize}

\begin{attention}[Point de vigilance Bac]
La réunification n'est pas une « fusion » mais une \textbf{fédération} (1961). L'unité unitaire ne viendra qu'en 1972 (Référendum Ahidjo). Ne pas confondre les dates : 1960 (Indépendance Est), 1961 (Réunification/Fédération), 1972 (Unité/État unitaire).
\end{attention}

\begin{exercice}[À toi de jouer]
1. Ordre chronologique : Loi-cadre Defferre, Indépendance 1er janv 1960, Référendum 11 fév 1961, Réunification 1er oct 1961, Conférence de Foumban. 2. Pourquoi le Nord britannique a-t-il choisi le Nigéria ? 3. Cite deux caractéristiques de l'État fédéral (1961-1972).
\end{exercice}

\begin{corrige}[Exercice]
1. Loi-cadre (juin 1956) -> Indépendance (1 janv 1960) -> Référendum (11 fév 1961) -> Conférence Foumban (juil 1961) -> Réunification (1er oct 1961). 2. Affinités culturelles/religieuses (islam), liens économiques (commerce bétail/arachide vers Nord Nigéria), peur de la domination francophone/du Sud. 3. Deux États fédérés (Est/Ouest) avec gouvernements/assemblées propres ; Président fédéral + VP ; Bilinguisme officiel ; Bijuridisme.
\end{corrige}

\coursec{Construire l'État-nation : Défis de l'unité, de la réconciliation et de la mémoire (1961-Aujourd'hui)}

\courssub{De la Fédération à l'État unitaire (1961-1972) : l'unification par le haut}

La fédération est perçue par Ahidjo comme une transition. Les tensions apparaissent vite :
\begin{itemize}
\item \textbf{1962} : Référendum constitutionnel (renforcement pouvoir fédéral).
\item \textbf{1966} : Fusion partis politiques -> \cle{UNC} (Union Nationale Camerounaise), parti unique.
\item \textbf{1968-1970} : Démission de Foncha (VP) et Solomon Tandeng Muna (PM Ouest) : protestation contre la centralisation excessive, non-respect de l'esprit fédéral (bijuridisme, bilinguisme).
\item \textbf{20 mai 1972} : \cle{Référendum} sur l'État unitaire. « Oui » massif (officiel > 99\%). Naissance de la \cle{République Unie du Cameroun}. Disparition des États fédérés, du poste de VP, de l'Assemblée fédérale.
\item \textbf{1984} : \cle{Loi n°84-001} : Changement de nom -> \cle{République du Cameroun} (suppression « Uni »). Symbole : retour à l'appellation 1960, gommage de l'héritage fédéral.
\end{itemize}

\begin{savaistu}[Le savais-tu ?]
Le \textbf{20 mai} est la \cle{Fête nationale} du Cameroun. Elle commémore le référendum de 1972 (instauration de l'État unitaire), et non l'indépendance (1er janv 1960) ni la réunification (1er oct 1961). Ce choix officiel souligne la primauté accordée à l'unité unitaire sur la diversité fédérale dans le récit national dominant.
\end{savaistu}

\courssub{La question de la réconciliation : l'héritage de la guerre (UPC/Maquis)}

L'indépendance n'a pas apaisé le conflit UPC / État.
\begin{itemize}
\item \textbf{1960-1971} : Poursuite du maquis (Ouandié, Woungly). Répression féroce (armée camerounaise + conseillers français). \cle{État d'urgence} dans l'Ouest (Bamiléké).
\item \textbf{1971} : Exécution publique d'Ernest Ouandié (Bafoussam). Fin militaire de la rébellion.
\item \textbf{Années 90} : Retour au multipartisme (1990). \cle{Loi n°91-022} (1991) : Réhabilitation officielle de \cle{Ruben Um Nyobé} et des « martyrs de la libération ». Cérémonies officielles à Éséka (tombe Um Nyobé).
\item \textbf{2010} : Loi portant réhabilitation des « nationalistes » (Moumié, Ouandié, etc.). Geste mémoriel fort, mais contentieux sur les réparations et la vérité historique (archives, disparus).
\end{itemize}

\courssub{Le défi anglophone et la décentralisation (2016-Aujourd'hui) : l'héritage de 1961 en question}

Le sentiment d'abandon du Sud britannique (ex-État de l'Ouest) resurgit violemment en 2016.
\begin{itemize}
\item \textbf{Crise anglophone} : Grèves avocats/enseignants (common law, éducation) -> Répression -> Radicalisation (sécessionnistes « Ambazonie ») -> Conflit armé (Nord-Ouest, Sud-Ouest). Crise humanitaire (déplacés, écoles fermées).
\item \textbf{Réponses de l'État} : Commission nationale bilinguisme/multiculturalisme (2017). \cle{Grand Dialogue National} (oct. 2019). \cle{Statut spécial} pour les régions NW/SW (Loi déc. 2019). \cle{Décentralisation} : Élections régionales (déc. 2020), transfert compétences/ressources.
\item \textbf{Enjeu central} : Comment concilier \cle{unité nationale} (héritage Ahidjo/Biya) et \cle{diversité} (héritage fédéral/foncha) ? La gestion du \cle{bijuridisme} et du \cle{bilinguisme} réel (pas seulement symbolique) est le test de la maturité démocratique.
\end{itemize}

\begin{popfigure}
\begin{center}
\begin{tikzpicture}[
node distance=0.8cm,
state/.style={draw=popdark, fill=popblueL, rounded corners, text width=3.2cm, align=center, font=\small, minimum height=1.2cm},
trans/.style={draw=poporange, fill=poporangeL, rounded corners, text width=3.2cm, align=center, font=\small, minimum height=1cm},
arr/.style={-{Stealth[length=2.5mm]}, thick, popdark}
]
\node[state, align=center] (fed){\textbf{Rép. Fédérale (1961-72)}\\ 2 États (Est/Ouest)\\ VP Foncha\\ Bijuridisme effectif};
\node[trans, right=of fed, align=center] (ref72){\textbf{Référendum 20 mai 1972}\\ « Oui » à l'unitaire\\ Fin fédéralisme};
\node[state, right=of ref72, align=center] (uni){\textbf{Rép. Unie (1972-84)}\\ 1 État unitaire\\ Centralisation\\ UNC parti unique};
\node[state, right=of uni, align=center] (rep){\textbf{Rép. du Cameroun (1984-)}\\ Loi 84-001\\ Multipartisme (1990)\\ Crise anglophone (2016-)};
\node[trans, below=of uni, align=center] (dec){\textbf{Décentralisation (2019-20)}\\ Statut spécial NW/SW\\ Régions, Conseillers};
\draw[arr] (fed) -- (ref72);
\draw[arr] (ref72) -- (uni);
\draw[arr] (uni) -- (rep);
\draw[arr, dashed] (uni) -- (dec);
\draw[arr, dashed] (rep) -- (dec);
\end{tikzpicture}
\end{center}
\legende{Fig. 7 -- Évolution constitutionnelle : du fédéralisme (1961) à l'unitarisme (1972/84) et aux tentatives de redécentralisation (2019). Flèches pleines : ruptures constitutionnelles ; flèches pointillées : réformes récentes.}
\end{popfigure}

\begin{methode}[Répondre à une question de synthèse sur la construction nationale]
Sujet type : « Le Cameroun, une nation inachevée ? » ou « Unité et diversité au Cameroun depuis 1961 ».
\begin{enumerate}
\item \textbf{Analyser les termes} : « Inachevée » = processus continu, tensions persistantes. « Unité/Diversité » = couple conceptuel clé (fédéralisme vs unitarisme, bijuridisme).
\item \textbf{Structurer en 3 temps} : 1. Construction étatique (1960-1972 : fédéralisme -> unitarisme). 2. Gestion de la mémoire conflictuelle (UPC, maquis, réhabilitation tardive). 3. Gestion de la diversité linguistique/juridique (crise anglophone, décentralisation).
\item \textbf{Nuancer} : Éviter le manichéisme. L'unitarisme a permis l'intégration administrative ; le fédéralisme a garanti l'entrée du Sud. La réhabilitation (1991, 2010) montre une évolution. La décentralisation (2020) est une réponse institutionnelle à la crise.
\item \textbf{Conclure} : La nation camerounaise se construit \emph{par} la gestion de ses contradictions (double héritage, mémoire blessée, diversité). Le Probatoire attend une analyse équilibrée, factuelle, datée.
\end{enumerate}
\end{methode}

\begin{exercice}[Entraînement Bac - Sujet de synthèse]
\textbf{Sujet.} « De la tutelle à la décentralisation (1946-2020), le Cameroun a-t-il réussi son unité nationale ? » Vous organiserez votre réponse en trois parties : 1. La construction de la souveraineté (1946-1961). 2. Le choix de l'unitarisme et ses limites (1961-1990). 3. Les défis contemporains de l'unité (1990-2020).
\end{exercice}

\begin{corrige}[Sujet de synthèse - Éléments de corrigé]
\textbf{Intro.} Accroche : double héritage colonial. Problématique : tension entre unité juridique et unité réelle. Plan annoncé.
\textbf{I. 1946-1961 : La souveraineté par la réunification.}
A. Tutelle ONU : cadre juridique unique pour deux territoires. B. Nationalisme pluriel (UPC radical / Modérés légalistes) -> Pression internationale + Négociation. C. 1960 Indépendance Est + 1961 Référendum/Réunification = Souveraineté complète sur le triangle national.
\textbf{II. 1961-1990 : L'unitarisme comme ciment forcé.}
A. Fédération 1961 : compromis fragile (Foumban). B. Glissement unitaire : Parti unique 1966, Référendum 1972, Loi 1984. C. Limites : Élimination opposition (UPC/Ouandié), Frustration anglophone (démission Foncha 1970), Gommage mémoire (Um Nyobé tabou).
\textbf{III. 1990-2020 : L'unité mise à l'épreuve de la démocratie et de la diversité.}
A. Multipartisme 1990 : Libertés mais fragmentation ethnorégionale. B. Réhabilitation mémoire (Lois 1991, 2010) : Apaisement symbolique partiel. C. Crise anglophone 2016 : Échec de l'assimilation unitaire, revendication fédéralisme/identité. D. Décentralisation 2019-2020 : Réponse institutionnelle (statut spécial, régions), mais application inachevée (insécurité, ressources).
\textbf{Conclusion.} Unité \textbf{juridique} acquise (frontières, État, ONU). Unité \textbf{politique/sociale} inachevée : défi de la réconciliation mémorielle (UPC) et de la reconnaissance effective de la diversité (anglophone). La décentralisation est le chantier actuel de cette « nation inachevée ».
\end{corrige}

\begin{aretenir}[Bilan général de la leçon]
\begin{itemize}
\item \textbf{1946-1960} : Tutelle ONU = cadre légal + double héritage (France centralisatrice / UK indirect rule/négligence). Tribunes ONU (pétitions Um Nyobé) = levier independence.
\item \textbf{1948-1955} : Structuration nationalisme : UPC (radical, réunification) vs Modérés (institutionnel). 1955 = basculement répression/clandestinité.
\item \textbf{Figures} : Um Nyobé/Moumié/Ouandié (sacrifice, radicalité) ; Ahidjo/Mbida/Foncha (négociation, construction État). Complémentarité historique.
\item \textbf{1960-1961} : Indépendance (1er janv 60) + Réunification (1er oct 61) via référendum ONU. Fédération (Foumban).
\item \textbf{1972-1984} : Unitarisme (Référendum 72, Loi 84). Centralisation, parti unique, gommage fédéralisme.
\item \textbf{1990-Aujourd'hui} : Multipartisme, Réhabilitation mémoire (Um Nyobé), Crise anglophone (héritage 1961 non résolu), Décentralisation (chantier ouvert). La nation se construit encore.
\end{itemize}
\end{aretenir}

\coursec{L'éveil et la structuration du nationalisme camerounais (1945-1955)}

Nous avons vu dans la section précédente que le Cameroun, placé sous \cle{tutelle} de l'ONU en 1946, n'était plus une simple colonie : la France et la Grande-Bretagne devaient désormais rendre compte de leur administration devant l'Assemblée générale des Nations unies et préparer le territoire à l'autonomie. Ce statut nouveau ouvre une brèche juridique et politique. C'est précisément dans cette brèche que va s'engouffrer, entre 1945 et 1955, un mouvement inédit : le \cle{nationalisme camerounais}. Comment un peuple dispersé, divisé entre deux administrations coloniales, privé de droits politiques, parvient-il à organiser en quelques années des partis de masse capables de porter sa voix jusqu'à la tribune de l'ONU ? C'est toute la question de cette section.

\begin{definition}[Nationalisme camerounais]
Le \cle{nationalisme camerounais} est le mouvement politique qui, après la Deuxième Guerre mondiale, revendique la \textbf{souveraineté nationale}, l'\textbf{unité du territoire} (réunification des deux Cameroun) et la \textbf{fin de la domination coloniale}. Il n'est pas un bloc monolithique : il est traversé par des courants idéologiques divers — \textbf{marxiste} (indépendance immédiate, rupture avec la France), \textbf{libéral et modéré} (autonomie progressive dans le cadre de l'Union française) et \textbf{fédéraliste ou intégrationniste} (au Cameroun britannique, choix entre réunification et rattachement au Nigéria).
\end{definition}

\courssub{Les facteurs externes : un monde qui change}

Le nationalisme camerounais ne naît pas dans le vide. Il s'enracine dans un contexte international bouleversé par la guerre de 1939-1945.

\textbf{La Charte de l'Atlantique (août 1941).} Signée par Roosevelt et Churchill, elle proclame le « droit des peuples à disposer d'eux-mêmes ». Même si Churchill prétend qu'elle ne s'applique pas aux colonies, les élites africaines instruites s'en emparent immédiatement : si les Européens se battent contre le nazisme au nom de la liberté, pourquoi refuser cette liberté aux Africains ?

\textbf{La Conférence de Brazzaville (janvier-février 1944).} Réunie par le général de Gaulle, elle reconnaît que la France doit des réformes à l'Afrique qui l'a aidée à se libérer : suppression du travail forcé (1946), abolition de l'\cle{indigénat} (code répressif réservé aux Africains), représentation des colonies au Parlement français. Mais elle exclut toute indépendance : « les fins de l'œuvre de civilisation [...] écartent toute idée d'autonomie ». C'est une porte entrouverte, pas ouverte.

\textbf{La Guerre froide (1947-1991).} L'affrontement entre le bloc américain et le bloc soviétique transforme chaque colonie en enjeu stratégique. Les mouvements indépendantistes peuvent chercher des appuis à Moscou ou jouer de la peur du communisme pour faire pression sur les Occidentaux. L'UPC sera d'ailleurs accusée — à tort ou à raison — de « communisme » par l'administration française, ce qui justifiera sa répression.

\textbf{La décolonisation asiatique.} L'indépendance de l'\textbf{Inde} (1947), de l'Indonésie (1949), la victoire de \textbf{Diên Biên Phu} au Vietnam (1954) prouvent qu'un empire colonial peut être vaincu. Enfin, la \textbf{Conférence de Bandung} (avril 1955) réunit 29 pays d'Asie et d'Afrique et affirme le droit des peuples colonisés à l'indépendance. Le Cameroun n'y est pas représenté, mais l'événement a un immense retentissement chez les militants.

\courssub{Les facteurs internes : le terreau des mécontentements}

L'intuition est simple : on ne fait pas un mouvement de masse avec des idées seules. Il faut des souffrances concrètes, quotidiennes, que le discours nationaliste va transformer en revendication politique.

\begin{itemize}
\item \textbf{La frustration des anciens combattants.} Des dizaines de milliers de \cle{tirailleurs} camerounais ont combattu pour la France (1940-1945). De retour au pays, ils retrouvent l'impôt, le travail forcé et l'humiliation, alors qu'ils ont vu des Européens mourir à leurs côtés comme des égaux. Leur pension est inférieure à celle des soldats français. Ils deviennent des relais naturels du nationalisme.
\item \textbf{La pression fiscale.} L'\cle{impôt de capitation} (impôt par tête) écrase les paysans. Pour le payer, il faut de l'argent, donc vendre ses récoltes à bas prix aux commerçants européens ou partir travailler dans les plantations. L'impôt est l'outil de l'exploitation économique.
\item \textbf{L'expropriation des terres.} Les grandes \textbf{plantations industrielles} (bananeraies, hévéas, palmiers à huile) s'approprient les meilleures terres, notamment dans le pays bamiléké et autour de Douala. Des villages entiers sont déplacés. La colère paysanne fournira à l'UPC ses bataillons les plus déterminés.
\item \textbf{Le racisme social.} Quartiers séparés, mépris dans les administrations, salaires inégaux à travail égal : la discrimination quotidienne radicalise les élites comme les travailleurs.
\item \textbf{Le rôle de l'Église.} Les missions catholiques et protestantes ont créé des écoles. Les \textbf{catéchistes} et les instituteurs, premiers lettrés des villages, lisent les journaux, rédigent les pétitions et deviennent les cadres locaux des partis.
\item \textbf{Le rôle des syndicats.} Les syndicats (liés à la CGT ou à Force Ouvrière en France) apprennent aux travailleurs l'organisation, la grève, la négociation : une école de militantisme.
\end{itemize}

\begin{attention}[Piège fréquent]
Ne réduisez jamais le nationalisme camerounais à l'UPC seule. Le nationalisme « \textbf{modéré} » (Ahidjo, Fouda, Betote Akoa et le BDC) revendique aussi l'émancipation du Cameroun, mais par la voie légale et l'autonomie progressive. C'est d'ailleurs ce courant qui conduit légalement le pays à l'indépendance du \textbf{1er janvier 1960}. Au Bac, un devoir qui oppose les deux courants avec nuance est toujours mieux noté qu'un récit héroïque à sens unique.
\end{attention}

\courssub{La naissance des partis politiques au Cameroun français}

\textbf{L'UPC (Union des Populations du Cameroun), 10 avril 1948.} Fondée à Douala, dirigée par \textbf{Ruben Um Nyobé}, l'UPC est un parti \textbf{radical} : il exige l'\textbf{indépendance immédiate} et la \textbf{réunification} des deux Cameroun. D'inspiration panafricaniste et proche des milieux communistes français, il refuse toute participation aux institutions coloniales qu'il juge factices. Son organisation en cellules de base (comités de village, de quartier, d'entreprise) en fait une machine politique redoutable.

\textbf{Le BDC (Bloc Démocratique Camerounais), 1949.} Créé autour de Louis-Paul Aujoulat et de notables comme André Fouda, le BDC est un parti \textbf{modéré} et pro-français : il accepte le cadre de l'Union française, participe aux élections et vise une autonomie interne progressive.

\textbf{Les Paysans Indépendants, 1951.} Mouvement rural modéré, il regroupe des planteurs (notamment du Centre et du Sud) hostiles à la radicalité de l'UPC et soucieux de défendre leurs intérêts économiques dans le cadre français.

\begin{propriete}[L'originalité de l'UPC]
L'UPC est le \textbf{premier parti nationaliste à structure nationale} — implanté dans les deux zones, française et britannique — et à \textbf{audience internationale} : il porte la cause camerounaise devant l'ONU et se lie aux mouvements panafricains naissants. Aucun autre parti camerounais de l'époque ne combine ces deux dimensions.
\end{propriete}

\begin{exemplebox}[Un parti de masse]
Les adhésions de l'UPC sont estimées à \textbf{300 000 militants en 1954} (source : rapports de la police coloniale, qui surveillait le parti de près). C'est alors le premier parti de masse d'Afrique noire francophone. Ce chiffre, même s'il doit être manié avec prudence (source policière, cartes parfois collectives), témoigne d'une implantation populaire exceptionnelle, surtout dans les pays bassa, bamiléké et à Douala.
\end{exemplebox}

\courssub{Au Cameroun britannique : deux options opposées}

Sous tutelle britannique, le débat est différent : il porte sur le \textbf{destin du territoire} une fois le Nigéria indépendant.

\begin{itemize}
\item \textbf{Le KNC (Kamerun National Congress), 1951}, dirigé par \textbf{John Ngu Foncha}, défend la \textbf{réunification} avec le Cameroun français pour former un seul État camerounais.
\item \textbf{Le KPP (Kamerun People's Party), 1953}, dirigé par \textbf{Emmanuel Endeley}, préfère l'\textbf{intégration au Nigéria} voisin (ou le maintien du statut quo), jugée plus favorable économiquement.
\end{itemize}

Ce clivage explique la question posée lors du plébiscite de février 1961, que nous étudierons dans la section 4.

\begin{popfigure}
\begin{center}
\begin{tikzpicture}[x=1cm,y=1cm,>=Stealth,font=\small]
% Axe chronologique proportionnel : 1945 à 1955 = 10 ans -> 14cm
\draw[->,thick,popdark] (0,0) -- (14.2,0);
% Graduations tous les 2 ans (2.8cm)
\foreach \x/\a in {0/1945, 2.8/1947, 5.6/1949, 8.4/1951, 11.2/1953, 14/1955}{
  \draw[thick,popdark] (\x,-0.12) -- (\x,0.12);
  \node[below,font=\bfseries,text=popdark] at (\x,-0.35) {\a};
}
% Événements haut (alternance haut/bas pour lisibilité)
% 1945 Grève Douala
\draw[popblue,thick] (0,0) -- (0,1.6);
\node[above,font=\scriptsize,text=popblue,align=center,text width=2.8cm] at (0,1.7) {Grève générale de Douala (sept. 1945)};
% 1947 Inde, Guerre froide
\draw[popblue,thick] (2.8,0) -- (2.8,2.4);
\node[above,font=\scriptsize,text=popblue,align=center,text width=2.8cm] at (2.8,2.5) {Indépendance de l'Inde (1947) ; Début Guerre froide};
% 1948 UPC (avril) -> x ~ 1.4
\draw[poporange,thick] (1.4,0) -- (1.4,1.6);
\node[above,font=\scriptsize,text=poporange,align=center,text width=2.6cm] at (1.4,1.7) {Création UPC (10 avril 1948)};
% 1949 BDC
\draw[popblue,thick] (4.2,0) -- (4.2,2.4);
\node[above,font=\scriptsize,text=popblue,align=center,text width=2.6cm] at (4.2,2.5) {Création BDC (1949)};
% 1951 KNC, Paysans Indép.
\draw[popgreen,thick] (8.4,0) -- (8.4,1.6);
\node[above,font=\scriptsize,text=popgreen,align=center,text width=2.8cm] at (8.4,1.7) {KNC (1951) ; Paysans Indépendants (1951)};
% 1952-54 Um Nyobé ONU
\draw[poporange,thick] (9.8,0) -- (9.8,2.4);
\node[above,font=\scriptsize,text=poporange,align=center,text width=2.8cm] at (9.8,2.5) {Um Nyobé à l'ONU (1952, 1953, 1954)};
% 1953 KPP
\draw[popgreen,thick] (11.2,0) -- (11.2,3.4);
\node[above,font=\scriptsize,text=popgreen,align=center,text width=2.6cm] at (11.2,3.5) {Création KPP (1953)};
% 1954 Diên Biên Phu, 300k adhérents
\draw[popmuted,thick] (12.6,0) -- (12.6,-1.4);
\node[below,font=\scriptsize,text=popmuted,align=center,text width=2.8cm] at (12.6,-1.5) {Diên Biên Phu (1954) ; 300 000 adhérents UPC};
% 1955 Bandung (avril)
\draw[popblue,thick] (13.3,0) -- (13.3,1.6);
\node[above,font=\scriptsize,text=popblue,align=center,text width=2.6cm] at (13.3,1.7) {Conférence de Bandung (avril 1955)};
% 1955 Émeutes mai, interdiction UPC (juillet)
\draw[poppink,thick] (14,0) -- (14,2.4);
\node[above,font=\scriptsize,text=poppink,align=center,text width=2.6cm] at (14,2.5) {Émeutes mai 1955 ; Interdiction UPC (13 juil. 1955)};
% Événements bas : Tutelle ONU déc 1946 -> x ~ 0.56
\draw[popmuted,thick] (0.56,0) -- (0.56,-1.0);
\node[below,font=\scriptsize,text=popmuted,align=center,text width=2.8cm] at (0.56,-1.1) {Tutelle ONU (déc. 1946)};
% Premières pétitions 1948-1950 ~ x 1.4 à 4.2
\draw[popmuted,thick] (2.8,0) -- (2.8,-1.0);
\node[below,font=\scriptsize,text=popmuted,align=center,text width=2.8cm] at (2.8,-1.1) {Premières pétitions UPC à l'ONU (1948-1950)};
\end{tikzpicture}
\end{center}
\legende{Figure 4 — Frise chronologique proportionnelle 1945-1955 : conférences internationales (bleu), créations de partis (orange, vert), répression (rose), contexte de tutelle (gris).}
\end{popfigure}

\begin{popfigure}
\begin{center}
\begin{tikzpicture}[font=\small,>=Stealth,
  racine/.style={draw=popdark,fill=popcream,rounded corners,align=center,text width=4.2cm,font=\small\bfseries},
  rouge/.style={draw=poppink,fill=poppinkL,rounded corners,align=center,text width=4.6cm},
  bleu/.style={draw=popblue,fill=popblueL,rounded corners,align=center,text width=4.6cm},
  vert/.style={draw=popgreen,fill=popgreenL,rounded corners,align=center,text width=4.6cm},
  orange/.style={draw=poporange,fill=poporangeL,rounded corners,align=center,text width=4.6cm}]
\node[racine] (r) at (0,0) {Nationalisme camerounais (1948-1960)};
\node[rouge, align=center] (upc) at (-6.5,-2.5){\textbf{UPC} (1948)\\ Um Nyobé\\ Indépendance immédiate, réunification, panafricanisme};
\node[bleu, align=center] (bdc) at (-1.5,-2.5){\textbf{BDC} (1949)\\ Aujoulat, Fouda\\ Autonomie dans l'Union française};
\node[bleu, align=center] (pi) at (3.5,-2.5){\textbf{Paysans Indépendants} (1951)\\ Défense des planteurs, modérés};
\node[vert, align=center] (knc) at (-4.0,-5.0){\textbf{KNC} (1951) — Foncha\\ Réunification des deux Cameroun};
\node[orange, align=center] (kpp) at (2.0,-5.0){\textbf{KPP} (1953) — Endeley\\ Intégration au Nigéria / statut quo};
\draw[->,thick,popdark] (r) -- (upc);
\draw[->,thick,popdark] (r) -- (bdc);
\draw[->,thick,popdark] (r) -- (pi);
\draw[->,thick,popdark] (r.south) |- (knc.north);
\draw[->,thick,popdark] (r.south) |- (kpp.north);
% Légende couleurs à droite, espacée
\node[rouge,text width=3.8cm,font=\scriptsize] at (7.5,-2.0) {Rouge : nationalistes radicaux (indépendance immédiate)};
\node[bleu,text width=3.8cm,font=\scriptsize] at (7.5,-3.2) {Bleu : modérés / pro-français (autonomie progressive)};
\node[vert,text width=3.8cm,font=\scriptsize] at (7.5,-4.4) {Vert : fédéralistes (réunification, zone britannique)};
\node[orange,text width=3.8cm,font=\scriptsize] at (7.5,-5.6) {Orange : intégrationnistes (rattachement au Nigéria)};
\end{tikzpicture}
\end{center}
\legende{Figure 5 — Organigramme des partis politiques camerounais (1948-1960) et leurs options idéologiques.}
\end{popfigure}

\courssub{L'apport des syndicats et la répression coloniale}

Les syndicats sont la pépinière du nationalisme. La \textbf{grève générale de Douala (septembre 1945)}, menée par les cheminots et les dockers, paralyse le port et le rail : elle montre que les travailleurs camerounais peuvent faire plier l'administration. De cette expérience naît l'\textbf{USC (Union des Syndicats Confédérés du Cameroun)}, étroitement liée à l'UPC : les syndicats fournissent au parti ses organisateurs, ses finances et ses réseaux dans les villes.

La réponse coloniale est la \textbf{répression} : arrestations de meneurs, licenciements, \textbf{déportations} de militants vers des camps d'internement comme la « Ferme Modèle » de Dschang, où des nationalistes sont détenus sans procès dans des conditions très dures. Loin d'étouffer le mouvement, cette répression nourrit la colère et grossit les rangs de l'UPC.

\courssub{La pétition à l'ONU : la voix du Cameroun à New York}

Le statut de tutelle offre une arme juridique inédite : le \textbf{droit de pétition} devant le Conseil de tutelle de l'ONU. L'UPC l'utilise massivement — des milliers de pétitions affluent à New York. Mieux : \textbf{Ruben Um Nyobé} se rend en personne devant l'Assemblée générale des Nations unies en \textbf{1952, 1953 et 1954} pour y défendre le « droit des peuples à disposer d'eux-mêmes » et dénoncer les violations de l'accord de tutelle par la France. L'impact médiatique est considérable : pour la première fois, un Camerounais expose la cause de son peuple devant le monde entier, et la France est mise en accusation sur la scène internationale.

\begin{savaistu}[« Mpodol » à la tribune de l'ONU]
Ruben Um Nyobé est surnommé « \textbf{Mpodol} », c'est-à-dire en langue bassa « celui qui parle pour le peuple ». Lors de son célèbre discours à l'ONU en 1952, il s'exprime d'abord en \textbf{bassa}, sa langue maternelle, avant de se traduire lui-même en français : un geste politique fort, affirmant la dignité des langues africaines face à la langue du colonisateur.
\end{savaistu}

\begin{popfigure}
\begin{center}
\begin{tikzpicture}[font=\small,>=Stealth,scale=0.9]
% Cadre photo
\draw[fill=popcream,draw=popline] (0,0) rectangle (13,5.5);
% Estrade
\draw[fill=poporangeL,draw=poporange] (8.8,0.5) rectangle (12.2,1.5);
\node[font=\scriptsize,text=popdark] at (10.5,1.0) {Estrade};
% Orateur
\draw[fill=popdark] (10.5,1.85) circle (0.18);
\draw[thick,popdark] (10.5,1.67) -- (10.5,2.6);
\draw[thick,popdark] (10.5,2.3) -- (11.2,2.7);
\draw[thick,popdark] (10.5,2.3) -- (9.8,2.6);
\node[font=\scriptsize\bfseries,text=popdark,align=center] at (10.5,3.1) {L'orateur (1)};
% Drapeau / banderole
\draw[thick,poppink] (9.0,3.8) -- (12.0,3.8);
\node[font=\scriptsize\bfseries,text=poppink] at (10.5,4.15) {Banderole : « Indépendance — Réunification » (2)};
% Foule dense
\foreach \x/\y in {
  1.0/0.7, 2.0/0.6, 3.0/0.8, 4.0/0.65, 5.0/0.75, 6.0/0.6, 7.0/0.8,
  1.5/1.6, 2.5/1.7, 3.5/1.55, 4.5/1.7, 5.5/1.6, 6.5/1.7,
  2.0/2.5, 3.2/2.6, 4.4/2.5, 5.6/2.6, 6.8/2.5
}{
  \draw[fill=popblue] (\x,\y+0.35) circle (0.12);
  \draw[thick,popblue] (\x,\y+0.22) -- (\x,\y-0.2);
}
\node[font=\scriptsize,text=popblue,align=center] at (4.0,3.5) {Foule de militants et de paysans (3)};
% Femmes au premier rang
\draw[fill=poppink] (0.8,2.8) circle (0.12); \draw[thick,poppink] (0.8,2.67) -- (0.8,2.2);
\draw[fill=poppink] (1.4,2.9) circle (0.12); \draw[thick,poppink] (1.4,2.77) -- (1.4,2.3);
\draw[fill=poppink] (2.0,2.85) circle (0.12); \draw[thick,poppink] (2.0,2.72) -- (2.0,2.25);
\node[font=\scriptsize,text=poppink,align=center,text width=2.8cm] at (1.4,3.9) {Femmes militantes, en pagne, au premier rang (4)};
% Affluence villageoise
\node[font=\scriptsize,text=popmuted,align=center,text width=3.2cm] at (7.5,4.9) {Affluence venue des villages environnants, à pied ou à vélo (5)};
% Flèche nord symbolique
\draw[->,thick,popdark] (12.2,5.0) -- (12.2,5.4);
\node[font=\tiny,text=popdark] at (12.2,5.5) {N};
\end{tikzpicture}
\end{center}
\legende{Figure 6 — Analyse d'une photographie : meeting de l'UPC à Eseka (1953), reconstitution schématique. (1) L'orateur sur une estrade de fortune : le parti s'implante en milieu rural. (2) Les mots d'ordre écrits : indépendance et réunification. (3) Une foule dense de paysans : le caractère de masse du parti. (4) La présence active des femmes. (5) La mobilisation dépasse le cadre urbain : le nationalisme gagne les campagnes.}
\end{popfigure}

\courssub{Complément Bac : le rôle occulté des femmes}

Les manuels classiques racontent souvent le nationalisme comme une affaire d'hommes. C'est une erreur historique. Des femmes comme \textbf{Marthe Moumié}, militante de l'UPC, participent aux meetings, transportent les tracts, cachent les militants recherchés et financent le parti par le commerce. Dans l'Ouest, les \textbf{pétitions des femmes} bamiléké et bafoussam affluent vers l'ONU : elles dénoncent l'expropriation des terres et l'arrestation de leurs maris et fils. Après l'interdiction de l'UPC en 1955, ce sont souvent les femmes qui maintiennent les réseaux clandestins. Retenir ce fait, c'est comprendre que le nationalisme fut un mouvement \textbf{populaire et total}, engageant toute la société.

\begin{methode}[Comparer deux programmes politiques]
Au Bac, on peut vous demander de comparer le \emph{Manifeste de l'UPC} (1948) et le \emph{Programme du BDC} (1950). Procédez en trois étapes :
\begin{enumerate}
\item \textbf{Le vocabulaire} : relevez les mots-clés. L'UPC écrit « indépendance », « réunification », « évacuation des troupes » ; le BDC écrit « autonomie », « évolution », « Union française ».
\item \textbf{Les cibles sociales} : à qui s'adresse chaque texte ? L'UPC parle aux paysans expropriés, aux travailleurs, aux jeunes ; le BDC aux notables, aux élus, aux planteurs intégrés au commerce colonial.
\item \textbf{Le rapport à la France} : rupture et méfiance chez l'UPC ; coopération et gradualisme chez le BDC.
\end{enumerate}
Concluez en montrant que les deux partis partagent le même objectif final — l'émancipation du Cameroun — mais divergent sur la \textbf{voie} et le \textbf{rythme}.

\begin{center}
\begin{tabularx}{\linewidth}{|l|X|X|}\hline
\rowcolor{popblueL}
\textbf{Critère} & \textbf{UPC (Manifeste 1948)} & \textbf{BDC (Programme 1950)} \\ \hline
Objectif proclamé & Indépendance immédiate et totale & Autonomie progressive dans l'Union française \\ \hline
Réunification & Revendication centrale & Acceptée mais secondaire \\ \hline
Moyens d'action & Lutte de masse, pétitions ONU, refus institutions & Participation aux élections, négociation avec Paris \\ \hline
Base sociale & Paysans expropriés, ouvriers, jeunes, femmes & Notables, chefs traditionnels, planteurs intégrés \\ \hline
Rapport à la France & Rupture, méfiance, anti-impérialisme & Coopération, gradualisme, fidélité à l'Union française \\ \hline
\end{tabularx}
\end{center}
\end{methode}

\begin{exoresolu}[Analyse de document : Um Nyobé à l'ONU (1952)]
\textbf{Document.} « Nous demandons que soit fixée une date pour l'indépendance du Cameroun et la réunification des deux territoires. Le peuple camerounais ne veut plus être administré : il veut se gouverner lui-même. » (Ruben Um Nyobé, discours devant l'Assemblée générale de l'ONU, 1952, adapté.)

\textbf{Questions.} 1) Qui est l'auteur et dans quel cadre s'exprime-t-il ? 2) Quelles sont les deux revendications exprimées ? 3) Pourquoi ce discours a-t-il un retentissement international ?
\tcblower
\textbf{Solution détaillée.}
1) L'auteur est Ruben Um Nyobé, secrétaire général de l'UPC, principal parti nationaliste camerounais. Il s'exprime devant l'Assemblée générale de l'ONU, dans le cadre du droit de pétition prévu par le régime de tutelle institué en 1946 : le Cameroun n'est plus une colonie mais un territoire sous tutelle, ce qui autorise ses représentants à saisir directement l'ONU.
2) Les deux revendications sont : l'\textbf{indépendance} du Cameroun avec une date fixée à l'avance, et la \textbf{réunification} des deux territoires (Cameroun français et Cameroun britannique) en un seul État.
3) Le retentissement est international car, pour la première fois, la France est publiquement mise en accusation devant une instance mondiale pour sa gestion du territoire. En pleine Guerre froide et au lendemain des indépendances asiatiques, ce discours fait du Cameroun un symbole de la lutte anticoloniale et oblige la France à accélérer les réformes — tout en durcissant sa répression contre l'UPC, interdite en juillet 1955.
\end{exoresolu}

\begin{aretenir}[L'essentiel de la section]
Entre 1945 et 1955, le nationalisme camerounais naît de la rencontre entre un \textbf{contexte international favorable} (Charte de l'Atlantique, Brazzaville, décolonisation asiatique, Guerre froide, Bandung) et des \textbf{frustrations internes} (anciens combattants, impôt de capitation, expropriations, racisme). Il s'organise en partis : l'\textbf{UPC} (1948, radicale, Um Nyobé), le \textbf{BDC} (1949, modéré), les \textbf{Paysans Indépendants} (1951) ; au Cameroun britannique, le \textbf{KNC} (réunification) et le \textbf{KPP} (intégration au Nigéria). Soutenu par les syndicats (USC) et porté jusqu'à la tribune de l'ONU par Um Nyobé (1952-1954), le mouvement est réprimé : l'UPC est interdite le \textbf{13 juillet 1955}, ce qui ouvre la phase de radicalisation étudiée dans la section suivante.
\end{aretenir}

\coursec{Dossier 3 : Les grandes figures du nationalisme camerounais — Portraits croisés}

Dans la section précédente, nous avons vu comment, entre 1945 et 1955, le nationalisme camerounais s'est éveillé puis structuré : création de l'UPC en 1948, pétitions à l'ONU, syndicalisme à Douala, répression de 1955. Mais une histoire ne se résume pas à des dates et à des partis : elle est portée par des \cle{hommes et des femmes} qui font des choix, prennent des risques et paient parfois de leur vie leur engagement. Ce dossier propose des \cle{portraits croisés} des grandes figures du nationalisme camerounais. L'objectif n'est pas de classer les uns en « héros » et les autres en « traîtres », mais de comprendre la \cle{diversité des stratégies} — lutte armée, diplomatie internationale, action parlementaire, syndicalisme — qui ont toutes, à leur manière, conduit le Cameroun vers la souveraineté.

\courssub{Comment étudier une figure historique ?}

Avant d'aborder les portraits, il faut se doter d'un outil d'analyse commun. Une figure historique ne se juge pas seulement sur ses intentions, mais sur son parcours complet : d'où vient-elle, comment s'engage-t-elle, quelle stratégie choisit-elle, quelle action la rend célèbre, et quelle trace laisse-t-elle dans la mémoire nationale ?

\begin{methode}[Construire une fiche biographique « Figure historique »]
Pour analyser rigoureusement une grande figure du nationalisme, suivre cinq étapes :
\begin{enumerate}
\item \textbf{Origines et formation} : date et lieu de naissance, milieu social, études, métier.
\item \textbf{Engagement politique} : date d'entrée en politique, parti ou syndicat, fonctions occupées.
\item \textbf{Stratégie et idéologie} : légalisme ou lutte armée, indépendance immédiate ou progressive, unité ou fédéralisme.
\item \textbf{Action majeure} : l'événement qui fait d'elle une figure historique (discours, bataille, vote, négociation).
\item \textbf{Postérité et héritage} : sort final (mort, exil, pouvoir), place dans la mémoire nationale, réhabilitation éventuelle.
\end{enumerate}
Cette grille sera appliquée à chaque portrait de ce dossier et pourra être réutilisée à l'examen.
\end{methode}

\courssub{Les figures de la lutte armée et de la diplomatie clandestine : l'UPC}

\textbf{Ruben Um Nyobé (1913-1958), le « père de l'indépendance ».}\par
Né en 1913 à Song Mpeck (région du Centre, département du Sanaga-Maritime), Ruben Um Nyobé est un petit fonctionnaire devenu syndicaliste. Élu secrétaire général de l'\cle{UPC} en 1948, il en fait le principal parti nationaliste. Sa stratégie est d'abord \cle{légaliste} : il porte trois fois la revendication d'indépendance et de réunification devant l'ONU à New York (1952, 1953, 1954), avec une éloquence qui impressionne les diplomates. Mais après l'interdiction de l'UPC par la France en juillet 1955, il choisit la \cle{lutte armée} et entre dans le maquis de la forêt du Sanaga-Maritime. Traqué, il est \cle{abattu le 13 septembre 1958} par l'armée française près de Boumnyebel. Longtemps effacé de la mémoire officielle, il est aujourd'hui reconnu comme le \cle{père de l'indépendance} camerounaise.

\textbf{Félix-Roland Moumié (1925-1960), le diplomate assassiné.}\par
Médecin de formation, Félix-Roland Moumié succède à Um Nyobé à la tête de l'UPC en 1958. Depuis l'exil, il mène une intense \cle{diplomatie internationale} : il obtient des soutiens en Chine, au Ghana de Nkrumah, en Guinée de Sékou Touré et auprès du FLN algérien. Cette présence sur la scène mondiale inquiète la France, qui vient d'accorder l'indépendance au Cameroun. Le 3 novembre 1960, Moumié est \cle{empoisonné au thallium} dans un restaurant de Genève par un agent des services secrets français (le SDECE). Il meurt le lendemain, à 35 ans. Son assassinat, reconnu officiellement par la France des décennies plus tard, illustre la violence de la guerre secrète menée contre l'UPC.

\textbf{Ernest Ouandié (1924-1971), le dernier résistant.}\par
Chef militaire de l'UPC, Ernest Ouandié poursuit la lutte armée \emph{après} l'indépendance de 1960, considérant que celle-ci est une « indépendance confisquée » au profit d'un pouvoir soutenu par la France. Traqué pendant plus de dix ans dans les maquis de l'Ouest, il est capturé en 1970, jugé lors d'un procès expéditif et \cle{exécuté publiquement} à Bafoussam le 15 janvier 1971. Sa mort marque la fin de la rébellion upéciste et fait de lui le symbole de la \cle{résistance post-indépendance}.

\begin{definition}[Le maquis]
Le \cle{maquis} désigne la guérilla nationaliste menée par l'UPC dans les forêts du Sanaga-Maritime, du pays bamiléké et du Mungo, entre 1955 et 1971. Les maquisards, cachés dans la brousse et soutenus par une partie des populations rurales, attaquent les symboles de l'administration coloniale puis du jeune État camerounais. Le maquis est réprimé d'abord par l'armée française, puis par l'armée camerounaise avec l'aide de conseillers français.
\end{definition}

\begin{exemplebox}[Le coût humain de la répression du maquis]
Selon les estimations des historiens, la répression du maquis upéciste entre 1957 et 1960 aurait fait entre \cle{10 000 et 100 000 morts}, selon les sources. Des centaines de villages ont été détruits et leurs habitants déplacés de force dans le cadre de la politique de \cle{regroupement forcé} (ou « villagisation »), destinée à couper les maquisards de leurs soutiens. Ces chiffres, longtemps tabous, montrent que la marche vers l'indépendance camerounaise fut l'une des plus sanglantes d'Afrique coloniale.
\end{exemplebox}

\begin{savaistu}[Un « gouvernement en exil » reconnu à l'étranger]
En 1958, Félix-Roland Moumié obtient du Ghana de Kwame Nkrumah l'ouverture de la première représentation diplomatique de l'UPC à l'étranger. Grâce à ses voyages, l'UPC devient une sorte de \cle{gouvernement en exil}, reconnu et soutenu par plusieurs pays du bloc de l'Est (Chine, URSS) et du mouvement des non-alignés (Ghana, Guinée, Égypte). C'est une preuve que la lutte pour l'indépendance du Cameroun s'inscrivait pleinement dans la guerre froide.
\end{savaistu}

\courssub{Les figures de la voie légale et parlementaire}

\textbf{André-Marie Mbida (1917-1980), le premier Premier ministre.}\par
Catholique, modéré, André-Marie Mbida devient en 1957 le \cle{premier Premier ministre africain} du Cameroun, à la tête du premier gouvernement autonome. Hostile à l'UPC et partisan d'une indépendance progressive en lien avec la France, il se heurte à son propre camp : en février 1958, il est renversé par une \cle{motion de censure} de l'Assemblée territoriale (l'« affaire Mbida »). Il devient ainsi le premier chef de gouvernement africain destitué par son parlement — un événement qui montre la précocité de la vie parlementaire camerounaise.

\textbf{Ahmadou Ahidjo (1924-1989), l'artisan de l'indépendance négociée.}\par
Né à Garoua, Ahmadou Ahidjo est député à l'Assemblée française avant de devenir Premier ministre en février 1958, puis \cle{premier Président de la République} à l'indépendance du 1er janvier 1960. Sa stratégie est celle de l'\cle{indépendance négociée} avec la France, dans la légalité institutionnelle. Il est aussi l'artisan, avec Foncha, de la \cle{réunification} de 1961 et le partisan d'une unité forte, aboutissant à l'\cle{État unitaire} de 1972. Son parcours illustre la voie parlementaire et étatique vers la souveraineté.

\textbf{John Ngu Foncha (1916-1999), le père de la réunification.}\par
Instituteur anglophone, John Ngu Foncha devient Premier ministre du \cle{Southern Cameroons} (Cameroun britannique) en 1959. Convaincu que l'avenir de sa région est aux côtés du Cameroun francophone, il fait campagne pour le « OUI » lors du \cle{référendum du 11 février 1961} organisé sous l'égide de l'ONU. Après la réunification du 1er octobre 1961, il devient \cle{Vice-président} de la République fédérale aux côtés d'Ahidjo. Il démissionnera en 1970, déçu par l'évolution vers l'État unitaire.

\textbf{Les autres figures.}\par
D'autres acteurs ont joué un rôle décisif : \cle{Charles Assalé}, Premier ministre du Cameroun oriental dans la République fédérale ; \cle{Vincent de Paul Ahanda}, ministre et diplomate ; \cle{Solomon Tandeng Muna}, successeur de Foncha à la tête du Cameroun occidental puis Vice-président ; et \cle{Paul Soppo Priso}, syndicaliste de Douala, figure du nationalisme modéré et de l'action sociale.

\begin{propriete}[La réunification : une convergence, pas l'œuvre d'un seul homme]
La réunification de 1961 n'est pas l'œuvre d'un seul leader. Elle résulte de la \cle{convergence tactique} entre Ahmadou Ahidjo, qui apporte la \cle{légalité étatique} (le Cameroun indépendant, reconnu par l'ONU), et John Ngu Foncha, qui apporte la \cle{légitimité électorale} (le vote du Southern Cameroons), le tout sous le parapluie de l'ONU qui organise et contrôle le référendum. Sans l'un de ces trois éléments, la réunification n'aurait pas eu lieu.
\end{propriete}

\begin{attention}[Le piège historiographique : « collabos » contre « vrais nationalistes »]
Il est tentant mais faux d'opposer artificiellement des « collaborateurs » (Ahidjo, Mbida) à des « vrais nationalistes » (Um Nyobé, Moumié). Pendant des décennies, l'histoire officielle a gommé le rôle de l'UPC et de ses martyrs ; à l'inverse, une lecture simpliste réduit Ahidjo à un « homme de la France ». L'histoire critique actuelle réhabilite la \cle{complexité des parcours} : tous voulaient la souveraineté du Cameroun, mais ils divergeaient sur les \cle{moyens} (lutte armée ou négociation), le \cle{calendrier} (immédiat ou progressif) et les \cle{alliances} internationales. À l'examen, évite tout jugement manichéen et analyse les stratégies.
\end{attention}

\courssub{Galerie de portraits : une vue d'ensemble}

\begin{popfigure}
\begin{center}
\begin{tabularx}{\linewidth}{|l|X|X|l|}
\hline
\rowcolor{popblueL} \textbf{Figure} & \textbf{Dates / Parti} & \textbf{Rôle principal} & \textbf{Sort final} \\
\hline
Ruben Um Nyobé & 1913-1958 / UPC & Secrétaire général, porte-parole à l'ONU, chef du maquis & Mort au combat (1958) \\
\hline
Félix-Roland Moumié & 1925-1960 / UPC & Président de l'UPC, diplomate en exil & Assassiné à Genève (1960) \\
\hline
Ernest Ouandié & 1924-1971 / UPC & Chef militaire du maquis & Exécuté (1971) \\
\hline
Ahmadou Ahidjo & 1924-1989 / UC puis UNC & Premier Président, indépendance négociée, État unitaire & Pouvoir puis exil (1982) \\
\hline
John Ngu Foncha & 1916-1999 / KNDP & Artisan du « OUI » à la réunification, Vice-président fédéral & Retraite politique (1970) \\
\hline
André-Marie Mbida & 1917-1980 / Démocrates & Premier Premier ministre africain (1957-1958) & Mort naturelle (1980) \\
\hline
\end{tabularx}
\end{center}
\legende{Figure 7 — Galerie de portraits des grandes figures du nationalisme camerounais : dates, parti, rôle principal et sort final. On remarque la diversité des destins : mort au combat, assassinat, exécution, pouvoir, retraite.}
\end{popfigure}

\courssub{Une carte mentale : quatre voies vers l'indépendance}

Pour mémoriser ces portraits, il est utile de classer les figures non pas par « camps », mais par \cle{type d'action}. Quatre grandes voies convergent vers l'indépendance : la diplomatie internationale, la lutte armée, l'action parlementaire et le syndicalisme.

\begin{popfigure}
\begin{center}
\begin{tikzpicture}[mindmap, grow cyclic, every node/.style={concept, text=white, font=\small}, concept color=popblue, level 1/.append style={level distance=3.6cm, sibling angle=90}, level 2/.append style={level distance=2.6cm, sibling angle=45, concept color=poporange, font=\scriptsize}]
\node[concept color=popdark, font=\small\bfseries]{Indépendance du Cameroun}
  child[concept color=popblue]{ node {Diplomatie (ONU, exil)} 
      child { node {Um Nyobé (New York 1952-1954)} }
      child { node {Moumié (Accra, Conakry)} } }
  child[concept color=popgreen]{ node {Maquis (forêt)} 
      child { node {Um Nyobé (Sanaga-Maritime)} }
      child { node {Ouandié (Ouest)} } }
  child[concept color=popgold]{ node {Parlement (Yaoundé, Buéa)} 
      child { node {Ahidjo (indépendance négociée)} }
      child { node {Foncha (réunification)} }
      child { node {Mbida (1\textsuperscript{er} gouvernement)} } }
  child[concept color=poppurple]{ node {Syndicat (Douala)} 
      child { node {Soppo Priso} } };
\end{tikzpicture}
\end{center}
\legende{Figure 8 — Carte mentale : les quatre voies d'action des figures du nationalisme camerounais. La diplomatie à l'ONU, le maquis dans les forêts, le travail parlementaire à Yaoundé et Buéa, et le syndicalisme à Douala convergent vers un même objectif : la souveraineté.}
\end{popfigure}

\courssub{Analyser les discours des figures : une méthode appliquée}

Pour comprendre les stratégies des leaders, rien ne vaut l'analyse de leurs propres paroles. Deux discours fondateurs peuvent être comparés : le discours d'Ahmadou Ahidjo du 1er janvier 1960 (proclamation de l'indépendance du Cameroun oriental) et le discours de John Ngu Foncha du 1er octobre 1961 (réunification à Buéa). Une capsule vidéo avec extraits sonores est proposée en complément (code QR fourni par l'enseignant).

\begin{popfigure}
\begin{center}
\begin{tikzpicture}[scale=0.95]
\draw[fill=popblueL, rounded corners] (0,0) rectangle (6.6,3.4);
\node[font=\small\bfseries, text=popdark] at (3.3,3.0) {Ahidjo, 1\textsuperscript{er} janvier 1960};
\node[font=\scriptsize, text width=6cm, align=left] at (3.3,1.5) {Thèmes : souveraineté retrouvée,\\ ordre, travail, fidélité aux\\ engagements internationaux.\\ Ton : solennel, rassurant,\\ tourné vers l'État à construire.};
\draw[fill=poporangeL, rounded corners] (7.4,0) rectangle (14,3.4);
\node[font=\small\bfseries, text=popdark] at (10.7,3.0) {Foncha, 1\textsuperscript{er} octobre 1961};
\node[font=\scriptsize, text width=6cm, align=left] at (10.7,1.5) {Thèmes : fraternité retrouvée,\\ fin de la séparation coloniale,\\ espoir de l'unité.\\ Ton : ému, historique, tourné\\ vers le peuple réuni.};
\draw[-{Stealth}, very thick, popdark] (3.3,-0.4) -- (3.3,-1.4);
\draw[-{Stealth}, very thick, popdark] (10.7,-0.4) -- (10.7,-1.4);
\draw[fill=popgreenL, rounded corners] (3.5,-2.6) rectangle (10.5,-1.6);
\node[font=\small\bfseries, text=popdark] at (7,-2.1) {Convergence : la République fédérale du Cameroun};
\end{tikzpicture}
\end{center}
\legende{Figure 9 — Analyse rhétorique comparée des discours d'Ahidjo (1960) et de Foncha (1961) : deux styles, deux légitimités, une même destination. Extraits vidéo consultables via le code QR fourni en classe.}
\end{popfigure}

\begin{exoresolu}[Analyse de document : le discours de Foncha du 1er octobre 1961]
\textbf{Énoncé.} Document : extrait du discours de John Ngu Foncha à Buéa, le 1er octobre 1961 : « Aujourd'hui, les deux Cameroun, séparés par l'histoire coloniale depuis plus de quarante ans, redeviennent un seul peuple. Ce jour marque la fin d'une longue attente et le début d'une nation unie. » Question : dégagez l'idée principale de ce discours et montrez en quoi il traduit la stratégie politique de son auteur.
\tcblower
\textbf{Solution détaillée.}
\begin{enumerate}
\item \textbf{Idée principale} : Foncha célèbre la réunification des deux Cameroun comme la fin d'une séparation artificielle imposée par la colonisation et la naissance d'une nation une.
\item \textbf{Stratégie de l'auteur} : en parlant de « séparation par l'histoire coloniale », Foncha présente la réunification comme un retour à l'ordre naturel, ce qui légitime le vote « OUI » du référendum du 11 février 1961 dont il fut le principal artisan.
\item \textbf{Conclusion} : ce discours traduit la stratégie de Foncha, fondée sur la légitimité électorale et la voie légale sous contrôle de l'ONU, par opposition à la lutte armée de l'UPC. Il confirme que la réunification est une convergence entre la légalité étatique d'Ahidjo et la légitimité populaire de Foncha.
\end{enumerate}
\end{exoresolu}

\begin{aretenir}[L'essentiel du dossier]
\begin{itemize}
\item \cle{Ruben Um Nyobé} (1913-1958) : secrétaire général de l'UPC, porte-parole à l'ONU puis chef du maquis, « père de l'indépendance ».
\item \cle{Félix-Roland Moumié} (1925-1960) : président de l'UPC en exil, diplomate, assassiné au thallium à Genève par le SDECE.
\item \cle{Ernest Ouandié} (1924-1971) : dernier chef militaire du maquis, exécuté publiquement à Bafoussam.
\item \cle{Ahmadou Ahidjo} (1924-1989) : artisan de l'indépendance négociée (1960), de la réunification et de l'État unitaire (1972).
\item \cle{John Ngu Foncha} (1916-1999) : artisan du vote « OUI » de 1961, Vice-président fédéral.
\item \cle{André-Marie Mbida} (1917-1980) : premier Premier ministre africain (1957), renversé par motion de censure.
\item Le nationalisme camerounais a emprunté \cle{quatre voies complémentaires} : diplomatie, maquis, parlement, syndicat. Il faut éviter le jugement manichéen et analyser la complexité des parcours.
\end{itemize}
\end{aretenir}

\begin{exercice}[Pour s'entraîner]
\begin{enumerate}
\item Construis, à l'aide de la méthode en cinq étapes, la fiche biographique complète de Félix-Roland Moumié.
\item Explique pourquoi il est réducteur d'opposer « vrais » et « faux » nationalistes camerounais. Appuie ta réponse sur deux portraits du dossier.
\item À l'aide de la carte mentale (figure 8), rédige un paragraphe montrant que l'indépendance du Cameroun est le résultat de stratégies complémentaires.
\end{enumerate}
\end{exercice}

\begin{corrige}[Exercice « Pour s'entraîner »]
\begin{enumerate}
\item Fiche de Moumié : (1) né en 1925, médecin ; (2) militant UPC, président du parti après la mort d'Um Nyobé en 1958 ; (3) stratégie : diplomatie internationale depuis l'exil (Chine, Ghana, Guinée, Algérie) ; (4) action majeure : obtention de la première représentation diplomatique de l'UPC à Accra ; (5) postérité : assassiné au thallium à Genève en 1960 par le SDECE, martyr du nationalisme.
\item L'opposition est réductrice car Ahidjo et Foncha, par la voie légale, ont obtenu l'indépendance (1960) et la réunification (1961), tandis que l'UPC, par la lutte armée et la diplomatie, a porté la revendication dès 1948 et contraint la France à accélérer le processus. Les deux stratégies visaient la souveraineté, par des moyens différents.
\item Paragraphe attendu : la diplomatie (Um Nyobé à l'ONU, Moumié en exil) a internationalisé la question camerounaise ; le maquis (Ouandié) a montré la détermination d'une partie du peuple ; le parlement (Ahidjo, Mbida, Foncha) a construit les institutions de la souveraineté ; le syndicat (Soppo Priso) a mobilisé les travailleurs. Ces quatre voies convergent vers l'indépendance de 1960 et la réunification de 1961.
\end{enumerate}
\end{corrige}

\coursec{De la revendication à la souveraineté : Étapes clés (1955-1961)}

Après avoir rencontré les grandes figures du nationalisme camerounais — Um Nyobè, Moumié, Ahidjo, Foncha —, il faut maintenant suivre le chemin concret que leurs luttes ont tracé : comment passe-t-on, en six ans seulement, d'un territoire sous tutelle à un État indépendant et réunifié ? Cette section reconstitue, étape par étape, la mécanique politique qui a conduit de l'insurrection de 1955 à la naissance de la République Fédérale du Cameroun en 1961. Retiens bien la logique d'ensemble : \cle{répression} du nationalisme radical, \cle{réformes} accordées par la France, \cle{élections}, puis \cle{indépendance} et \cle{réunification}.

\courssub{l'insurrection de l'UPC et le début de la guerre cachée}

L'année 1955 est un tournant tragique. L'UPC, parti le plus populaire du Cameroun français, réclame l'indépendance immédiate et la réunification des deux Cameroun. Face au blocage de l'administration française et à la répression des militants, des émeutes éclatent en \cle{mai 1955} à Douala, Yaoundé et dans la Sanaga-Maritime : ce sont les \cle{événements de mai 1955}, qui font plusieurs dizaines de morts.

La réponse de la France est radicale : le \cle{13 juillet 1955}, l'UPC et ses organisations amies (JDC, UDEFEC, OFUC) sont \cle{interdites} par décret. Les leaders du parti, dont Ruben Um Nyobè, refusent de se rendre. Une partie se réfugie dans le \cle{maquis} (forêts de la Sanaga-Maritime et du pays bassa) : c'est le début de la \cle{guerre cachée}, un conflit armé entre les maquisards upécistes et l'armée française, qui durera plusieurs années. D'autres dirigeants passent la frontière et s'exilent à \cle{Kumba}, au Cameroun britannique, avant de gagner l'étranger (Guinée, Ghana, Égypte).

\begin{attention}[Ne pas confondre les dates]
L'interdiction de l'UPC date de \textbf{juillet 1955} (au Cameroun français), pas de 1956. Et l'UPC n'a jamais disparu : interdite, elle poursuit la lutte depuis le maquis et l'exil, ce qui explique que l'indépendance de 1960 soit contestée par une partie des nationalistes.
\end{attention}

\courssub{la loi-cadre Defferre, la grande réforme}

Pour calmer les tensions dans tout son empire, la France adopte le \cle{23 juin 1956} la loi-cadre proposée par le ministre Gaston Defferre.

\begin{definition}[La loi-cadre Defferre (1956)]
Texte fondateur de la décolonisation française en Afrique. Elle accorde aux territoires d'outre-mer et sous tutelle :
\begin{itemize}
\item l'\cle{autonomie interne} : chaque territoire se dote d'un gouvernement responsable devant une assemblée élue ;
\item le \cle{suffrage universel} : tous les adultes votent, sans distinction de statut ;
\item la fin définitive du \cle{code de l'indigénat} et du collège électoral discriminatoire.
\end{itemize}
Au Cameroun, elle entraîne la création de l'\cle{Assemblée législative du Cameroun} (ALCAM).
\end{definition}

L'intuition est simple : plutôt que d'affronter les nationalismes partout, la France choisit de former des élites africaines « raisonnables » à qui elle transfère progressivement le pouvoir, tout en gardant la main sur la défense et les affaires étrangères. La loi-cadre est donc à la fois une avancée démocratique réelle et un instrument de contrôle.

\courssub{1957-1958 : des élections à l'indépendance proclamée}

Les \cle{élections législatives du 23 décembre 1957} donnent la victoire à l'Union camerounaise (UC) d'\cle{Ahmadou Ahidjo} et à ses alliés. Mais c'est \cle{André-Marie Mbida} qui devient Premier ministre : il est le \cle{premier Africain} à diriger un gouvernement camerounais.

Mbida se montre réticent à réclamer l'indépendance immédiate et hostile à la réunification. Il perd rapidement sa majorité : en \cle{février 1958}, son gouvernement est renversé par l'ALCAM, et \cle{Ahmadou Ahidjo} devient Premier ministre. L'année 1958 voit ensuite deux événements décisifs :

\begin{itemize}
\item le \cle{28 septembre 1958}, le référendum constitutionnel français (proposition de la Communauté du général de Gaulle) : le Cameroun vote massivement \cle{« OUI »}, contrairement à la Guinée de Sékou Touré ;
\item le \cle{12 novembre 1958}, l'ALCAM adopte une résolution proclamant la volonté d'indépendance du Cameroun, fixée au \cle{1\textsuperscript{er} janvier 1960}.
\end{itemize}

\begin{aretenir}[Les dates clés 1955-1960]
\begin{itemize}
\item Mai 1955 : insurrection UPC ; juillet 1955 : interdiction de l'UPC.
\item 23 juin 1956 : loi-cadre Defferre.
\item 23 décembre 1957 : élections, Mbida Premier ministre.
\item Février 1958 : Ahidjo Premier ministre ; 28 septembre 1958 : « OUI » à la Communauté ; 12 novembre 1958 : proclamation de l'indépendance pour le 1\textsuperscript{er} janvier 1960.
\item 1\textsuperscript{er} janvier 1960 : indépendance du Cameroun français.
\end{itemize}
\end{aretenir}

\courssub{1\textsuperscript{er} janvier 1960 : l'indépendance du Cameroun français}

Le \cle{1\textsuperscript{er} janvier 1960}, la tutelle française prend fin : naît la \cle{République du Cameroun}, avec \cle{Ahmadou Ahidjo} comme premier Président. Le nouvel État adhère à l'\cle{ONU le 20 septembre 1960}, ce qui consacre sa souveraineté internationale.

Mais cette indépendance est incomplète à deux titres : d'une part, la guerre contre les maquis upécistes se poursuit (Um Nyobè est tué en septembre 1958, mais la rébellion continue) ; d'autre part, la question du \cle{Cameroun britannique} reste entière. Or la Charte de l'UPC et l'ONU exigeaient la \cle{réunification} des deux territoires.

\courssub{Le sort du Cameroun britannique : le référendum du 11 février 1961}

Sous tutelle britannique, administré avec le Nigéria voisin, le Cameroun britannique doit choisir son destin. Après la \cle{conférence de Londres (1959)}, l'ONU organise un plébiscite le \cle{11 février 1961}. Deux questions sont posées aux électeurs, séparément dans le Nord et le Sud : souhaitez-vous être rattaché au \cle{Nigéria} indépendant, ou à la \cle{République du Cameroun} ?

\begin{exemplebox}[Les résultats chiffrés du référendum de 1961]
Participation : environ 85~\%.
\begin{itemize}
\item \textbf{Nord} (Adamawa, Bénoué) : 103 311 voix pour le Nigéria (environ 60~\%) contre 69 092 pour la réunification (environ 40~\%).
\item \textbf{Sud} (Grassfields, forêt) : 145 996 voix pour la réunification (70,5~\%) contre 61 092 pour le Nigéria (29,5~\%).
\end{itemize}
Conséquence : le Nord britannique rejoint le Nigéria (1\textsuperscript{er} juin 1961), le Sud britannique rejoint le Cameroun (1\textsuperscript{er} octobre 1961).
\end{exemplebox}

\begin{attention}[Comprendre le vote du Nord sans caricature]
Le Nord britannique (Adamawa, Bénoué) a choisi le Nigéria pour des raisons \cle{religieuses} (islam partagé avec le Nord-Nigéria), \cle{économiques} (commerce du bétail orienté vers Lagos) et \cle{politiques} (influence des émirs pro-nigérians). Ce n'est pas un rejet du Cameroun, mais un \cle{choix pragmatique} fondé sur des intérêts concrets.
\end{attention}

\begin{popfigure}
\begin{center}
\begin{tikzpicture}[scale=0.9, font=\small]
% Cameroun (français)
\fill[popblueL] (0,0) -- (0,5.5) -- (2.2,5.8) -- (3.4,4.6) -- (3.0,2.6) -- (3.4,0.8) -- (2.4,-0.6) -- (0.6,-0.4) -- cycle;
\draw[popblue, thick] (0,0) -- (0,5.5) -- (2.2,5.8) -- (3.4,4.6) -- (3.0,2.6) -- (3.4,0.8) -- (2.4,-0.6) -- (0.6,-0.4) -- cycle;
\node[popdark] at (1.7,2.6) {\textbf{République}};
\node[popdark] at (1.7,2.1) {\textbf{du Cameroun}};
\node[popdark, font=\scriptsize] at (1.7,1.6) {(ex-français, 1960)};
\node[popink] at (1.2,0.4) {$\bullet$ Yaoundé};
% Cameroun britannique Sud
\fill[popgreenL] (-2.6,0.6) -- (-2.6,3.0) -- (-0.4,3.2) -- (-0.2,0.8) -- cycle;
\draw[popgreen, thick] (-2.6,0.6) -- (-2.6,3.0) -- (-0.4,3.2) -- (-0.2,0.8) -- cycle;
\node[popdark, font=\scriptsize] at (-1.5,2.6) {\textbf{Sud britannique}};
% Camembert Sud : 70.5% réunification (vert), 29.5% orange
\fill[popgreen] (-1.5,1.7) circle (0.55);
\fill[poporange] (-1.5,1.7) -- ++(90:0.55) arc (90:90-106:0.55) -- cycle;
\node[white, font=\scriptsize\bfseries] at (-1.5,1.55) {70,5\%};
% Cameroun britannique Nord
\fill[poporangeL] (-2.6,3.2) -- (-2.6,5.6) -- (-0.4,5.8) -- (-0.4,3.4) -- cycle;
\draw[poporange, thick] (-2.6,3.2) -- (-2.6,5.6) -- (-0.4,5.8) -- (-0.4,3.4) -- cycle;
\node[popdark, font=\scriptsize] at (-1.5,5.3) {\textbf{Nord britannique}};
% Camembert Nord : 60% Nigéria (orange), 40% vert
\fill[poporange] (-1.5,4.4) circle (0.55);
\fill[popgreen] (-1.5,4.4) -- ++(90:0.55) arc (90:90-144:0.55) -- cycle;
\node[white, font=\scriptsize\bfseries] at (-1.5,4.25) {60\%};
% Nigéria
\fill[popcream] (-6.4,0.4) rectangle (-2.8,5.8);
\draw[popmuted] (-6.4,0.4) rectangle (-2.8,5.8);
\node[popmuted] at (-4.6,3.1) {\textbf{NIGÉRIA}};
% Ligne de partage Nord/Sud
\draw[popink, very thick, dashed] (-2.7,3.1) -- (-0.1,3.3);
% Flèches de rattachement
\draw[-{Stealth[length=4mm]}, poporange, very thick] (-2.7,4.5) -- (-3.6,4.5);
\node[poporange, font=\scriptsize, align=center] at (-4.6,4.9) {1\textsuperscript{er} juin 1961\\vers le Nigéria};
\draw[-{Stealth[length=4mm]}, popgreen, very thick] (-0.1,1.9) -- (0.7,1.9);
\node[popgreen!60!black, font=\scriptsize, align=center] at (1.9,3.6) {1\textsuperscript{er} oct. 1961\\réunification};
% Nord
\draw[-{Stealth}, popdark] (4.2,5.2) -- (4.2,6.0) node[above, font=\scriptsize]{N};
\end{tikzpicture}
\end{center}
\legende{Figure 10 — Le référendum du 11 février 1961. Le Nord britannique (orange) vote à 60~\% pour le Nigéria ; le Sud britannique (vert) vote à 70,5~\% pour la réunification avec la République du Cameroun. La ligne pointillée matérialise le partage Nord/Sud. Croquis schématique.}
\end{popfigure}

\courssub{Foumban et la naissance de la République Fédérale (1961)}

Pour préparer la réunification, Ahidjo et \cle{John Ngu Foncha}, leader du Sud britannique, négocient une constitution commune : c'est la \cle{Constitution de Foumban} (août 1961). Le \cle{1\textsuperscript{er} octobre 1961}, la réunification devient effective : naît la \cle{République Fédérale du Cameroun}, avec Ahidjo comme Président fédéral et Foncha comme Vice-Président.

\begin{definition}[La Fédération (1961-1972)]
État composé de deux \cle{États fédérés} : le \cle{Cameroun oriental} (ex-français, capitale Yaoundé) et le \cle{Cameroun occidental} (ex-britannique, capitale Buéa). Chaque État fédéré possède sa constitution, son gouvernement et son assemblée propres, mais le pouvoir \cle{fédéral} est fort, incarné par le Président Ahidjo.
\end{definition}

\begin{savaistu}[Une constitution négociée en 10 jours]
La Constitution de Foumban (août 1961) a été négociée en \textbf{10 jours seulement}. Elle instaure un fédéralisme « à la camerounaise » : un Président fédéral très puissant, disposant même du droit de \cle{dissoudre} les assemblées des États fédérés. Ce déséquilibre prépare la disparition du fédéralisme en 1972.
\end{savaistu}

\begin{popfigure}
\begin{center}
\begin{tikzpicture}[font=\scriptsize, node distance=0.35cm and 0.5cm,
etape/.style={draw=popblue, fill=popblueL, rounded corners=2pt, align=center, minimum width=2.6cm, minimum height=0.85cm, thick},
etape2/.style={draw=popgreen, fill=popgreenL, rounded corners=2pt, align=center, minimum width=2.6cm, minimum height=0.85cm, thick},
fleche/.style={-{Stealth[length=3mm]}, thick, popdark}]
\node[etape, align=center] (a){Loi-cadre Defferre\\23 juin 1956};
\node[etape, right=of a, align=center] (b){Élections\\23 déc. 1957};
\node[etape, right=of b, align=center] (c){Gouvernement\\Mbida};
\node[etape, below=0.5cm of c, align=center] (d){Gouvernement Ahidjo\\fév. 1958};
\node[etape, left=of d, align=center] (e){Indépendance\\1\textsuperscript{er} janv. 1960};
\node[etape2, left=of e, align=center] (f){Référendum UK\\11 fév. 1961};
\node[etape2, below=0.5cm of f, align=center] (g){Conférence\\de Foumban\\août 1961};
\node[etape2, right=of g, align=center] (h){République Fédérale\\1\textsuperscript{er} oct. 1961};
\draw[fleche] (a) -- (b);
\draw[fleche] (b) -- (c);
\draw[fleche] (c) -- (d);
\draw[fleche] (d) -- (e);
\draw[fleche] (e) -- (f);
\draw[fleche] (f) -- (g);
\draw[fleche] (g) -- (h);
\end{tikzpicture}
\end{center}
\legende{Figure 11 — Diagramme de flux : les étapes institutionnelles de la marche vers l'indépendance et la réunification (1956-1961).}
\end{popfigure}

\courssub{Travaux pratiques : croquis « Le Cameroun à l'indépendance (1961) »}

\begin{methode}[Réaliser un croquis cartographique guidé (type Bac)]
\begin{enumerate}
\item \textbf{Tracer le fond de carte} : contour simplifié du Cameroun fédéral, ligne de partage Nord/Sud du référendum, frontière avec le Nigéria.
\item \textbf{Placer les villes} : Yaoundé (capitale fédérale et du Cameroun oriental), Buéa (capitale du Cameroun occidental), Douala (port).
\item \textbf{Figurer les richesses} : symboles localisés — café (ouest et centre), coton (nord), pétrole (zone côtière).
\item \textbf{Soigner la légende} : titre, orientation (nord), échelle indicative, légende numérotée.
\end{enumerate}
\end{methode}

\begin{popfigure}
\begin{center}
\begin{tikzpicture}[scale=0.95, font=\small]
\fill[popcream] (0,0) -- (0,6) -- (2.4,6.4) -- (3.8,5) -- (3.3,2.8) -- (3.8,1) -- (2.6,-0.6) -- (0.6,-0.4) -- cycle;
\draw[popdark, thick] (0,0) -- (0,6) -- (2.4,6.4) -- (3.8,5) -- (3.3,2.8) -- (3.8,1) -- (2.6,-0.6) -- (0.6,-0.4) -- cycle;
% Ligne de partage référendum (ouest)
\draw[popink, thick, dashed] (0.1,3.4) -- (1.6,3.6);
\node[popink, font=\scriptsize, rotate=4] at (0.85,3.85) {partage Nord/Sud 1961};
% Frontière fédérale (ex-britannique)
\draw[popgreen, thick] (0.05,0.2) -- (0.05,5.9);
\node[popgreen!60!black, font=\scriptsize, rotate=90] at (-0.25,3.1) {Cameroun occidental};
\node[popblue, font=\scriptsize] at (2.1,3.1) {Cameroun oriental};
% Villes
\fill[popink] (1.6,1.2) circle (0.09); \node[right, font=\scriptsize] at (1.75,1.2) {Yaoundé (cap. fédérale)};
\fill[popink] (0.35,1.6) circle (0.09); \node[left, font=\scriptsize] at (0.25,1.6) {Buéa};
\fill[popblue] (0.9,0.5) circle (0.09); \node[right, font=\scriptsize] at (1.0,0.45) {Douala (port)};
% Richesses
\node[font=\scriptsize] at (0.7,2.6) {\textbf{Café}};
\draw[popgold, thick] (0.55,2.35) circle (0.12);
\node[font=\scriptsize] at (2.2,1.9) {\textbf{Café}};
\draw[popgold, thick] (2.05,1.65) circle (0.12);
\node[font=\scriptsize] at (2.3,5.2) {\textbf{Coton}};
\draw[popmuted, thick] (2.15,4.95) circle (0.12);
\node[font=\scriptsize] at (1.35,0.05) {\textbf{Pétrole}};
\fill[popdark] (1.2,-0.15) circle (0.1);
% Nord
\draw[-{Stealth}, popdark] (4.3,5.6) -- (4.3,6.4) node[above, font=\scriptsize]{N};
% Échelle
\draw[thick] (0.2,-0.9) -- (1.2,-0.9) node[midway, below, font=\scriptsize] {0 \quad 200 km};
\end{tikzpicture}
\end{center}
\legende{Figure 12 — Croquis guidé : le Cameroun à l'indépendance (1961). 1 : ligne de partage du référendum ; 2 : limite des deux États fédérés ; 3 : capitales ; 4 : principales richesses (café, coton, pétrole).}
\end{popfigure}

\begin{methode}[Commenter une carte de référendum (type Bac)]
\begin{enumerate}
\item \textbf{Présenter le document} : nature (carte de résultats électoraux), titre, date (11 février 1961), échelle, contexte (plébiscite organisé par l'ONU au Cameroun britannique).
\item \textbf{Analyser les résultats} : dégager la géographie du vote — le Nord, musulman et tourné vers l'économie nigériane, choisit le Nigéria ; le Sud, chrétien et animiste, lié aux populations du Cameroun francophone, choisit la réunification. Croiser religion, ethnies et économie.
\item \textbf{Dégager les conséquences politiques} : partage du Cameroun britannique, nouvelles frontières, naissance de l'État fédéral en octobre 1961.
\end{enumerate}
\end{methode}

\begin{exoresolu}[Analyser les résultats du référendum de 1961]
Énoncé : au Sud britannique, 145 996 électeurs votent pour la réunification et 61 092 pour le Nigéria. Calcule la part de chaque option et explique le résultat.
\tcblower
Solution : total des suffrages : $145\,996 + 61\,092 = 207\,088$. Part de la réunification :
\[ \frac{145\,996}{207\,088} \times 100 \approx 70,5\ \% \]
Part du rattachement au Nigéria : $100 - 70,5 = 29,5\ \%$. Le Sud choisit donc nettement la réunification : ses populations (chrétiennes et animistes des Grassfields) sont culturellement proches des groupes du Cameroun francophone (Bamiléké notamment), et le leader John Ngu Foncha a mené campagne pour l'option camerounaise. Le 1\textsuperscript{er} octobre 1961, ce territoire devient le Cameroun occidental au sein de la République Fédérale.
\end{exoresolu}

\begin{exercice}[Pour s'entraîner]
\begin{enumerate}
\item Construis une frise chronologique de 1955 à 1961 avec huit dates essentielles de cette section.
\item Rédige un paragraphe argumenté : « La loi-cadre Defferre est-elle une victoire du nationalisme camerounais ? »
\item Explique pourquoi le Nord et le Sud du Cameroun britannique ont fait des choix opposés en 1961.
\end{enumerate}
\end{exercice}

\begin{corrige}[Exercice — indications]
1. Frise attendue : mai-juillet 1955 (insurrection, interdiction UPC) ; 23 juin 1956 (loi-cadre) ; 23 décembre 1957 (élections, Mbida) ; février 1958 (Ahidjo) ; 28 septembre 1958 (référendum) ; 12 novembre 1958 (proclamation) ; 1\textsuperscript{er} janvier 1960 (indépendance) ; 11 février et 1\textsuperscript{er} octobre 1961 (plébiscite, réunification).
2. Oui partiellement : elle accorde suffrage universel et autonomie interne (revendications nationalistes), mais l'UPC, qui portait ces exigences, est interdite et exclue du processus ; la France garde la défense et la diplomatie. Nuance exigée.
3. Nord : islam, émirs pro-nigérians, commerce du bétail vers Lagos ; Sud : proximités ethniques et religieuses avec le Cameroun francophone, campagne de Foncha.
\end{corrige}

\begin{aretenir}[L'essentiel de la section]
De 1955 à 1961, le Cameroun passe de la répression du nationalisme radical (UPC interdite, guerre cachée) à une décolonisation négociée : loi-cadre (1956), gouvernements autonomes (Mbida puis Ahidjo), indépendance du Cameroun français (1\textsuperscript{er} janvier 1960), puis réunification partielle après le référendum de 1961 et la Constitution de Foumban. La République Fédérale du Cameroun naît le 1\textsuperscript{er} octobre 1961 — mais sans le Nord britannique, devenu nigérian.
\end{aretenir}

\coursec{Construire l'État-nation : Défis de l'unité, de la réconciliation et de la mémoire (1961-1972 / Aujourd'hui)}

\courssub{Héritages divergents : le défi de l'intégration de deux États en un}

L'accession à la souveraineté internationale (1er janvier 1960 pour le Cameroun oriental, 1er octobre 1961 pour le Cameroun occidental) ne règle pas la question de la construction nationale. Elle inaugure au contraire la période la plus délicate : celle de la fusion de deux entités politiques, juridiques, administratives et culturelles façonnées par un demi-siècle de colonisations distinctes. La conférence de Foumban (juillet 1961) acte la forme fédérale, mais le contenu de l'unité reste à inventer au quotidien.

\begin{definition}[Héritages coloniaux divergents : le double legs]
Le Cameroun indépendant hérite de deux systèmes étatiques antinomiques :
\begin{itemize}
    \item \textbf{Droit et Justice :} Au Cameroun oriental (ex-français), le \textbf{Code civil} (droit écrit, romano-germanique) et le droit coutumier codifié coexistent sous le contrôle d'une magistrature unique. Au Cameroun occidental (ex-britannique), la \textbf{Common Law} (droit de la jurisprudence, \emph{stare decisis}) s'applique, avec des \emph{Native Courts} gérant le droit coutumier local.
    \item \textbf{Administration :} Modèle \textbf{centralisé} (préfets, sous-préfets nommés par Yaoundé) à l'Est ; modèle \textbf{décentralisé} (Native Authorities, conseils de district élus, pouvoir traditionnel fort) à l'Ouest.
    \item \textbf{Langue et Éducation :} Système francophone (programmes français, français langue d'enseignement) vs système anglophone (programmes britanniques, anglais langue d'enseignement, examen \emph{GCE O/A Level}).
    \item \textbf{Police/Sécurité :} Gendarmerie nationale (militaire, française) vs Police fédérale (civile, britannique).
\end{itemize}
\end{definition}

\begin{attention}[Piège de l'anachronisme : juger 1961 avec les yeux de 2024]
Il est tentant de qualifier l'unification de « néocoloniale » ou d'« assimilation forcée » à la lumière des standards actuels (droits des minorités, justice transitionnelle, fédéralisme asymétrique). Or, en 1961, l'impératif dominant chez les élites des deux bords (Ahidjo, Foncha, Muna) est la \textbf{survie de l'État} face aux risques de balkanisation (sécession du Nord, rébellion UPC). L'histoire explique les choix contraints par le contexte (Guerre froide, fragilité économique, absence d'administration formée), elle ne les juge pas a posteriori.
\end{attention}

\courssub{L'unification administrative et juridique (1961-1972) : une centralisation progressive}

La Fédération (1961-1972) est une architecture institutionnelle complexe : un président fédéral, deux premiers ministres (Est/Ouest), deux assemblées législatives, une assemblée fédérale. Mais très vite, la logique unitaire l'emporte sur l'esprit fédéral.

\begin{methode}[Analyser l'érosion du fédéralisme (1961-1972)]
Pour comprendre le passage de la Fédération à l'État Unitaire, suivez ces étapes :
\begin{enumerate}
    \item \textbf{Identifier les leviers constitutionnels :} Article 13 de la Constitution fédérale (pouvoir d'initiative législative du Président fédéral), Article 20 (pouvoir réglementaire), contrôle du budget fédéral.
    \item \textbf{Repérer les actes majeurs d'unification :}
    \begin{itemize}
        \item 1962 : Création de la \textbf{Gendarmerie nationale} (remplace la Police fédérale de l'Ouest, symbole d'uniformisation sécuritaire).
        \item 1963 : \textbf{Code de la nationalité} unique.
        \item 1965 : Création du \textbf{Parti unique} (UNC = KNDP + UC), fusion des partis politiques est/ouest.
        \item 1966 : \textbf{Loi sur l'organisation judiciaire} : harmonisation progressive, mais maintien de la dualité juridictionnelle (Cour d'appel de l'Ouest appliquant la Common Law).
        \item 1968 : Suppression du poste de Vice-Président fédéral (occupé par Foncha), remplacé par un Premier ministre fédéral.
    \end{itemize}
    \item \textbf{Évaluer le coût politique :} Démission de Foncha (1968), puis de Muna (1968), sentiment d'aliénation de l'élite anglophone.
    \item \textbf{Conclure sur le référendum du 20 mai 1972 :} Consultation organisée sur la question : « Voulez-vous que le Cameroun soit un État unitaire ? ». Résultat officiel : 99,99\% « Oui ». Fin de la Fédération, naissance de la \textbf{République Unie du Cameroun}.
\end{enumerate}
\end{methode}

\begin{popfigure}
\begin{tikzpicture}[scale=0.85, every node/.style={font=\small}]
    % Styles
    \tikzset{
        header/.style={fill=popblue, text=white, font=\bfseries\small, minimum height=0.7cm, align=center},
        colA/.style={fill=popblueL, minimum height=0.6cm, align=center},
        colB/.style={fill=poporangeL, minimum height=0.6cm, align=center},
        colC/.style={fill=popgreenL, minimum height=0.6cm, align=center},
        row/.style={minimum height=0.55cm, align=left},
        grid/.style={draw=popline, thin}
    }
    % Tableau : 4 colonnes (Critère, Fédéralisme 1961, Unitarisme 1972, Décentralisation 1996/2019)
    \matrix (tab) [matrix of nodes, nodes in empty cells, column sep=-\pgflinewidth, row sep=-\pgflinewidth,
        column 1/.style={nodes={row, text width=2.8cm, fill=popcream, font=\bfseries\small}},
        column 2/.style={nodes={row, text width=3.5cm}},
        column 3/.style={nodes={row, text width=3.5cm}},
        column 4/.style={nodes={row, text width=3.5cm}},
        row 1/.style={nodes={header, text width=3.5cm}},
        row 2/.style={nodes={colA}},
        row 3/.style={nodes={colB}},
        row 4/.style={nodes={colC}},
    ]{
        \textbf{Critère} & \textbf{Fédéralisme (1961-1972)} & \textbf{Unitarisme (1972-1996)} & \textbf{Décentralisation (1996 / 2019)} \\
        \textbf{Niveaux de pouvoir} & 2 États fédérés (Est/Ouest) + Fédéral. Assemblées distinctes. & 1 État central. Préfectures/Départements (administration déconcentrée). & État central + \textbf{Collectivités Territoriales Décentralisées (CTD)} : Régions, Communes. Conseil régional élu. \\
        \textbf{Langues officielles} & Anglais (Ouest) / Français (Est) au niveau fédéré. Français dominant au fédéral. & Français langue unique de l'administration centrale. Anglais langue officielle constitutionnelle (Art. 1er) mais usage marginalisé. & Bilinguisme officiel réaffirmé. \textbf{Commission Bilingue (CNBBM)} créée (2017). Statut spécial N-O/S-O (2019). \\
        \textbf{Justice} & Dualité : Cour Fédérale + Cours d'Appel Est (Droit civil) / Ouest (Common Law). & Unification formelle (Cour Suprême 1972). Maintien de fait de la Common Law dans l'Ouest (Section judiciaire). & Cour Suprême + Cour Constitutionnelle (2018). \textbf{Harmonisation des codes} (OHADA, Code pénal 2016). Dualité persistante (Common Law vs Droit civil). \\
        \textbf{Police / Sécurité} & Police Fédérale (Ouest) + Gendarmerie (Est). Fusion rapide en Gendarmerie Nationale (1962). & Gendarmerie Nationale unique (militaire). Police nationale (civile) créée plus tard. & Maintien Gendarmerie/Police. \textbf{Police de proximité} voulue par décentralisation (municipale). \\
        \textbf{Budget / Fiscalité} & Budgets fédérés autonomes + Budget fédéral. Ressources pétrolières (offshore Ouest) gérées fédéralement. & Budget unique centralisé. Recettes pétrolières (majoritairement zone anglophone) gérées par Yaoundé. & \textbf{Autonomie financière des CTD} (Loi 2009/2019). Partage recettes (centres d'impôts). Transferts de l'État (Dotation générale). \\
    };
    % Grille
    \draw[popline] (tab.north west) rectangle (tab.south east);
    \foreach \i in {1,...,5} {
        \draw[popline] (tab-\i-1.north west) -- (tab-\i-4.north east);
    }
    \foreach \j in {1,...,4} {
        \draw[popline] (tab-1-\j.north west) -- (tab-5-\j.south west);
    }
\end{tikzpicture}
\legende{Figure 13 : Schéma comparatif des trois architectures institutionnelles successives du Cameroun. Source : Constitutions 1961, 1972, 1996 (rév. 2008), Lois sur la décentralisation (2004, 2009, 2019).}
\end{popfigure}

\courssub{Gestion du passé : silences, amnisties et justice inachevée (1960-1990)}

La construction de l'État-nation s'est accompagnée d'une gestion autoritaire de la mémoire et de la dissidence. L'objectif affiché était l'unité nationale ; le résultat fut souvent l'impunité et le refoulement.

\begin{exemplebox}[La répression de l'UPC et la guerre du maquis (1955-1971)]
\begin{itemize}
    \item \textbf{1955-1960 :} Insurrection nationaliste (UPC) réprimée par l'administration française puis par l'armée camerounaise (création du BIR précurseur). Estimations : 10 000 à 100 000 morts (historiographie divergente).
    \item \textbf{1960 (Indépendance) :} Amnistie partielle. Les leaders historiques (Um Nyobé tué en 1958, Félix-Roland Moumié assassiné à Genève en 1960) sont absents. Ernest Ouandié et Abel Kingué poursuivent la lutte armée.
    \item \textbf{1964 :} Nouvelle amnistie. Ouandié se rend, puis repart au maquis.
    \item \textbf{1970-1971 :} Capture d'Ouandié (août 1970). Procès expéditif (Tribunal militaire de Yaoundé). \textbf{Exécution publique par peloton d'exécution (15 janvier 1971)} à Bafoussam. Symbole fort : fin de la rébellion armée, mais aussi fin de toute opposition armée légitime.
    \item \textbf{Silence officiel :} Jusqu'aux années 1990, l'histoire officielle occulte la guerre du maquis. Les manuels scolaires présentent l'indépendance comme « pacifique » (négociée).
\end{itemize}
\end{exemplebox}

\begin{aretenir}[Les amnisties : outil de pacification ou d'oubli imposé ?]
\begin{itemize}
    \item \textbf{1960 / 1964 / 1990 :} Mesures ciblées (droits civiques, libération prisonniers politiques) mais excluant souvent les « crimes de sang » ou les leaders historiques.
    \item \textbf{1991 (Conférence Tripartite / Comité Constitutionnel) :} Ouverture timide. Réhabilitation symbolique d'Um Nyobé (stèle à Eseka en 2007, rue à Yaoundé).
    \item \textbf{Absence de Commission Vérité et Réconciliation (CVR) :} Contrairement à l'Afrique du Sud (1995) ou au Ghana (2002), le Cameroun n'a jamais institué de mécanisme officiel d'enquête sur les violences de la décolonisation (1955-1971) ni sur les disparitions des années 1980-90 (Opération « Épervier »).
\end{itemize}
\end{aretenir}

\courssub{Enjeux mémoriels actuels : réhabilitation, restitution, place des maquisards}

Depuis les années 2000, une dynamique mémorielle portée par la société civile, les historiens et une partie de la classe politique oblige l'État à revisiter son récit national.

\begin{savaistu}[La réhabilitation d'Um Nyobé : un processus lent et symbolique]
\begin{itemize}
    \item \textbf{2007 :} Inauguration de la \textbf{stèle à Eseka} (lieu de sa mort) par le Premier Ministre Ephraïm Inoni. Première reconnaissance officielle du « Héros national ».
    \item \textbf{2010 (50e anniversaire) :} Discours du Président Biya : « Um Nyobé est un grand nationaliste ». Timide, sans qualification de « héros de l'indépendance ».
    \item \textbf{2020-2024 :} Demandes récurrentes de \textbf{restitution des archives} (Archives nationales d'Outre-mer à Aix-en-Provence, archives militaires de Vincennes) et des \textbf{restes mortels} (crâne d'Um Nyobé ? corps de Moumié ?). La France a restitué des archives numérisées (2021), mais les originaux et les biens culturels (trône Bangwa, etc.) restent en France.
    \item \textbf{Monument de la Réunification (Yaoundé, 1974, sculpture Gédéon Mpando) :} Symbolise l'union Est/Ouest (deux spirales qui fusionnent). Lieu de mémoire officiel, mais contesté par certains anglophones (symbole de l'absorption) et par les partisans de l'UPC (absence des maquisards).
\end{itemize}
\end{savaistu}

\begin{popfigure}
\begin{tikzpicture}[scale=1.1, every node/.style={font=\small}]
    % Représentation schématique du Monument de la Réunification (Yaoundé) et Stèle Eseka
    % Monument : deux spirales qui montent et se rejoignent
    \draw[popblue, very thick, domain=0:6.28, samples=200, smooth, variable=\t]
        plot ({1.5*cos(\t r) + 0.15*\t}, {1.5*sin(\t r) + 0.15*\t}, {0.3*\t});
    \draw[poporange, very thick, domain=0:6.28, samples=200, smooth, variable=\t]
        plot ({-1.5*cos(\t r) - 0.15*\t}, {1.5*sin(\t r) + 0.15*\t}, {0.3*\t});
    % Base
    \draw[popdark, thick] (-2.5,-0.5) -- (2.5,-0.5) -- (2.5,0) -- (-2.5,0) -- cycle;
    \node[popdark, align=center] at (0, -1.2) {\textbf{Monument de la Réunification (Yaoundé)}\\\emph{Spirales Est/Ouest fusionnant vers le haut}};
    % Stèle Eseka (décalée à droite)
    \begin{scope}[xshift=7cm]
        \draw[popgreen!70!black, thick] (-0.5,0) -- (-0.5,3) -- (0.5,3) -- (0.5,0) -- cycle;
        \draw[popgreen!70!black, thick] (-0.7,3) -- (0.7,3) -- (0,3.8) -- cycle;
        \node[popgreen!70!black, align=center, font=\tiny] at (0, 1.5) {UM NYOBÉ\\1913 - 1958};
        \node[popdark, align=center] at (0, -0.8) {\textbf{Stèle d'Eseka (2007)}\\\emph{Reconnaissance officielle tardive}};
    \end{scope}
    % Flèches symboliques
    \draw[->, popmuted, thick] (3.5, 1.5) -- (4.5, 1.5) node[midway, above, font=\tiny] {Mémoire partagée ?};
\end{tikzpicture}
\legende{Figure 15 : Analyse symbolique des lieux de mémoire officiels. À gauche, le Monument de la Réunification (Yaoundé, 1974, sculpteur Gédéon Mpando) : esthétique moderniste, fusion de deux entités. À droite, la Stèle d'Um Nyobé à Eseka (2007) : sobriété, ancrage territorial, réhabilitation individuelle. Le décalage temporel (30 ans) illustre la difficulté d'intégrer la lutte armée au récit national.}
\end{popfigure}

\courssub{Le bilinguisme officiel : ciment constitutionnel vs réalités asymétriques}

L'article 1er de la Constitution (1996, révisée 2008) proclame : « Les langues officielles de la République du Cameroun sont le français et l'anglais. Elles ont le même statut. L'État assure la promotion du bilinguisme sur tout le territoire. » La réalité administrative, judiciaire et éducative contredit souvent cet égalitarisme formel.

\begin{definition}[Bilinguisme officiel : principe et asymétrie]
\begin{itemize}
    \item \textbf{Principe constitutionnel (Art. 1er) :} Français et Anglais sont langues officielles à égalité. Tout citoyen a droit à s'adresser à l'administration dans l'une ou l'autre langue.
    \item \textbf{Réalité structurelle (Asymétrie forte) :}
    \begin{itemize}
        \item \textbf{Administration centrale :} Français langue de travail quasi exclusive (notes de service, correspondance, formation ENAM).
        \item \textbf{Justice :} Arrêts de la Cour Suprême/Cour Constitutionnelle rédigés en français. Procédures en Common Law (Nord-Ouest/Sud-Ouest) souvent traduites \emph{a posteriori}. Pénurie de magistrats bilingues compétents en Common Law.
        \item \textbf{Éducation :} Deux sous-systèmes parallèles (Francophone / Anglophone) qui s'ignorent. Peu d'écoles véritablement bilingues (immersion). Examens distincts (Baccalauréat vs GCE).
    \end{itemize}
\end{itemize}
\end{definition}

\begin{exemplebox}[Indicateurs de la fracture bilinguiste]
\begin{itemize}
    \item \textbf{Budget CNBBM (Commission Nationale pour le Bilinguisme et le Multiculturalisme, créée 2017) :} \textbf{$\sim$ 2 milliards FCFA/an} (budget de fonctionnement). Soit $\sim$ 0,01\% du budget national. Montre la priorité politique réelle.
    \item \textbf{Taux de réussite au Baccalauréat :} Écart structurel de \textbf{15 à 20 points} en faveur du sous-système francophone (ex: Session 2023 : Francophone $\sim$ 65\% vs Anglophone $\sim$ 45\%).
    \item \textbf{Fonction publique :} $\sim$ 85\% des cadres supérieurs sont francophones (concours en français, formation francophone).
    \item \textbf{Décrets d'application :} Loi d'orientation n°98/004 (éducation) prévoit le bilinguisme. Décrets d'application (organisation des examens, statut enseignants) tardifs ou incomplets.
\end{itemize}
\end{exemplebox}

\begin{methode}[Argumenter sur une question d'actualité : « Le fédéralisme est-il la solution à la crise anglophone ? »]
Sujet type Bac / Concours. Suivez cette structure rigoureuse :
\begin{enumerate}
    \item \textbf{Contexte historique (Introduction) :} Rappeler Foumban (1961), le pacte fédéral (2 États, 2 systèmes), l'unilatéralisme de 1972 (Référendum), la révision de 1996 (Décentralisation, pas fédéralisme), la crise 2016 (revendications corporatistes $\to$ politiques).
    \item \textbf{Arguments POUR le retour au fédéralisme (ou fédéralisme asymétrique) :}
    \begin{itemize}
        \item Respect de l'identité juridique (Common Law) et éducative (GCE) anglophone.
        \item Proximité démocratique : gestion locale des ressources (pétrole, bois, agriculture), police de proximité.
        \item Réponse au sentiment d'annexion (« marginalisation ») exprimé depuis 2016.
        \item Précédent historique : la Fédération a fonctionné (1961-1972) avec croissance économique.
    \end{itemize}
    \item \textbf{Arguments CONTRE (Risques et limites) :}
    \begin{itemize}
        \item Risque de balkanisation / sécession (précédent du Sud-Soudan, menace Biafra).
        \item Coût institutionnel élevé (doublons : 2 gouvernements, 2 parlements, 2 administrations) pour un pays à revenu intermédiaire.
        \item Précédent 1972 : le fédéralisme n'a pas empêché la centralisation (parti unique, pouvoir personnel).
        \item La décentralisation (Loi 2019, Statut Spécial N-O/S-O) offre un cadre juridique actuel pour l'autonomie sans rupture étatique.
    \end{itemize}
    \item \textbf{Synthèse personnelle nuancée (Conclusion) :} Le fédéralisme n'est pas une baguette magique. La solution passe par une \textbf{décentralisation effective} (transferts de compétences \textbf{et} de ressources financières réels), la \textbf{reconnaissance institutionnelle de la Common Law} (Cour d'appel dédiée, formation barreau), et une \textbf{justice transitionnelle} sur les violences 2016-2024. L'unité se construit dans la reconnaissance de la diversité, non dans son déni.
\end{enumerate}
\end{methode}

\courssub{La crise anglophone (2016-...) : révélateur des failles de l'État-nation}

Ce qui commence en octobre 2016 par une grève corporatiste (avocats, enseignants du Nord-Ouest/Sud-Ouest) contre l'érosion du système juridique et éducatif anglophone bascule en conflit armé (fin 2017). Elle met à nu l'échec de l'intégration nationale.

\begin{popfigure}
\begin{tikzpicture}[scale=0.9, every node/.style={font=\small}]
    % Carte schématique Cameroun - Focus Nord-Ouest / Sud-Ouest
    % Contour très simplifié
    \draw[popline, thick] (0,0) -- (10,0) -- (10,8) -- (0,8) -- cycle; % Cadre
    % Contour pays (polygone grossier)
    \draw[popdark, thick, fill=popcream] 
        (1,1) -- (9,1) -- (9.5,2) -- (9,3.5) -- (8.5,5) -- (8,6.5) -- (6,7.5) -- (3,7.5) -- (1.5,6) -- (0.5,4) -- (1,2.5) -- cycle;
    % Régions (simplifiées)
    % Nord-Ouest (zone crise)
    \fill[popred!30, opacity=0.6] (2.5,4.5) -- (4.5,5.5) -- (4,6.5) -- (2,6) -- cycle;
    \node[popred, font=\bfseries] at (3.2, 5.5) {NORD-OUEST};
    % Sud-Ouest (zone crise)
    \fill[popred!30, opacity=0.6] (4.5,3.5) -- (6.5,3) -- (7,4.5) -- (5.5,5.5) -- (4.5,5.5) -- cycle;
    \node[popred, font=\bfseries] at (5.8, 4.2) {SUD-OUEST};
    % Villes clés
    \node[circle, fill=popdark, inner sep=1.5pt] (bamenda) at (3, 5.2) {};
    \node[anchor=west, font=\tiny] at (3.1, 5.2) {Bamenda};
    \node[circle, fill=popdark, inner sep=1.5pt] (buae) at (5.5, 3.8) {};
    \node[anchor=west, font=\tiny] at (5.6, 3.8) {Buea};
    \node[circle, fill=popblue, inner sep=1.5pt] (yaounde) at (5.5, 2.2) {};
    \node[anchor=north, font=\tiny, popblue] at (5.5, 2.0) {Yaoundé};
    % Flèches déplacés internes (IDP)
    \draw[->, poporange, thick, decorate, decoration={snake, amplitude=1pt, segment length=4pt}] (3, 5.2) -- (2, 3.5) node[left, font=\tiny, poporange] {IDP $\to$ Littoral/Centre};
    \draw[->, poporange, thick, decorate, decoration={snake, amplitude=1pt, segment length=4pt}] (5.5, 3.8) -- (6.5, 2.5) node[right, font=\tiny, poporange] {IDP $\to$ Douala};
    % Écoles fermées (symboles)
    \foreach \x/\y in {2.8/5.5, 3.5/5.8, 5.2/4.5, 6.0/4.0} {
        \draw[popred, thick] (\x-0.1,\y-0.1) -- (\x+0.1,\y+0.1);
        \draw[popred, thick] (\x-0.1,\y+0.1) -- (\x+0.1,\y-0.1);
    }
    \node[popred, font=\tiny, align=center] at (7.5, 6.5) {$\times$ = Écoles fermées\\(années 2017-2023)};
    % Légende
    \begin{scope}[yshift=-2.5cm]
        \node[anchor=west] at (0,0) {\textbf{Légende :}};
        \fill[popred!30] (0,-0.4) rectangle (0.5,-0.1); \node[anchor=west, font=\tiny] at (0.6,-0.25) {Zones touchées (N-O / S-O)};
        \draw[->, poporange, thick, decorate, decoration={snake, amplitude=1pt, segment length=4pt}] (0,-0.65) -- (0.5,-0.65); \node[anchor=west, font=\tiny] at (0.6,-0.65) {Déplacés internes (IDP > 700 000 OCHA)};
        \draw[popred, thick] (0,-0.9) -- (0.2,-0.65); \draw[popred, thick] (0,-0.65) -- (0.2,-0.9); \node[anchor=west, font=\tiny] at (0.6,-0.77) {Écoles fermées / « Ghost Towns »};
        \node[circle, fill=popblue, inner sep=1.5pt] at (0,-1.15) {}; \node[anchor=west, font=\tiny] at (0.6,-1.15) {Capitale (Yaoundé)};
    \end{scope}
    % Flèche Nord
    \draw[->, popdark] (8.5, 7.2) -- (8.5, 7.8) node[above, font=\tiny] {N};
\end{tikzpicture}
\legende{Figure 14 : Carte schématique de la crise anglophone (2016-...). Les régions du Nord-Ouest et du Sud-Ouest (zones rouges) concentrent les combats, les déplacements de populations (flèches orange vers Douala/Yaoundé/Littoral) et la fermeture massive d'écoles (croix rouges). Source : OCHA, UNHCR, Rapports ONG (2023-2024).}
\end{popfigure}

\begin{exoresolu}[Analyse de la réponse politique : Le Grand Dialogue National (GDN, 30 sept - 4 oct 2019)]
\textbf{Énoncé :} En quoi le Grand Dialogue National (GDN) de 2019 constitue-t-il une tentative de réponse politique à la crise anglophone ? Quelles en sont les limites structurelles ?

\textbf{Solution détaillée :}
\begin{enumerate}
    \item \textbf{Contexte :} Crise ouverte depuis oct. 2016. Radikalisation (sécessionnistes « Ambazonie »), réponse militaire (Opération « Bamenda Clean », « Bui Clean »), crise humanitaire (IDP, réfugiés au Nigeria). Pression internationale (UA, UE, ONU, France, USA).
    \item \textbf{Mesures phares du GDN (Décret présidentiel nov. 2019) :}
    \begin{itemize}
        \item \textbf{Statut Spécial pour le N-O et S-O :} Assemblées régionales élargies, Conseil des Chefs traditionnels, Police régionale, gestion locale de l'éducation (sous-système anglophone préservé).
        \item \textbf{Commission Nationale pour le Bilinguisme et le Multiculturalisme (CNBBM) :} Renforcée, saisine facilitée.
        \item \textbf{Décentralisation accélérée :} Élections régionales (déc. 2020), transferts de compétences (santé, éducation de base, routes rurales) et de ressources (centres d'impôts régionaux).
        \item \textbf{Justice :} Création de la Chambre des Comptes (Cour Suprême), Section Common Law à la Cour Suprême (décret 2020), recrutement magistrats anglophones.
        \item \textbf{DDR (Désarmement, Démobilisation, Réintégration) :} Centres créés (Bamenda, Buea, Douala). Résultats mitigés (peu de renditions volontaires, méfiance).
    \end{itemize}
    \item \textbf{Limites structurelles (Analyse critique) :}
    \begin{itemize}
        \item \textbf{Absence des acteurs armés :} Sécessionnistes (Ambazonia Governing Council, ADF) et diaspora radicale exclus/boycottent. Dialogue « entre soi » (partis politiques légaux, société civile pro-gouvernement, chefs traditionnels).
        \item \textbf{Statut Spécial vs Fédéralisme :} Le Statut Spécial (Art. 62 Const. 1996) maintient la tutelle de l'État central (Gouverneur = représentant du Président, tutelle administrative). Pas d'autonomie législative propre (lois votées à Yaoundé).
        \item \textbf{Continuité répressive :} Parallèlement au GDN, poursuite des opérations militaires, arrestations arbitraires, lois antiterroristes (2014) utilisées contre opposants politiques/journalistes.
        \item \textbf{Application lente :} Élections régionales (2020) boycottées par l'opposition majeure (MRC). Transferts financiers effectifs aux Régions encore partiels (2024). CNBBM sous-dotée (2 Mrds FCFA).
    \end{itemize}
    \item \textbf{Conclusion :} Le GDN est une \textbf{réponse institutionnelle par le haut} (octroi), non une \textbf{négociation politique par le bas} (pacte). Il apaise la forme (cadre légal décentralisation) mais ne règle pas le fond (reconnaissance de l'identité politique anglophone, justice pour les victimes, partage du pouvoir). La crise persiste (2024) sous forme de guérilla à basse intensité et de crise de gouvernance.
\end{enumerate}
\end{exoresolu}

\courssub{Bilan : un processus inachevé}

\begin{propriete}[Théorème : L'indépendance juridique a précédé l'intégration nationale effective]
La construction de l'État-nation camerounais est un \textbf{processus inachevé}. L'accession à la souveraineté internationale (1960/1961) a créé un sujet de droit international, mais n'a pas produit une communauté politique unifiée (\emph{nation} au sens politique). La crise anglophone (2016-) en est la manifestation la plus aiguë : elle révèle que l'unité administrative (centralisation 1972, décentralisation 1996/2019) n'a pas résolu la question de la \textbf{reconnaissance mutuelle} des identités historiques (Common Law / Civil Law), de la \textbf{justice transitionnelle} (guerre du maquis, violences post-électorales 1992, 2008, 2018, crise N-O/S-O) et du \textbf{partage équitable des ressources et du pouvoir}. L'histoire du Cameroun contemporain est celle de la tension permanente entre un État fort, centralisateur, héritier de l'administration coloniale, et des sociétés locales plurielles qui négocient, résistent ou s'arment pour exister.
\end{propriete}

\begin{aretenir}[Clés de lecture pour le Bac - Section 5]
\begin{itemize}
    \item \textbf{1961-1972 :} Fédéralisme « à géométrie variable » $\to$ Unitarisme (Référendum 20 mai 1972). Centralisation du pouvoir (Ahidjo), unification partis (UNC), gendarmerie, nationalité.
    \item \textbf{Gestion mémoire :} Silences sur guerre UPC (1955-71), exécution Ouandié (1971), amnisties sélectives. Réhabilitation symbolique tardive (Um Nyobé 2007). Pas de CVR.
    \item \textbf{Bilinguisme :} Principe constitutionnel (Art. 1er) vs Asymétrie réelle (administration, justice, éducation, budget CNBBM dérisoire). Écart réussite Bac (15-20\%).
    \item \textbf{Crise anglophone (2016-) :} Déclencheur corporatiste $\to$ Conflit armé / Crise politique. Racines : non-respect engagements Foumban (Common Law, Éducation), marginalisation ressentie.
    \item \textbf{GDN 2019 :} Réponse par le haut (Statut Spécial, Décentralisation, DDR). Limites : exclusion acteurs armés, tutelle maintenue, application lente, répression parallèle.
    \item \textbf{Enjeu actuel :} Sortir de l'État-nation « inachevé » par justice transitionnelle, décentralisation effective (ressources + compétences), reconnaissance institutionnelle de la dualité juridique/éducative.
\end{itemize}
\end{aretenir}

\newpage

\coursec{Activité d'intégration}

\coursec{La marche vers l’indépendance : Activité d’intégration — Simulation ONU (1959)}

\courssub{Objectifs de l'activité}

Cette activité d'intégration clôt la leçon sur la marche du Cameroun vers l'indépendance. Elle place l'élève dans une situation authentique de simulation diplomatique : rédiger, en 1959, le rapport d'une mission fictive de l'Organisation des Nations Unies (ONU) chargée de statuer sur l'avenir du Cameroun britannique (\cle{Southern Cameroons} et \cle{Northern Cameroons}).

À l'issue de l'activité, l'élève doit être capable de :
\begin{itemize}
\item mobiliser les savoirs de la leçon : la tutelle de l'ONU, le nationalisme camerounais après 1945, les grandes figures du nationalisme, les étapes de 1955 à 1961 ;
\item analyser un dossier documentaire hétérogène (pétitions, rapports administratifs, statistiques, cartes, discours) ;
\item rédiger et justifier une conclusion argumentée, dans un style diplomatique adapté ;
\item appliquer les notions de l'unité (\cle{tutelle}, \cle{référendum}, \cle{réunification}, \cle{intégration}) à une situation concrète ;
\item travailler en groupe et restituer oralement un travail collectif.
\end{itemize}

\courssub{Situation}

\textbf{New York, siège de l'ONU, octobre 1959.} Vous êtes les quatre membres de la \cle{Mission des Nations Unies pour le Cameroun britannique}. Le Conseil de tutelle vous a donné le mandat suivant : enquêter sur place, puis recommander les modalités d'un plébiscite destiné à régler le statut définitif des deux territoires du Cameroun sous administration britannique, le \cle{Northern Cameroons} (administré avec le Nigeria) et le \cle{Southern Cameroons} (administré comme une province du Nigeria jusqu'en 1954, puis doté d'un statut autonome).

\textbf{Contexte rappelé par le Secrétaire général.} Depuis 1946, le Cameroun n'est plus un mandat de la Société des Nations (SDN) mais un territoire sous \cle{tutelle} de l'ONU, administré pour les quatre cinquièmes par la France (\cle{Cameroun oriental}) et pour un cinquième par le Royaume-Uni (\cle{Cameroun occidental}, deux bandes de territoire le long de la frontière nigériane). Le 1\textsuperscript{er} janvier 1960, le Cameroun français accédera à l'indépendance sous la présidence d'\cle{Ahmadou Ahidjo}, après avoir connu le soulèvement de l'\cle{UPC} (Union des populations du Cameroun), parti nationaliste créé en 1948, interdit par la France en 1955. La question se pose donc : que faire du Cameroun britannique ?

\begin{popfigure}
\begin{tikzpicture}[scale=0.85, every node/.style={font=\small}]
% Contour très simplifié du Cameroun (polygone approximatif)
\draw[popdark, thick, fill=popcream] (0,0) -- (6,0) -- (7,2) -- (6.5,5) -- (4,6) -- (1,5.5) -- (0.5,3) -- cycle;
% Frontière Est/Ouest (approximative)
\draw[popdark, dashed, thick] (2.5,0.5) -- (3,3.5) -- (3.5,5.5);
% Nigeria à l'Ouest
\draw[popmuted, fill=popblueL!30] (-2,0) rectangle (0,6);
\node[popdark] at (-1,5.5) {\textbf{Nigeria}};
% Cameroun Oriental (Français)
\draw[popgreenL!50, fill=popgreenL!20] (3.5,0.5) -- (7,2) -- (6.5,5) -- (4,6) -- (3.5,5.5) -- cycle;
\node[popgreen!80!black, align=center] at (5.5,3) {\textbf{Cameroun oriental}\\(Français)\\\textit{Indép. 1er janv. 1960}};
% Northern Cameroons
\draw[poporangeL!50, fill=poporangeL!20] (0,0) -- (2.5,0.5) -- (3,3.5) -- (1,3) -- cycle;
\node[poporange!80!black, align=center] at (1.2,1.5) {\textbf{Northern Cameroons}\\(Britannique)\\\textit{~700 000 hab.}};
% Southern Cameroons
\draw[poppurpleL!50, fill=poppurpleL!20] (1,3) -- (3,3.5) -- (3.5,5.5) -- (1,5.5) -- cycle;
\node[poppurple!80!black, align=center] at (2,4.5) {\textbf{Southern Cameroons}\\(Britannique)\\\textit{~800 000 hab.}};
% Villes repères
\node[circle,fill=popdark,inner sep=1.5pt] (yaounde) at (4.5,2.5) {};
\node[right] at (yaounde) {Yaoundé};
\node[circle,fill=popdark,inner sep=1.5pt] (buea) at (1.8,4) {};
\node[left] at (buea) {Buea};
\node[circle,fill=popdark,inner sep=1.5pt] (garoua) at (1.2,1) {};
\node[below] at (garoua) {Garoua};
% Flèche Nord
\begin{scope}[shift={(5.5,5.5)}, rotate=0]
\draw[popdark, thick, -{Stealth[length=3mm]}] (0,0) -- (0,0.8);
\node[above] at (0,0.8) {N};
\end{scope}
% Échelle indicative
\draw[popdark, thick] (0.2,0.2) -- (1.2,0.2);
\node[below] at (0.7,0.1) {~200 km};
\end{tikzpicture}
\legende{Carte schématique du Cameroun sous tutelle (1946-1960) : distinction entre le Cameroun oriental (français) et les deux composantes du Cameroun occidental (britannique). Les frontières sont approximatives.}
\end{popfigure}

\textbf{Votre dossier documentaire contient les pièces suivantes :}

\begin{enumerate}
\item \textbf{Document 1 — Extraits de pétitions (1948-1958).} Des pétitions de l'UPC réclament depuis 1948 « l'indépendance immédiate et la réunification des deux Cameroun ». Des pétitions du KNC (Kamerun National Congress) du Dr \cle{E. M. L. Endeley} demandent au contraire l'intégration du Southern Cameroons au Nigeria voisin. Le KNDP (Kamerun National Democratic Party) de \cle{John Ngu Foncha}, créé en 1955, réclame la réunification avec le Cameroun français devenu indépendant.
\item \textbf{Document 2 — Rapport administratif britannique (1958).} Le Royaume-Uni, puissance administrante, juge le territoire « économiquement non viable » et souhaite se retirer rapidement. Il souligne que le Northern Cameroons est intégré administrativement au Nigeria depuis des décennies.
\item \textbf{Document 3 — Données socio-économiques (1958).} Southern Cameroons : environ 800 000 habitants ; économie de plantation (bananes, palmiers à huile, hévéas) tournée vers le Nigeria ; revenu par habitant inférieur à celui du Nigeria ; réseau routier quasi inexistant vers l'est (Cameroun français). Northern Cameroons : environ 700 000 habitants, majoritairement musulmans, rattachés aux émirats du nord du Nigeria.
\item \textbf{Document 4 — Carte ethno-religieuse simplifiée.} Le Southern Cameroons est majoritairement chrétien et animiste, de langue anglaise ; le Northern Cameroons est majoritairement musulman, culturellement proche du nord du Nigeria.
\item \textbf{Document 5 — Extraits de discours (1959).} Ahmadou Ahidjo : « Le Cameroun britannique est une partie de notre territoire national amputée par la colonisation ; nous souhaitons la réunification dans la paix. » John Ngu Foncha : « Nous ne voulons ni rester une annexe du Nigeria, ni être absorbés : nous voulons choisir librement. » Endeley : « L'intégration au Nigeria garantit notre développement économique. »
\end{enumerate}

\textbf{Problème posé à la Mission :} faut-il organiser un plébiscite ? Si oui, avec quelle(s) question(s) ? Faut-il traiter le Northern et le Southern Cameroons ensemble ou séparément ? Quelles garanties offrir aux populations ? Votre rapport final sera présenté devant le Conseil de tutelle (la classe) et devra être historiquement rigoureux.

\courssub{Consignes}

Le travail se fait en groupes de quatre élèves, avec les rôles suivants : \textbf{Rapporteur} (rédaction et synthèse), \textbf{Expert juridique} (statut de tutelle, droit des peuples, mandat de l'ONU), \textbf{Expert socio-économique} (documents 3 et 4), \textbf{Expert politique} (documents 1 et 5 : partis et figures du nationalisme).

\begin{enumerate}
\item \textbf{Analyser la situation.} À partir des documents 2, 3 et 4, dégagez trois faits chiffrés qui éclairent le problème posé. Pourquoi la différence entre le Northern et le Southern Cameroons interdit-elle de traiter les deux territoires comme un seul bloc ?
\item \textbf{Identifier les acteurs.} Complétez un tableau présentant, pour chaque acteur (UPC, KNC/Endeley, KNDP/Foncha, Ahidjo, Royaume-Uni, France, Nigeria), sa position sur l'avenir du Cameroun britannique et l'argument principal qu'il avance.
\item \textbf{Formuler une recommandation motivée.} Choisissez l'une des options suivantes et justifiez-la en vous appuyant sur au moins trois arguments tirés des documents : (a) intégration au Nigeria ; (b) réunification avec la République du Cameroun ; (c) indépendance séparée ; (d) report de la décision. Précisez si votre recommandation vaut pour les deux territoires ou si elle les distingue.
\item \textbf{Proposer les modalités du scrutin.} Rédigez la ou les questions du plébiscite, la date envisagée, le corps électoral concerné et les conditions de contrôle international.
\item \textbf{Garantir les droits des minorités.} Proposez deux mesures concrètes de protection des minorités (linguistiques, religieuses, politiques) dans l'État qui accueillera le territoire.
\item \textbf{Restituer.} Présentez oralement votre rapport en cinq minutes devant le « Conseil de tutelle » (la classe), qui pourra vous interroger. Le rapport écrit (une page) est remis à l'enseignant.
\end{enumerate}

\begin{attention}[Pièges à éviter]
Ne projetez pas le résultat connu (le plébiscite du 11 février 1961) comme une évidence : en 1959, rien n'est joué. Ne confondez pas le Cameroun français (indépendant le 1\textsuperscript{er} janvier 1960) et le Cameroun britannique (qui choisit en 1961). Enfin, une recommandation sans argument documentaire précis (chiffre, citation, fait daté) ne sera pas recevable devant le Conseil de tutelle.
\end{attention}

\courssub{Ressources à mobiliser}

\begin{aretenir}[Notions et repères utiles]
\begin{itemize}
\item \textbf{Tutelle ONU (1946) :} régime international succédant au mandat SDN ; la puissance administrante (France, Royaume-Uni) doit conduire le territoire vers l'autonomie ou l'indépendance et rend compte au Conseil de tutelle ; les habitants peuvent adresser des pétitions à l'ONU.
\item \textbf{Nationalisme camerounais :} essor après 1945 (partis, syndicats, presse) ; l'UPC (1948, Ruben Um Nyobè) réclame indépendance et réunification ; répression après 1955.
\item \textbf{Figures :} Um Nyobè (UPC), Ahidjo (premier président du Cameroun indépendant), Foncha (KNDP, réunification), Endeley (KNC, intégration au Nigeria).
\item \textbf{Étapes clés :} 1946 (Tutelle ONU), 1948 (Création UPC), 1955 (Interdiction UPC), 1956 (Loi-cadre Defferre), 1958 (Autonomie interne Cameroun français), 1\textsuperscript{er} janvier 1960 (Indépendance Cameroun français), 11 février 1961 (Plébiscite Cameroun britannique), 1\textsuperscript{er} octobre 1961 (République fédérale du Cameroun).
\item \textbf{Savoir-faire :} analyser un document (nature, auteur, date, contexte, idée principale) ; rédiger une conclusion justifiée par des faits ; distinguer fait et opinion.
\end{itemize}
\end{aretenir}

\begin{popfigure}
\begin{tikzpicture}[scale=0.9, every node/.style={font=\small}]
% Frise chronologique 1946-1961
\draw[popdark, thick] (0,0) -- (15,0);
% Graduations
\foreach \x/\year in {0/1946, 1.5/1948, 3.5/1955, 4.5/1956, 6/1958, 7.5/1960, 8.5/1961} {
    \draw[popdark, thick] (\x,-0.15) -- (\x,0.15);
    \node[below, popdark] at (\x,-0.3) {\year};
}
% Événements au-dessus
\node[above, align=center, popblue] at (0,0.5) {Tutelle ONU};
\node[above, align=center, poppurple] at (1.5,1) {Création UPC\\(Um Nyobè)};
\node[above, align=center, poporange] at (3.5,0.5) {Interdiction UPC};
\node[above, align=center, popgreen] at (4.5,1) {Loi-cadre\\Defferre};
\node[above, align=center, popblue] at (6,0.5) {Autonomie\\interne};
\node[above, align=center, popgold] at (7.5,1) {Indépendance\\Cameroun français\\(Ahidjo)};
\node[above, align=center, poppink] at (8.5,0.5) {Plébiscite\\11 fév. 1961};
\node[above, align=center, popdark] at (10,1) {Rép. Fédérale\\1er oct. 1961};
% Flèches de liaison
\draw[popline, -{Stealth[length=2mm]}] (0,0.3) -- (1.5,0.8);
\draw[popline, -{Stealth[length=2mm]}] (1.5,0.8) -- (3.5,0.3);
\draw[popline, -{Stealth[length=2mm]}] (3.5,0.3) -- (4.5,0.8);
\draw[popline, -{Stealth[length=2mm]}] (4.5,0.8) -- (6,0.3);
\draw[popline, -{Stealth[length=2mm]}] (6,0.3) -- (7.5,0.8);
\draw[popline, -{Stealth[length=2mm]}] (7.5,0.8) -- (8.5,0.3);
\draw[popline, -{Stealth[length=2mm]}] (8.5,0.3) -- (10,0.8);
\end{tikzpicture}
\legende{Frise chronologique simplifiée : les grandes étapes de la décolonisation du Cameroun (1946-1961).}
\end{popfigure}

\courssub{Démarche et corrigé commenté}

\begin{exoresolu}[Exemple de rapport de la Mission ONU (corrigé type)]
Énoncé : rédiger le rapport final de la Mission (analyse, position des acteurs, recommandation motivée, garanties aux minorités) en mobilisant les cinq documents du dossier.
\tcblower
\textbf{Étape 1 — Analyse de la situation (faits et chiffres).} Trois faits ressortent du dossier. D'abord, le territoire est \emph{double et hétérogène} : le Southern Cameroons (environ 800 000 habitants, majoritairement chrétien et anglophone) n'a presque aucune route vers le Cameroun français, tandis que le Northern Cameroons (environ 700 000 habitants, majoritairement musulman) est intégré depuis des décennies à l'administration du nord du Nigeria (documents 2 et 4). Ensuite, le territoire est \emph{économiquement fragile} : le rapport britannique le juge « non viable » et son revenu par habitant est inférieur à celui du Nigeria (document 3). Enfin, la puissance administrante veut \emph{se retirer vite}, ce qui impose un calendrier court (document 2). Conclusion de l'étape : les deux territoires ne peuvent pas être traités comme un bloc unique ; une consultation séparée s'impose.

\textbf{Étape 2 — Position des acteurs.} L'UPC réclame depuis 1948 l'indépendance et la réunification (document 1) ; mais interdite en 1955, elle est affaiblie. Au Southern Cameroons, deux partis s'affrontent : le KNC d'Endeley pour l'intégration au Nigeria (argument économique), le KNDP de Foncha pour la réunification avec le Cameroun indépendant, tout en exigeant un choix libre (documents 1 et 5). Ahidjo, président du Cameroun indépendant, plaide pour la réunification « dans la paix » (document 5). Le Royaume-Uni veut partir ; la France et le Nigeria observent sans formellement s'opposer à une consultation.

\begin{center}
\begin{tabularx}{\linewidth}{|l|X|X|}\hline
\rowcolor{popblueL}
\textbf{Acteur} & \textbf{Position} & \textbf{Argument principal} \\ \hline
UPC (Um Nyobè) & Indépendance immédiate \& Réunification & Droit des peuples, unité nationale historique \\ \hline
KNC (Endeley) & Intégration au Nigeria & Viabilité économique, liens administratifs existants \\ \hline
KNDP (Foncha) & Réunification avec le Cameroun français & Choix libre, refus de l'annexion, identité commune \\ \hline
Ahmadou Ahidjo & Réunification & Intégrité territoriale, paix \\ \hline
Royaume-Uni & Retrait rapide, organisation du scrutin & Non-viabilité économique, fin de la tutelle \\ \hline
France & Soutien à l'indépendance du Cameroun français & Accords de coopération, stabilité régionale \\ \hline
Nigeria & Intégration (surtout Northern) & Unité administrative, économique, religieuse \\ \hline
\end{tabularx}
\end{center}

\textbf{Étape 3 — Recommandation motivée.} La Mission recommande la tenue d'un \cle{plébiscite sous contrôle de l'ONU}, organisé \emph{séparément} dans le Northern et le Southern Cameroons. Justification : (1) le principe de la tutelle impose de consulter les populations (droit des peuples à disposer d'eux-mêmes) ; (2) l'hétérogénéité des deux zones (document 4) rend impossible une question commune ; (3) aucune option ne fait l'unanimité chez les acteurs locaux (document 1), donc seul le scrutin peut trancher légitimement. L'option « indépendance séparée » est écartée en raison de la non-viabilité économique constatée (document 3) ; l'option « report » est écartée car elle prolongerait l'incertitude alors que le Cameroun français accède à l'indépendance le 1\textsuperscript{er} janvier 1960.

\textbf{Étape 4 — Modalités du scrutin.} Au Southern Cameroons, la question sera : « Souhaitez-vous l'indépendance par (a) l'intégration à la République fédérale du Nigeria, ou (b) la réunification avec la République du Cameroun ? ». Au Northern Cameroons, la même alternative est posée. Corps électoral : tous les adultes inscrits. Contrôle : observateurs de l'ONU, dépouillement public, campagne équitable garantie aux deux camps.

\textbf{Étape 5 — Garanties pour les minorités.} La Mission recommande : (1) la reconnaissance officielle de la langue anglaise et de la langue française dans l'État réunifié, pour protéger la minorité anglophone ; (2) la garantie de la liberté religieuse et du respect des chefferies traditionnelles, notamment dans le Northern Cameroons musulman ; (3) la représentation des populations concernées dans les institutions (assemblée régionale), afin qu'aucune communauté ne soit absorbée sans voix.

\textbf{Commentaire méthodologique.} Ce corrigé montre la démarche attendue : chaque recommandation s'appuie sur un document précis et daté ; la distinction Northern/Southern structure tout le raisonnement ; le style reste neutre et diplomatique (« la Mission recommande », « il apparaît que »). Historiquement, le plébiscite eut bien lieu le 11 février 1961 : le Northern Cameroons choisit le Nigeria (majorité pour l'intégration), le Southern Cameroons choisit la réunification (majorité pour la réunification), aboutissant à la République fédérale du Cameroun le 1\textsuperscript{er} octobre 1961 — ce qui valide \emph{a posteriori} la logique d'une consultation séparée.
\end{exoresolu}

\courssub{Bilan de l'activité}

Cette simulation a permis de mettre en évidence plusieurs acquis essentiels de la leçon.

\begin{itemize}
\item \textbf{La tutelle n'est pas la colonisation classique :} l'ONU contrôle, reçoit les pétitions et impose une consultation des populations — ce qui explique le rôle central du plébiscite de 1961 dans la décolonisation du Cameroun britannique.
\item \textbf{Le nationalisme camerounais est pluriel :} loin d'être un bloc unique, il se divise (UPC, KNC, KNDP) sur la question même de l'unité nationale ; les figures d'Um Nyobè, Ahidjo, Foncha et Endeley incarnent des options différentes.
\item \textbf{La réunification n'était pas évidente :} les données économiques et ethno-religieuses montrent que le choix du Nigeria était rationnel pour une partie des populations ; l'histoire du Cameroun aurait pu être autre.
\item \textbf{Savoir-faire travaillés :} analyser des documents croisés, distinguer faits et opinions, construire une argumentation justifiée, rédiger dans un registre adapté et défendre oralement une position — compétences directement mobilisables à l'épreuve d'Histoire du Probatoire et du Baccalauréat technique (commentaire de document et composition).
\end{itemize}

L'enseignant évalue le rapport écrit et la restitution orale à l'aide d'une grille critériée portant sur quatre critères : la rigueur historique (faits, dates, acteurs exacts), la qualité de l'argumentation (appui sur les documents), le style diplomatique (neutralité, clarté) et le respect du mandat de l'ONU (consultation des populations, garanties aux minorités).

\newpage

\coursec{Méthodes et exercices résolus}

\courssub{Méthodes et exercices résolus}

\begin{methode}[Analyser un document historique : identifier la nature, le contexte et le message]
\textbf{Objectif :} Maîtriser la grille d'analyse documentaire pour répondre aux sujets de type « Commentaire de document » au Probatoire/Baccalauréat.
\begin{enumerate}
    \item \textbf{Identification (Qui ? Quand ? Quoi ?) :} Nature du document (texte officiel, discours, tract, article de presse, carte, photo), auteur, date, lieu, source.
    \item \textbf{Contextualisation (Pourquoi ?) :} Rattacher le document au chapitre précis (ex: Tutelle ONU, Insurrection 1955, Conférence de Foumban). Citer 2--3 faits majeurs contemporains.
    \item \textbf{Analyse du contenu (Message) :} Dégager l'idée principale (thèse de l'auteur), les arguments/mots-clés (\cle{indépendance}, \cle{réunification}, \cle{démocratie}, \cle{violence}), le vocabulaire spécifique.
    \item \textbf{Portée et Limites (Intérêt historique) :} Ce que le document \emph{apporte} à l'historien (témoignage, propagande, trace administrative) et ce qu'il \emph{ne dit pas} (silences, parti pris, public visé).
    \item \textbf{Conclusion synthétique :} Une phrase résumant l'apport du document pour comprendre la marche vers l'indépendance.
\end{enumerate}
\cle{Réflexe Bac :} Toujours citer le document (« L'auteur écrit : "..." ») et dater les faits évoqués.
\end{methode}

\begin{exoresolu}[Commentaire d'un extrait du Manifeste de l'UPC (avril 1955)]
\textbf{Énoncé :} À l'aide du document ci-dessous et de vos connaissances, montrez comment l'UPC traduit les aspirations nationalistes au Cameroun sous tutelle dans les années 1955.

\textit{Document : Extrait du Manifeste de l'UPC, avril 1955.}
« ... Le peuple camerounais veut être libre. Il veut gérer lui-même ses affaires. Il en a le droit, comme tous les peuples du monde. La tutelle française a failli à sa mission civilisatrice... Nous réclamons l'indépendance immédiate et la réunification du Cameroun... »

\tcblower
\textbf{Solution détaillée :}

\textbf{1. Identification}
\begin{itemize}
    \item \textbf{Nature :} Texte politique, tract/manifeste de parti (source partisane).
    \item \textbf{Auteur :} Union des Populations du Cameroun (UPC), parti nationaliste radical créé en 1948.
    \item \textbf{Date :} Avril 1955 (quelques semaines avant l'insurrection de mai 1955).
    \item \textbf{Contexte immédiat :} Cameroun sous tutelle française (ONU), montée des tensions, refus de l'administration d'accorder l'autonomie, interdiction de l'UPC (juillet 1955).
\end{itemize}

\textbf{2. Contextualisation historique}
Le Cameroun est un Territoire sous Tutelle (TST) confié à la France par l'ONU (1946). L'Article 76 de la Charte de l'ONU prévoit l'évolution vers l'autonomie/indépendance. Depuis 1948, l'UPC (Ruben Um Nyobè, Félix-Roland Moumié) revendique l'indépendance immédiate et la réunification avec le British Southern Cameroons. L'année 1955 marque un tournant : échec des négociations, radicalisation, insurrection (mai 1955), répression, interdiction du parti (13 juillet 1955), début de la guerre de libération nationale (maquis).

\textbf{3. Analyse du contenu (Idées forces \& Arguments)}
\begin{itemize}
    \item \textbf{Thèse centrale :} Revendication de la \cle{souveraineté nationale} (\textit{« Le peuple camerounais veut être libre... gérer lui-même ses affaires »}).
    \item \textbf{Argument juridique/légitime :} Droit des peuples à l'autodétermination (\textit{« Il en a le droit, comme tous les peuples du monde »}) $\rightarrow$ référence implicite à la Charte de l'ONU (Art. 1, 76).
    \item \textbf{Critique du pouvoir colonial :} Échec de la \cle{mission civilisatrice} française (\textit{« La tutelle française a failli »}) $\rightarrow$ délégitimation de l'autorité tutélaire.
    \item \textbf{Revendications concrètes :} \cle{Indépendance immédiate} (rupture avec l'évolutionnisme modéré) et \cle{réunification} (dimension territoriale pan-camerounaise).
\end{itemize}

\textbf{4. Portée et Limites}
\begin{itemize}
    \item \textbf{Portée :} Témoigne de la radicalisation du nationalisme (rupture avec les modérés type ARCA/AMDC). Montre la double revendication (indépendance + réunification) qui structure le conflit. Preuve de la mobilisation populaire visée.
    \item \textbf{Limites :} Document de propagande (langage mobilisateur, pas d'analyse socio-économique). Ne reflète pas l'avis de \emph{tous} les Camerounais (ex: chefferies traditionnelles, élites modérées, populations du Sud britannique). Silence sur les moyens d'action (légal vs insurrectionnel).
\end{itemize}

\textbf{5. Conclusion}
Ce manifeste de l'UPC d'avril 1955 constitue l'acte de naissance politique de la revendication d'indépendance \emph{immédiate} et de \emph{réunification}, marquant la rupture définitive avec l'ordre colonial français et préfigurant le basculement dans la clandestinité et la lutte armée après l'interdiction du parti en juillet 1955.
\end{exoresolu}

\begin{methode}[Construire une frise chronologique commentée : situer les ruptures et continuités]
\textbf{Objectif :} Visualiser la dynamique temporelle (1946--1972) et hiérarchiser les dates charnières pour une dissertation ou un QCM chronologique.
\begin{enumerate}
    \item \textbf{Délimiter l'intervalle :} 1946 (Tutelle ONU) -- 1972 (Référendum État unitaire). Ajouter 1961 (Indépendance/Réunification) comme pivot central.
    \item \textbf{Sélectionner 8--10 dates "charnières" (Ruptures) :} Ne pas tout mettre. Choisir les changements de régime, lois, traités, événements violents.
    \item \textbf{Structurer en périodes (Phases) :} Grouper les dates sous 3--4 étiquettes de phase (ex: \textit{« Tutelle \& Émergence nationaliste (1946-1955) »}, \textit{« Guerre \& Indépendance (1955-1960) »}, \textit{« Réunification \& Construction étatique (1961-1972) »}).
    \item \textbf{Commenter chaque date :} 1 date = 1 événement + 1 conséquence courte (cause $\rightarrow$ effet).
    \item \textbf{Code couleur / Symboles :} \textcolor{popblue}{Bleu} = Juridique/Institutionnel (lois, référendums) ; \textcolor{poporange}{Orange} = Violences/Conflits ; \textcolor{popgreen}{Vert} = Figures/Acteurs clés ; \textcolor{poppurple}{Violet} = Actes internationaux (ONU).
\end{enumerate}
\cle{À retenir :} Les dates « pivots » (1955, 1960, 1961, 1972) sont obligatoires au Bac.
\end{methode}

\begin{exoresolu}[Frise chronologique : Les étapes majeures de la souveraineté (1946-1972)]
\textbf{Énoncé :} Réalisez une frise chronologique commentée des 10 dates indispensables pour comprendre l'accession du Cameroun à la souveraineté nationale et la construction de l'État unitaire (1946-1972). Pour chaque date, indiquez l'événement et sa signification historique en une phrase.

\tcblower
\textbf{Solution détaillée : Frise commentée (Modèle de rédaction au Bac)}

\begin{center}
\begin{tabularx}{\linewidth}{|c|X|X|}\hline
\textbf{Date} & \textbf{\'Evénement (Acte / Fait)} & \textbf{Signification historique (Cause $\rightarrow$ Conséquence)} \\ \hline
\textcolor{poppurple}{13 déc. 1946} & Accords de tutelle ONU / France. & Institutionnalise le statut international ; obligation pour la France de mener à l'autonomie (Art. 76 Charte ONU). \\ \hline
\textcolor{popgreen}{10 avr. 1948} & Création de l'UPC (Um Nyobè). & Naissance du nationalisme radical, unitaire et panafricain ; rupture avec l'assimilationnisme. \\ \hline
\textcolor{poporange}{Mai--Juil. 1955} & Insurrection nationaliste \& Interdiction UPC. & Basculement dans la clandestinité et la lutte armée (maquis) ; fin de la légalité politique pour les radicaux. \\ \hline
\textcolor{popblue}{23 déc. 1956} & Loi-cadre Defferre (application au Cameroun). & Octroie l'autonomie interne ; élection 1er gouvernement africain (André-Marie Mbida, puis Ahmadou Ahidjo). \\ \hline
\textcolor{popblue}{1er janv. 1960} & Proclamation de l'Indépendance (République du Cameroun). & Fin de la tutelle française ; accession à la souveraineté internationale (adhésion ONU sept. 1960). \\ \hline
\textcolor{poppurple}{11 fév. 1961} & Plebiscite ONU (British Cameroons). & Nord : rattachement au Nigeria ; Sud : réunification avec la République du Cameroun. \\ \hline
\textcolor{popblue}{1er oct. 1961} & Indépendance du Southern Cameroons \& Réunification. & Naissance de la \cle{République Fédérale du Cameroun} (2 États fédérés : Ouest / Est). \\ \hline
\textcolor{poporange}{1961--1970} & Guerre de l'Ouest (Répression maquis UPC). & Élimination physique des leaders (Moumié 1960, Ouandié 1971) ; consolidation du pouvoir Ahidjo. \\ \hline
\textcolor{popblue}{20 mai 1972} & Référendum sur l'État unitaire. & Abolition du fédéralisme ; naissance de la \cle{République Unie du Cameroun} (centralisation). \\ \hline
\textcolor{popblue}{1984} & Décret Ahidjo : \textit{République du Cameroun} (nom actuel). & Effacement de la référence « Unie » ; achèvement symbolique de la construction étatique centralisée. \\ \hline
\end{tabularx}
\end{center}

\textbf{Conseil de rédaction :} Au Bac, dessinez une flèche horizontale. Placez les dates au-dessus (juridique/international) et au-dessous (interne/conflit). Utilisez les codes couleur.
\end{exoresolu}

\begin{methode}[Rédiger une dissertation historique : Plan dialectique (Thèse / Antithèse / Synthèse)]
\textbf{Objectif :} Répondre à un sujet de type « Les défis de la construction de l'État-nation camerounais (1961-1972) » avec une argumentation structurée, nuancée et documentée.
\begin{enumerate}
    \item \textbf{Analyser le sujet (Mots-clés, Bornes, Commande) :}
        \begin{itemize}
            \item \cle{Mots-clés} : « Défis », « Construction », « État-nation », « Réconciliation », « Mémoire ».
            \item \cle{Bornes} : 1961 (Réunification) -- 1972 (État unitaire) [voire 1984].
            \item \cle{Commande} : « Montrez que... », « Dans quelle mesure... », « Analysez... » $\rightarrow$ Exige une démonstration en \textbf{3 parties équilibrées}.
        \end{itemize}
    \item \textbf{Problématiser (La question centrale) :} Ex: « Comment le Cameroun a-t-il surmonté les contradictions nées de la réunification pour bâtir un État unitaire, au prix de quels silences mémoriels ? »
    \item \textbf{Construire le Plan détaillé (A / B / C) :}
        \begin{itemize}
            \item \textbf{I. Thèse : L'héritage conflictuel de la réunification (Les défis structurels).} Fédéralisme bancal, disparités juridiques/administratives, opposition UPC/Ahidjo, guerre cachée.
            \item \textbf{II. Antithèse : La centralisation autoritaire comme réponse (La construction étatique).} Référendum 1972, parti unique (UNC), intégration administrative, développement infrastructurel, « paix » apparente.
            \item \textbf{III. Synthèse : Un État fort, une nation inachevée ? (Le pari mémoriel).} Unité administrative acquise, mais fractures linguistiques (anglophone/francophone), mémorielles (guerre UPC tuée), identitaires non résolues $\rightarrow$ racine de la crise actuelle (depuis 2016).
        \end{itemize}
    \item \textbf{Rédiger l'Introduction (Entonnoir) :} Accroche $\rightarrow$ Définitions $\rightarrow$ Contexte $\rightarrow$ Problématique $\rightarrow$ Annonce du plan.
    \item \textbf{Rédiger le Développement :} 1 paragraphe = 1 argument + 1 exemple précis (date, chiffre, nom, loi) + 1 mini-conclusion de paragraphe. \cle{Transition} obligatoire entre I, II, III.
    \item \textbf{Rédiger la Conclusion :} Réponse directe à la problématique $\rightarrow$ Bilan nuancé $\rightarrow$ Ouverture (ex: Crise anglophone 2016+, Commission Vérité/Réconciliation).
\end{enumerate}
\cle{Formule Introduction :} « Depuis la réunification du 1er octobre 1961, le Cameroun peine à transformer son unité juridique en unité nationale... »
\end{methode}

\begin{exoresolu}[Dissertation : « La construction de l'État-nation camerounais de 1961 à 1972 : entre réunification et centralisation »]
\textbf{Énoncé :} Sujet de dissertation type Bac. « Montrez que la construction de l'État-nation camerounais (1961-1972) a été marquée par la tension entre l'idéal réunificateur et la réalité de la centralisation autoritaire. »

\tcblower
\textbf{Solution détaillée : Plan rédigé (Squelette complet pour l'examen)}

\textbf{INTRODUCTION}
\textbf{[Accroche]} « Le Cameroun est né deux fois : une première fois le 1er janvier 1960, une seconde fois le 1er octobre 1961. » (Paul Biya, discours investiture).
\textbf{[Définitions]} \cle{État-nation} : adéquation entre un peuple (nation) et une structure politique souveraine (État). \cle{Réunification} : jonction du Cameroun oriental (ex-français) et du Southern Cameroons (ex-britannique).
\textbf{[Contexte]} 1961 : Fédération biculturelle (Common Law / Droit civil ; Anglais / Français ; Parlementarisme / Présidentialisme). 1972 : Référendum, État unitaire, parti unique.
\textbf{[Problématique]} Comment le passage de la fédération à l'État unitaire a-t-il résolu les défis institutionnels de la réunification tout en générant de nouvelles fractures identitaires et mémorielles ?
\textbf{[Annonce]} Nous verrons d'abord que la fédération portait en elle les germes de la dislocation (I), puis que l'État unitaire s'est imposé par la force juridique et politique (II), pour enfin analyser l'héritage ambigu d'une unité administrative sans réconciliation mémorielle (III).

\textbf{DÉVELOPPEMENT}

\textbf{I. La Fédération (1961-1972) : Un compromis bancal face aux défis de la réunification}
\textbf{A. Une architecture institutionnelle dualiste source de blocages.}
Constitution du 1er sept. 1961 (Foumban). 2 États fédérés (Ouest / Est) + Fédéral. \textbf{Exemple :} Deux Premiers ministres (Foncha / Ahidjo), deux assemblées, deux systèmes judiciaires/éducatifs. \textbf{Conséquence :} Paralysie décisionnelle, coût financier énorme, sentiment d'inégalité (Ouest se sent marginalisé).
\textbf{B. L'opposition radicale : L'UPC, « l'oubliée » de la réunification.}
L'UPC (maquis) refuse la fédération (néocolonialisme). \textbf{Exemple :} Assassinat de Félix-Roland Moumié (Genève, 1960), exécution d'Ernest Ouandié (Bafoussam, 1971). \textbf{Conséquence :} Guerre civile larvée dans l'Ouest (Bamiléké) ; répression militaire (Commandant Loumpet) ; silence officiel sur ces morts.
\textbf{C. La dérive unitaire avant l'heure (Ahidjo).}
1966 : Parti unique (UNC = fusion KNDP + UC + autres). 1968 : Fin du poste de VP de la Fédération (Foncha). \textbf{Argument :} Ahidjo utilise la menace UPC pour recentraliser le pouvoir avant 1972.

\textbf{II. Le référendum du 20 mai 1972 : L'avènement de l'État unitaire centralisé}
\textbf{A. Une consultation controversée.}
Question : « Voulez-vous... un État unitaire ? » (pas de choix « Non » clair, bulletins uniques). \textbf{Chiffre :} 99,99\% « Oui » (résultat type parti unique). \textbf{Acteurs :} Ahidjo (partisan) vs Foncha / Muna / Yamure (opposants fédéralistes, mis en retraite/assignés à résidence).
\textbf{B. La refonte institutionnelle radicale (Constitution 2 juin 1972).}
Suppression : Fédéralisme, Vice-Présidence, Assemblées fédérées, Premier Ministère. Création : Présidence toute-puissante (chef de l'État, du govt, des armées), Assemblée nationale unique, Cour Suprême. \textbf{Objectif affiché :} Efficacité, unité nationale, développement économique (Plan quinquennal).
\textbf{C. L'intégration administrative et symbolique.}
Remplacement des administrators anglais par des francophones dans l'Ouest. Harmonisation codes (pénal, travail). Changement drapeau (2 étoiles $\rightarrow$ 1 étoile). Hymne : « Ô Cameroun, berceau de nos ancêtres » (version unique).

\textbf{III. Bilan 1972 : Une unité administrative acquise, une réconciliation nationale différée}
\textbf{A. La « paix » imposée par le parti unique.}
Stabilité politique (Ahidjo 22 ans). Développement infrastructures (routes, barrages, SONARA). Intégration cadres nord/sud, est/ouest dans l'administration unique.
\textbf{B. Les fractures mémorielles non refermées (Le « non-dit »).}
\textbf{1. Guerre UPC :} Pas de procès équitable, pas d'amnistie générale, pas de réhabilitation officielle des « maquisards » (qualifiés de « rebelles »/« bandits » jusqu'aux années 90). \textbf{2. Question anglophone :} Perte du \cle{Common Law}, de l'éducation bilingue de qualité, marginalisation politique (postes clés aux francophones). Germe de la crise 2016.
\textbf{C. La nation en chantier : Unité \emph{juridique} $\neq$ Unité \emph{nationale}.}
L'État est fort (centralisé), la nation est fragile (identités plurielles non reconnues). Ouverture : Commission Vérité et Réconciliation (CVR, 2019-2021) tente de rouvrir le dossier mémoriel (guerre UPC, crises anglophones).

\textbf{CONCLUSION}
\textbf{[Réponse]} 1961-1972 : Passage d'une unité « confédérale » négociée (Foumban) à une unité « jacobine » imposée (20 mai). La centralisation a « fait l'État » (administration unique, territoire continu) mais a « défait la nation » (exclusion UPC, frustration anglophone).
\textbf{[Ouverture]} La crise anglophone actuelle (depuis 2016) et les débats sur la « mémoire » de la guerre d'indépendance prouvent que le chantier de 1972 reste inachevé. La construction de l'État-nation est un processus continu, non un décret.
\end{exoresolu}

\begin{methode}[Réaliser un portrait croisé / fiche biographique comparative (Dossier 3 : Figures du nationalisme)]
\textbf{Objectif :} Maîtriser le Dossier 3 du programme. Comparer les parcours, stratégies et héritages des figures majeures pour une question de synthèse ou un oral.
\begin{enumerate}
    \item \textbf{Sélectionner les 4--5 figures incontournables :} Ruben Um Nyobè, Félix-Roland Moumié, Ernest Ouandié, Ahmadou Ahidjo, John Ngu Foncha. (Ajouter : Martin Paul Samba, Rudolf Duala Manga Bell pour l'héritage précoce ; André-Marie Mbida pour la transition).
    \item \textbf{Remplir une grille comparative standardisée (Tableau) :}
        \begin{tabular}{|l|c|c|c|c|c|}\hline
        \textbf{Critères} & \textbf{Um Nyobè} & \textbf{Moumié} & \textbf{Ouandié} & \textbf{Ahidjo} & \textbf{Foncha} \\ \hline
        \textbf{Origine / Base} & Sanaga-Maritime (Bassa) & Noun (Bamileke) & Haut-Nkam (Bamileke) & Nord (Fulbé/Musulman) & Nord-Ouest (Anglophone) \\ \hline
        \textbf{Parti / Rôle} & UPC (Fondateur, SG) & UPC (Président, exil) & UPC (Dernier VP, Maquis) & UC / UNC (1er PM, 1er PR) & KNDP (PM Ouest, VP Féd.) \\ \hline
        \textbf{Stratégie} & Légalisme $\rightarrow$ ONU $\rightarrow$ Maquis & Diplomatie internationale (Pékin, Accra, Le Caire) & Lutte armée intérieure (Maquis) & Légalisme institutionnel, alliance France & Fédéralisme, défense acquis anglophones \\ \hline
        \textbf{Revendication centrale} & Indépendance immédiate + Réunification & Idem + Panafricanisme & Idem + Justice sociale & Autonomie $\rightarrow$ Indépendance $\rightarrow$ Unité & Réunification \textbf{fédérale} \\ \hline
        \textbf{Fin tragique / Sort} & Tuissé par armée française (1958) & Empoisonné par SDECE (Genève, 1960) & Exécuté public (Bafoussam, 1971) & Démission 1982, exil, mort 1989 & Retraite forcée 1972, mort 1999 \\ \hline
        \textbf{Héritage / Mémoire} & « Père de l'indépendance » (réhabilité 1991) & Héros panafricain, martyr & Symbole résistance armée & « Père de la Nation » (bâtisseur État) & Défenseur identité anglophone \\ \hline
        \end{tabular}
    \item \textbf{Rédiger la synthèse comparative :} Opposer \cle{Nationalisme radical (UPC)} vs \cle{Nationalisme modéré/institutionnel (Ahidjo/Foncha)}. Montrer la complémentarité/antagonisme : L'UPC a \emph{imposé} l'agenda (indépendance immédiate) ; Ahidjo a \emph{géré} la souveraineté (État).
    \item \textbf{Argumenter avec des preuves :} Citations (Um Nyobè à l'ONU 1952), Actes (Loi 1956, Référendum 1972), Dates précises.
\end{enumerate}
\cle{Attention :} Ne pas confondre \textit{indépendance} (1960, Ahidjo/UC) et \textit{réunification} (1961, UPC/Foncha). L'UPC veut les deux \emph{ensemble} et \emph{immédiatement}.
\end{methode}

\begin{exoresolu}[Portraits croisés : Um Nyobè vs Ahidjo, deux stratégies pour une indépendance]
\textbf{Énoncé :} À l'aide d'un tableau comparatif et d'un paragraphe de synthèse, opposez les parcours, les méthodes et les héritages de Ruben Um Nyobè et d'Ahmadou Ahidjo dans l'accession du Cameroun à la souveraineté (1948-1960).

\tcblower
\textbf{Solution détaillée :}

\textbf{1. Tableau comparatif}

\begin{center}
\begin{tabularx}{\linewidth}{|l|X|X|}\hline
\textbf{Critères de comparaison} & \textbf{Ruben Um Nyobè (UPC - Radical)} & \textbf{Ahmadou Ahidjo (UC/UNC - Modéré)} \\ \hline
\textbf{Origine sociale / Base} & Instituteur, syndicaliste (Bassa, Sanaga-Maritime). Ancré dans la paysannerie et les syndicats (USCC). & Fils de notable, cadre administration coloniale (Fulbé, Nord). Ancré dans chefferies traditionnelles, élites nordistes, administration. \\ \hline
\textbf{Vision de l'indépendance} & \textbf{Immédiate, totale, inconditionnelle.} Rupture nette avec la France. Non-négociable. & \textbf{Progressive, négociée, dans la coopération.} « Indépendance dans l'interdépendance » (liens France). \\ \hline
\textbf{Rapport à la Tutelle / ONU} & Utilise la tribune de l'ONU (1952, 1953, 1954) pour dénoncer la France. « Accuseur » du pouvoir tutélaire. & Travaille \emph{dans} les institutions de la Tutelle (ARCA, ATCAM, Assemblée législative). « Partenaire » de l'Administration. \\ \hline
\textbf{Stratégie d'action} & Mobilisation de masse (grèves, meetings), Pétitions ONU, puis Maquis (clandestinité 1955). & Action parlementaire, Négociations (Loi-cadre 1956, Conférences constitutionnelles), Alliances politiques. \\ \hline
\textbf{Rapport à la Réunification} & Cœur du projet : « Un Cameroun, un Peuple, une Nation » (intégral, y compris Nord britannique). & Secondaire au début (priorité à l'indépendance de l'Est). Accepte plébiscite 1961 (perte Nord). \\ \hline
\textbf{Sort personnel} & Tuissé par l'armée française (13 sept 1958) $\rightarrow$ \cle{Martyr}. & 1er Premier Ministre (1958), 1er Président (1960-1982) $\rightarrow$ \cle{Bâtisseur de l'État}. \\ \hline
\textbf{Héritage officiel} & « Héros national » (loi 1991), mais mémoire longtemps refoulée. Symbole de la \cle{radicalité}. & « Père de la Nation », « Architecte de l'unité ». Symbole de la \cle{stabilité} et de l'autoritarisme. \\ \hline
\end{tabularx}
\end{center}

\textbf{2. Paragraphe de synthèse (Modèle rédaction Bac)}

Si tous deux sont les « pères » de l'indépendance camerounaise, \textbf{Ruben Um Nyobè} et \textbf{Ahmadou Ahidjo} en incarnent les deux visages antagonistes. Um Nyobè, par sa radicalité morale et son recours à la légalité internationale (ONU), \textbf{impose l'idée même d'indépendance immédiate} comme un droit inaliénable, payant de sa vie (1958) la rupture avec l'ordre colonial. Ahidjo, pragmatique, use des institutions de la tutelle (Loi-cadre Defferre, 1956) pour \textbf{conquérir le pouvoir d'État par les urnes} et négocier le transfert de souveraineté (1er janv. 1960). 

Le paradoxe historique est que \textbf{l'UPC d'Um Nyobè a « gagné » le combat des idées (indépendance + réunification)} — le programme d'Ahidjo s'aligne in fine sur celui de l'UPC (Réunification 1961) —, alors qu'\textbf{Ahidjo a « gagné » le combat du pouvoir}, éliminant physiquement (Moumié, Ouandié) ou politiquement (interdiction UPC 1955) ses rivaux. La construction de l'État-nation camerounais porte ainsi la double empreinte : l'exigence éthique de l'indépendance (Um Nyobè) et la realpolitik de la gestion étatique (Ahidjo).
\end{exoresolu}

\begin{methode}[Analyser un croquis / une carte schématique : L'évolution territoriale et administrative]
\textbf{Objectif :} Lire et produire un croquis de géographie historique (évolution des frontières, capitales, zones de conflit, axes de communication) pour illustrer une dissertation ou répondre à un sujet de cartographie.
\begin{enumerate}
    \item \textbf{Choisir le fond de carte simplifié :} Contour du Cameroun actuel (triangle approximatif). Repères : Lac Tchad (N), Golfe de Guinée (SW), Fleuves (Bénoué, Sanaga, Nyong, Logone).
    \item \textbf{Définir la légende (3 à 4 items max) :} Utiliser des \textbf{surfaces} (zones), des \textbf{points} (villes), des \textbf{flèches} (flux/mouvements), des \textbf{traits} (frontières/axes).
    \item \textbf{Respecter la sémiologie graphique :}
        \begin{itemize}
            \item \textit{Surfaces (hachures/couleurs claires) :} Zones (ex: Zone de tutelle française / britannique ; Zones de maquis UPC ; Régions anglophones/francophones).
            \item \textit{Points (cercles proportionnels) :} Capitales (Yaoundé, Buea, Garoua), Lieux de mémoire (Bafoussam, Eseka, Foumban).
            \item \textit{Flèches (épaisseur = intensité) :} Réunification (flèche Nord $\rightarrow$ Sud), Répression (flèches vers maquis), Migration/Échanges.
            \item \textit{Traits :} Frontière 1916 (Picot), Frontière 1919 (Millet-Simon), Ligne de partage Eau (partage des eaux).
        \end{itemize}
    \item \textbf{Titrer précisément :} « Le Cameroun : de la partition (1919) à l'État unitaire (1972) » ou « Foyers de la résistance nationaliste (1955-1971) ».
    \item \textbf{Intégrer au texte :} « Comme le montre le croquis ci-contre, la zone de maquis (hachures orange) coïncide avec les terres bamiléké/bassa... »
\end{enumerate}
\cle{Règle d'or :} Une carte sans titre, sans légende, sans échelle (indicative) et sans flèche Nord vaut 0 point.
\end{methode}

\begin{exoresolu}[Croquis : Les deux temps de la réunification (1919 / 1961)]
\textbf{Énoncé :} Réalisez un croquis légendé intitulé « Le Cameroun : de la partition coloniale (1919) à la réunification (1961) ». Montrez les frontières successives, les capitales administratives et le résultat du plébiscite de 1961.

\tcblower
\textbf{Solution détaillée : Description du croquis attendu (Code TikZ indicatif + Légende)}

\textbf{Code TikZ (Structure simplifiée pour l'examen - à dessiner à main levée proprement) :}
\begin{popfigure}
\begin{tikzpicture}[scale=0.8]
% Contour Cameroun actuel (polygone très simplifié)
\draw[popdark, thick] (0,0) -- (6,0.5) -- (7,4) -- (5.5,6) -- (2,6.5) -- (0.5,4) -- cycle;
% Frontière 1919 (Millet-Simon) : trait pointillé vertical-ish
\draw[poporange, dashed, thick] (3.5,0.3) -- (4.5,5.5) node[above right, font=\tiny] {Frontière 1919};
% Zone Française (Est) - Hachures Nord-Ouest
\pattern[pattern=north west lines, pattern color=popblue!30] (0,0) -- (6,0.5) -- (7,4) -- (5.5,6) -- (2,6.5) -- (0.5,4) -- cycle;
% Zone Britannique (Ouest) - Hachures Nord-Est (simulé sur partie gauche)
\pattern[pattern=north east lines, pattern color=poporange!30] (0,0) -- (3.5,0.3) -- (4.5,5.5) -- (0.5,4) -- cycle;
% Capitales
\fill[popblue] (2.5,3) circle (3pt) node[left, font=\small] {Yaoundé (Fr)};
\fill[poporange] (4.5,4.5) circle (3pt) node[right, font=\small] {Buea (Br)};
\fill[popdark] (3.5,1.5) circle (3pt) node[below, font=\small] {Douala};
% Flèche Réunification 1961
\draw[popgreen, very thick, ->] (5.5, 2) -- (2.5, 3) node[midway, above, font=\tiny, sloped] {Réunification 1961};
% Plébiscite 1961 : Flèches Nord (Nigeria) / Sud (Cameroun)
\draw[poppurple, thick, ->] (5.5, 5.5) -- (7, 6.5) node[above right, font=\tiny] {Nord $\rightarrow$ Nigeria};
\draw[popgreen, thick, ->] (4.5, 4.5) -- (2.5, 3) node[left, font=\tiny] {Sud $\rightarrow$ Cameroun};
% Nord
\draw (6.5, -0.5) node[rotate=90, font=\tiny] {N};
\end{tikzpicture}
\legende{Croquis schématique : Partition 1919 (trait pointillé orange) et Réunification 1961 (flèches). Hachures bleues = Zone française (Est). Hachures orange = Zone britannique (Ouest). Points = Capitales administratives.}
\end{popfigure}

\textbf{Légende structurée (à écrire à côté du dessin) :}
\begin{itemize}
    \item \textbf{Fond de carte :} Contour actuel du Cameroun (traits pleins noirs).
    \item \textbf{Limites :} \textcolor{poporange}{Trait pointillé épais} = Frontière Millet-Simon (1919) partageant le Kamerun.
    \item \textbf{Surfaces (Zones) :}
        \begin{itemize}
            \item \tikz[baseline]\fill[pattern=north west lines, pattern color=popblue!30] (0,0) rectangle (0.3,0.2); Zone Française (Cameroun Oriental) : Capitale \textcolor{popblue}{Yaoundé}.
            \item \tikz[baseline]\fill[pattern=north east lines, pattern color=poporange!30] (0,0) rectangle (0.3,0.2); Zone Britannique (Cameroons) : Capitale \textcolor{poporange}{Buea}. Divisée en Northern / Southern Cameroons.
        \end{itemize}
    \item \textbf{Flèches (Dynamiques 1961) :}
        \begin{itemize}
            \item \textcolor{poppurple}{Flèche épaisse vers l'Ouest (Nigeria)} : Rattachement du Northern Cameroons (Juin 1961).
            \item \textcolor{popgreen}{Flèche épaisse vers l'Est (Yaoundé)} : Réunification du Southern Cameroons (1er Oct 1961).
        \end{itemize}
    \item \textbf{Symboles ponctuels :} \textcolor{popblue}{\textbullet} Yaoundé ; \textcolor{poporange}{\textbullet} Buea ; \textcolor{popdark}{\textbullet} Douala (capitale économique).
\end{itemize}

\textbf{Commentaire type (3 lignes) :} « Ce croquis illustre la double rupture : la partition arbitraire de 1919 (frontière Millet-Simon) qui divise les peuples (Bamileke, Fulbé, Mafa), et la réunification partielle de 1961 où le Nord choisit le Nigeria, seul le Sud rejoignant la République du Cameroun, créant l'asymétrie fédéraliste (2 États inégaux) source de tensions futures. »
\end{exoresolu}

\begin{methode}[Argumenter à partir de sources orales / Mémoire : La méthode de l'historien face aux « non-dits »]
\textbf{Objectif :} Traiter les questions sur la « mémoire » et la « réconciliation » (partie 5 du plan) en croisant sources officielles, témoignages oraux et archives.
\begin{enumerate}
    \item \textbf{Identifier les « lieux de mémoire » (Pierre Nora) :} Monuments (Monument de la Réunification, Yaoundé ; Statue Um Nyobè, Eseka), Dates commémoratives (20 mai, 1er janv, 1er oct), Noms de rues/lycées.
    \item \textbf{Distinguer 3 niveaux de mémoire :}
        \begin{itemize}
            \item \cle{Mémoire officielle (État) :} Héroïsation d'Ahidjo, « Paix » 1972, Silences sur guerre UPC (qualifiée de « rébellion »).
            \item \cle{Mémoire des vainqueurs (Elites francophones) :} Narratif de la « construction nationale » réussie.
            \item \cle{Mémoire des vaincus / minoritaires (UPC, Anglophones) :} Revendication de reconnaissance (Um Nyobè « héros »), Justice pour crimes (Bafoussam, 1971), Revendication fédéralisme/autonomie.
        \end{itemize}
    \item \textbf{Croiser les sources (Triangulation) :}
        \begin{itemize}
            \item Archives coloniales (SHD Vincennes, Archives nationales Yaoundé) : Rapports militaires, correspondance gouverneur.
            \item Archives ONU (New York) : Rapports missions de visite, pétitions UPC.
            \item \textbf{Sources orales :} Entretiens avec anciens maquisards, familles victimes, administrateurs. \textbf{Éthique :} Recoupement, anonymat si risque, contextualisation (mémoire sélective, oubli, reconstruction a posteriori).
        \end{itemize}
    \item \textbf{Rédiger l'argumentation :} « Si le discours officiel célèbre l'unité (20 mai), l'analyse des sources orales des rescapés du maquis (ex: témoignages collected par l'historien M. Ngoh) révèle une mémoire traumatique de la répression (napalm, regroupements de villages), exclue du récit national jusqu'aux années 1990. »
\end{enumerate}
\cle{Compétence Bac :} « Montrez que l'histoire officielle a occulté... » $\rightarrow$ Utiliser le verbe \textbf{occulter}, \textbf{silencier}, \textbf{réhabiliter}, \textbf{instrumentaliser}.
\end{methode}

\begin{exoresolu}[Analyse de sources : L'instrumentalisation de la mémoire d'Um Nyobè (1958 - Aujourd'hui)]
\textbf{Énoncé :} À partir des trois extraits ci-dessous, montrez comment la mémoire de Ruben Um Nyobè a été instrumentalisée par le pouvoir et réappropriée par la société civile.

\textit{Doc 1 (1958, Communiqué officiel français) :} « Le bandit Um Nyobè, chef des rebelles communistes, a été abattu lors d'une opération de police. »
\textit{Doc 2 (1972, Discours Ahidjo, 20 mai) :} « Nous avons réalisé l'unité nationale. Les troubler-fêtes ont été neutralisés. Le Cameroun marche. »
\textit{Doc 3 (1991, Loi n°91/022) :} « Est déclaré héros national tout Camerounais qui s'est illustré par son action en faveur de l'indépendance... Ruben Um Nyobè est réhabilité. »

\tcblower
\textbf{Solution détaillée :}

\textbf{1. Analyse comparative des sources}

\begin{center}
\begin{tabularx}{\linewidth}{|l|X|X|X|}\hline
\textbf{Critères} & \textbf{Doc 1 (1958 : Colonial)} & \textbf{Doc 2 (1972 : Ahidjo)} & \textbf{Doc 3 (1991 : Biya / Démocratie)} \\ \hline
\textbf{Qualification} & « Bandit », « Rebelle communiste » & « Troubler-fêtes » (euphémisme), « Neutralisés » & « Héros national », « Action en faveur de l'indépendance » \\ \hline
\textbf{Intention / Fonction} & \textbf{Délégitimer} : Criminaliser la lutte armée, justifier la répression (guerre sale), effacer la dimension politique. & \textbf{Occulter / Unifier par le silence} : Gommer le conflit interne pour légitimer l'État unitaire (parti unique). La « paix » = absence d'opposition. & \textbf{Réhabiliter / Apaiser} : Répondre aux revendications de la société civile (années 90, « Villes mortes », Conférence nationale avortée). Geste mémoriel. \\ \hline
\textbf{Contexte politique} & Guerre froide, lutte anti-communiste, fin tutelle. & Consolidation pouvoir personnel, État unitaire, parti unique. & Libéralisation politique (multipartisme), pression opposition (SDF), enjeu anglophone. \\ \hline
\end{tabularx}
\end{center}

\textbf{2. Synthèse argumentée (Réponse type Bac)}

La mémoire de Ruben Um Nyobè illustre parfaitement l'\textbf{instrumentalisation politique de l'histoire} au Cameroun.

Dans un premier temps (Doc 1), le pouvoir colonial \textbf{criminalise} le nationaliste radical pour mieux délégitimer sa revendication d'indépendance immédiate. L'étiquette « communiste » (Guerre froide) vise à rallier l'opinion occidentale et à justifier une répression militaire hors cadre judiciaire (exécution sommaire, exposition du corps à Eseka).

Dans un second temps (Doc 2), le pouvoir postcolonial (Ahidjo) pratique l'\textbf{occultation mémorielle}. En qualifiant les maquisards de simples « trouble-fêtes » neutralisés, le discours officiel efface la guerre civile (1955-1971) pour fonder la légitimité de l'État unitaire sur le consensus apparent (« Le Cameroun marche »). Um Nyobè devient un « non-personnage » de l'histoire officielle, absent des manuels, des rues, des monuments.

Enfin (Doc 3), l'ouverture démocratique des années 1990 contraint le régime (Biya) à une \textbf{réhabilitation juridique et symbolique} (Loi 91/022). C'est une réappropriation par le haut, sous pression de la société civile (associations « Mémoire 58 », historiens, diaspora), pour tenter de solder le contentieux mémoriel. Toutefois, cette réhabilitation reste souvent formelle : absence de réparations aux familles, silence sur les responsabilités françaises, non-intégration complète dans les programmes scolaires jusqu'à récemment.

\textbf{Conclusion :} Le parcours mémoriel d'Um Nyobè (Bandit $\rightarrow$ Invisible $\rightarrow$ Héros) révèle que l'histoire nationale camerounaise est un champ de bataille où l'État gère sa légitimité en sélectionnant ses « pères fondateurs » selon les enjeux du moment présent.
\end{exoresolu}

\begin{methode}[Répondre à une question de synthèse / QCM argumenté : Mobiliser les connaissances précises]
\textbf{Objectif :} Réussir les questions courtes (QCM, Vrai/Faux justifié, Questions de cours directes) qui testent la précision factuelle (dates, noms, lois, articles).
\begin{enumerate}
    \item \textbf{Lire la question 2 fois :} Repérer le \textbf{mot-clé opératoire} (Citer, Expliquer, Comparer, Pourquoi, Conséquence).
    \item \textbf{Activer le « tiroir » mental correspondant :} (Tutelle ? Nationalisme ? Figures ? Indépendance ? Réunification ? État unitaire ?).
    \item \textbf{Appliquer la règle des « 3 P » pour la réponse courte :}
        \begin{itemize}
            \item \textbf{Précision :} Date exacte, Nom propre exact, Chiffre, Article de loi.
            \item \textbf{Pertinence :} Répondre \emph{uniquement} à la question posée (pas de hors-sujet).
            \item \textbf{Preuve :} Citer la source ou le fait qui valide l'affirmation.
        \end{itemize}
    \item \textbf{Structurer la réponse :} Phrase affirmative $\rightarrow$ Argument/Exemple $\rightarrow$ (Si « Pourquoi ») Mécanisme causal.
    \item \textbf{Relire :} Vérifier l'orthographe des noms propres (Um Nyobè, Moumié, Ouandié, Ahidjo, Foncha, Mbida, Foumban, Bafoussam, Eseka).
\end{enumerate}
\cle{Pièges classiques :} Confondre \textit{Indépendance} (1er janv 1960) et \textit{Réunification} (1er oct 1961). Confondre \textit{Fédéralisme} (1961) et \textit{État unitaire} (1972). Oublier le plébiscite du Nord (Nigeria).
\end{methode}

\begin{exoresolu}[Questions de synthèse / QCM argumenté : Vérification des acquis]
\textbf{Énoncé :} Répondez de manière précise et justifiée aux questions suivantes (niveau Bac).
\begin{enumerate}
    \item Quelle est la date exacte de l'entrée en vigueur de la Tutelle de l'ONU sur le Cameroun français ? Quel article de la Charte de l'ONU la fonde ?
    \item Citez les deux partis politiques modérés légalistes qui ont participé aux élections de 1952 et 1956 au Cameroun sous tutelle française.
    \item Pourquoi l'UPC boycotte-t-elle les élections de 1956 (Loi-cadre Defferre) ?
    \item Quelle est la question exacte posée lors du plébiscite du 11 février 1961 au Southern Cameroons ?
    \item Citez 3 différences institutionnelles majeures entre la Constitution fédérale de 1961 et la Constitution de l'État unitaire de 1972.
\end{enumerate}

\tcblower
\textbf{Solution détaillée :}

\begin{enumerate}
    \item \textbf{13 décembre 1946} (Accord de tutelle signé à New York, approuvé par l'Assemblée générale de l'ONU le 13 déc. 1946, entrée en vigueur). Fondement : \textbf{Article 76 de la Charte de l'ONU} (Chapitre XII : Système de tutelle) $\rightarrow$ Objectif : « favoriser l'évolution progressive [des territoires] vers l'autonomie ou l'indépendance ».
    \item \textbf{L'ARCA} (Assemblée Républicaine du Cameroun, André-Marie Mbida) et \textbf{L'ATCAM} (Alliance Traditionnelle Camerounaise, Ahmadou Ahidjo). *Note : L'UPC est interdite en juil. 1955, ne participe pas à 1956.*
    \item \textbf{3 raisons principales :} (1) La loi-cadre ne donne que l'\emph{autonomie interne} (pas l'indépendance), pouvoir résiduel à la France (Défense, Diplomatie, Monnaie). (2) Le mode de scrutin (collèges, suffrage restreint) favorise les notables/modérés. (3) L'UPC est \textbf{interdite} (décret 13 juil. 1955) et ses leaders en clandestinité/maquis ; participation impossible.
    \item \textbf{« Désirez-vous l'indépendance en rejoignant la République fédérale du Cameroun ? »} (Deux choix : OUI / NON). *Précision : Au Northern Cameroons, la question était : « Désirez-vous l'indépendance en rejoignant la Fédération du Nigeria ? »*
    \item \textbf{1. Structure :} Fédéralisme (2 États : Ouest/Est + Fédéral) $\rightarrow$ Unitarisme (1 État, 10 provinces puis régions). \textbf{2. Exécutif :} Président Fédéral + 2 Premiers Ministres (Vice-Président Fédéral) $\rightarrow$ Président de la République unique (Chef de l'État, du Govt, des Armées), suppression PM/VP. \textbf{3. Législatif :} Assemblée Fédérale + 2 Assemblées Fédérées $\rightarrow$ Assemblée Nationale unique (unicaméral). \textbf{4. Judiciaire :} Cour Fédérale + Cours Fédérées $\rightarrow$ Cour Suprême unique. \textbf{5. Parti :} Multipartisme (de fait UC/KNDP) $\rightarrow$ Parti Unique (UNC) constitutionnalisé (Art. 14 const. 72).
\end{enumerate}
\end{exoresolu}

\begin{attention}[Les 5 erreurs qui coûtent des points au Bac (Histoire - Terminale Tech)]
\begin{enumerate}
    \item \textbf{Confondre les bornes chronologiques et les statuts juridiques :} Écrire « Indépendance en 1961 » (c'est la réunification) ou « Fin de la tutelle en 1972 » (fin en 1960). \textbf{Reflexe :} \cle{1er janv 1960 = Indépendance / Fin Tutelle} ; \cle{1er oct 1961 = Réunification / Indépendance Sud} ; \cle{20 mai 1972 = État Unitaire}.
    \item \textbf{Traiter l'UPC comme un bloc monolithique ou « communiste » sans nuance :} Dire « L'UPC était communiste » (réduction propagande coloniale). \textbf{Reflexe :} L'UPC est \cle{nationaliste, radical, panafricain, marxisant (après 1955)}. Distinguer Um Nyobè (légaliste ONU), Moumié (diplomatie internationale), Ouandié (maquis intérieur). Ne pas oublier l'aile modérée (ARCA, UC, KNDP).
    \item \textbf{Oublier la dimension internationale (ONU) :} Traiter l'histoire comme une affaire purement franco-camerounaise. \textbf{Reflexe :} Citer systématiquement : \cle{Missions de visite ONU} (1949, 1952, 1955, 1958), \cle{Pétitions Um Nyobè} (1952, 53, 54), \cle{Plébiscite 1961} (supervisé ONU). C'est la spécificité du Cameroun (Tutelle, pas Colonie).
    \item \textbf{Rédiger une dissertation sans problématique ni plan visible :} Réciter le cours en vrac (« Ensuite... Ensuite... »). \textbf{Reflexe :} \textbf{Introduction} (Accroche, Déf, Problématique, Plan) \textbf{OBLIGATOIRE}. \textbf{Transitions} explicites entre I, II, III. \textbf{Conclusion} (Bilan + Ouverture). Paragraphes : 1 idée = 1 paragraphe = 1 exemple daté.
    \item \textbf{Négliger la partie « Mémoire / Réconciliation » (Partie 5 du programme) :} Croire que l'histoire s'arrête en 1972 ou 1984. \textbf{Reflexe :} Maîtriser le vocabulaire : \cle{Justice transitionnelle}, \cle{Réhabilitation (Loi 1991)}, \cle{Commission Vérité Réconciliation (2019)}, \cle{Crise anglophone (2016-)}, \cle{Lieux de mémoire}. L'examinateur attend une analyse critique du rapport au passé.
\end{enumerate}
\end{attention}

\newpage

\coursec{Exercices}

\courssub{Je vérifie mes connaissances}

\begin{exercice}[QCM : La Tutelle et les institutions]
Parmi les affirmations suivantes, une seule est exacte. La recopier en la complétant.
\begin{enumerate}
\item L'Organisation des Nations Unies (ONU) place le Cameroun sous tutelle en \textbf{1946} / 1950 / 1955.
\item La puissance administrante du Cameroun français est la \textbf{France} / la Grande-Bretagne / la Belgique.
\item La loi-cadre \textbf{Defferre} / Cadre / d'Outre-mer de 1956 instaure le suffrage universel et l'autonomie interne.
\item L'Union des Populations du Cameroun (UPC) est fondée en \textbf{1948} / 1952 / 1955 par Ruben Um Nyobé et d'autres nationalistes.
\end{enumerate}
\end{exercice}

\begin{exercice}[Vrai / Faux justifié : Le nationalisme et l'UPC]
Dire si les affirmations sont vraies ou fausses. Justifier chaque réponse par une phrase précise (date, acteur, fait).
\begin{enumerate}
\item L'UPC revendique d'abord l'indépendance immédiate et la réunification des deux Cameroun dès sa création en 1948.
\item Ruben Um Nyobé intervient à la IV\textsuperscript{e} Commission de l'ONU à New York en 1952 et 1954 pour défendre la cause camerounaise.
\item L'insurrection de mai 1955 est déclenchée par l'UPC après son interdiction par le Haut-commissaire Roland Pré en juillet 1955.
\item Félix-Roland Moumié succède à Um Nyobé à la tête de l'UPC (Secrétaire général) après l'assassinat de ce dernier le 13 septembre 1958.
\end{enumerate}
\end{exercice}

\begin{exercice}[Définitions et repères chronologiques]
Définir précisément les termes suivants en les situant dans le temps (date) et l'espace (contexte camerounais) :
\begin{enumerate}
\item \cle{Tutelle de l'ONU} (différence avec le mandat de la SDN) : Statut international créé en 1946 (Accords de tutelle signés 1946, approuvés 1947) remplaçant le mandat de la SDN ; obligation pour la puissance administrante de préparer le territoire à l'autonomie ou l'indépendance, sous contrôle de l'ONU (missions de visite, rapports annuels), contrairement au mandat qui n'imposait qu'un rapport annuel sans perspective d'indépendance fixe.
\item \cle{Loi-cadre Defferre} (23 juin 1956) : Loi française (Gaston Defferre, ministre de la France d'Outre-mer) instaurant le suffrage universel, l'autonomie interne (Assemblée territoriale, Conseil de gouvernement responsable devant elle) et posant le cadre de la future indépendance dans l'Empire français.
\item \cle{« Guerre cachée » ou « maquis »} (1955--1971) : Conflit armé opposant l'UPC (maquisards) aux forces françaises puis camerounaises, principalement dans le Sanaga-Maritime, le Bassa, le Bamileke et le Mungo ; répression massive, regroupements de populations, assassinats des leaders (Um Nyobé 1958, Moumié 1960, Ouandié 1971).
\item \cle{Référendum du 11 février 1961} : Scrutin organisé par l'ONU au Cameroun britannique sous supervision britannique. Question : « Voulez-vous réaliser l'indépendance en rejoignant la République fédérale du Cameroun ? » (Sud) / « ... en rejoignant la Fédération du Nigeria ? » (Nord). Le Nord vote l'intégration au Nigeria (60\%), le Sud la réunification (79\%).
\end{enumerate}
\end{exercice}

\begin{exercice}[Calculs directs et données chiffrées]
À partir des données ci-dessous, effectuer les calculs demandés.
\begin{itemize}
\item \textbf{Élections législatives du 23 décembre 1956 (Cameroun français) :} Inscrits : 1\,200\,000 ; Votants : 720\,000 ; Voix UPC (interdite, boycott) : 0 ; Voix UC (Union Camerounaise, Ahidjo) : 300\,000.
\item \textbf{Référendum du 11 février 1961 :}
      \begin{center}
      \begin{tabular}{lcc}
      & Nord camerounais & Sud camerounais \\ \hline
      Inscrits & 180\,000 & 250\,000 \\
      Votants & 170\,000 & 230\,000 \\
      « Intégration au Nigeria » & 60\% & 21\% \\
      « Réunification au Cameroun » & 40\% & 79\% \\
      \end{tabular}
      \end{center}
\end{itemize}
\begin{enumerate}
\item Calculer le taux de participation aux élections de 1956.
\item Calculer le nombre de voix pour la réunification au Sud.
\item Relever le pourcentage de voix pour l'intégration au Nigeria au Nord.
\item En déduire le choix majoritaire de chaque zone et la conséquence territoriale immédiate.
\end{enumerate}
\end{exercice}

\courssub{Je m'entraîne}

\begin{exercice}[Analyse de document : Discours d'Um Nyobé à l'ONU (1952)]
\textbf{Document :} Extrait de la pétition de Ruben Um Nyobé à la IV\textsuperscript{e} Commission de l'ONU, New York, décembre 1952.
\begin{quote}
« Le peuple camerounais réclame l'indépendance immédiate de son pays. [...] La France a violé les termes de l'accord de tutelle en refusant d'associer les véritables représentants du peuple à l'administration du territoire. [...] Nous demandons le départ des troupes françaises et l'organisation d'élections libres sous contrôle de l'ONU pour une Assemblée constituante. »
\end{quote}
\begin{enumerate}
\item Identifier l'auteur, la nature, la date et le destinataire du document.
\item Relever deux griefs formulés par l'UPC contre la puissance administrante.
\item Expliquer la revendication « élections libres sous contrôle de l'ONU ».
\item Quelle est la portée historique de cette intervention ? (Réponse en 3 lignes maximum).
\end{enumerate}
\end{exercice}

\begin{exercice}[Cartographie : Le référendum de 1961]
Sur le croquis simplifié du Cameroun britannique ci-dessous (fourni en annexe), réaliser la légende organisée suivante :
\begin{enumerate}
\item Tracer la limite entre Nord et Sud camerounais (ligne pointillée).
\item Hachurer en rouge la zone ayant voté « Intégration au Nigeria » (Nord).
\item Hachurer en vert la zone ayant voté « Réunification » (Sud).
\item Positionner les villes de Garoua, Bamenda, Buea, Kumba, Mamfe.
\item Tracer une flèche symbolisant le rattachement du Sud à la République du Cameroun (Yaoundé).
\item Titre, échelle indicative, orientation (flèche Nord), légende complète.
\end{enumerate}
\textit{Indication : Utiliser le fond de carte fourni par l'enseignant.}
\end{exercice}

\begin{exercice}[Tableau comparatif : Les grandes figures du nationalisme (Dossier 3)]
Compléter le tableau synthétique ci-dessous en mobilisant les connaissances du Dossier 3.

\begin{center}
\begin{tabularx}{\linewidth}{|l|X|X|X|X|}
\hline
\rowcolor{popblueL}
\textbf{Figure} & \textbf{Formation / Débuts politiques} & \textbf{Parti / Mouvement dirigé} & \textbf{Revendications principales} & \textbf{Sort / Fin de parcours} \\ \hline
Ruben Um Nyobé & Instituteur (École normale de Foulassi) ; syndicaliste USCC/CGT & UPC (fondateur, SG 1948) & Indépendance immédiate, réunification, réformes sociales, non-violence initiale & Assassiné par l'armée française le 13 sept. 1958 (maquis Sanaga-Maritime) \\ \hline
Félix-Roland Moumié & Médecin (Faculté de médecine de Dakar) ; jeunesse UPC & UPC (SG en exil 1958--1960) & Poursuite lutte armée, internationalisation (Pékin, Accra, Conakry), réunification & Empoisonné (thallium) par les services secrets français le 3 nov. 1960 à Genève \\ \hline
Ernest Ouandié & Instituteur ; vice-président UPC & UPC (chef militaire maquis, 1958--1970) & Lutte armée jusqu'à l'indépendance réelle, réunification, justice sociale & Capturé août 1970, procès public, exécuté par peloton d'exécution le 15 janv. 1971 (Bafoussam) \\ \hline
Ahmadou Ahidjo & École fils de chefs, École normale Yaoundé ; fonctionnaire, député & UC (1958), UNC (1966) & Autonomie interne (1956), indépendance dans l'ordre et la paix (1960), unité nationale & 1er Président de la République (1960--1982), démission 1982, mort en exil 1989 \\ \hline
John Ngu Foncha & Instituteur, études au Nigeria/UK ; syndicaliste & KNDP (Sud, 1955), Premier Ministre Ouest (1959) & Réunification par fédéralisme, autonomie régionale, droits minorités & VP Fédéral (1961--1966), PM Ouest (1959--1968), retraite, mort 2010 \\ \hline
\end{tabularx}
\end{center}
\end{exercice}

\begin{exercice}[Savoir-faire : Argumentation historique]
« L'accession du Cameroun à la souveraineté internationale (1\textsuperscript{er} janvier 1960) est le fruit d'une conjonction entre la lutte armée de l'UPC et la voie négociée d'Ahmadou Ahidjo. »
À l'aide d'exemples précis (dates, accords, acteurs), nuancer cette affirmation en distinguant l'indépendance du Cameroun français (1960) et la réunification (1961). Réponse structurée en deux paragraphes argumentés (environ 15 lignes au total).
\end{exercice}

\courssub{Je me prépare au Bac}

\begin{exercice}[Sujet type Baccalauréat : Dissertation structurée]
\textbf{Sujet :} En vous appuyant sur les documents 1 à 4 ci-après et vos connaissances, montrez que l'accession du Cameroun à la souveraineté internationale résulte d'une conjonction de luttes nationalistes et de négociations diplomatiques.

\textbf{Document 1 :} Extrait du discours de Ruben Um Nyobé à l'ONU (1952) -- cf. exercice précédent.
\textbf{Document 2 :} Extrait de la Loi-cadre Defferre (23 juin 1956) : « Il est institué dans chaque territoire d'Outre-mer une Assemblée territoriale élue au suffrage universel... Le Chef du gouvernement est responsable devant l'Assemblée. »
\textbf{Document 3 :} Résultats du référendum du 11 février 1961 au Cameroun britannique (voir tableau exercice 4).
\textbf{Document 4 :} Extrait de la Constitution de la République Fédérale du Cameroun (1er septembre 1961, Conférence de Foumban) : « La République Fédérale du Cameroun est constituée de deux États fédérés : l'État de l'Est (ex-Cameroun français) et l'État de l'Ouest (ex-Cameroun britannique méridional). »

\textbf{Travail demandé :}
\begin{enumerate}
\item Analyser le sujet (définir les termes, poser le problème, annoncer le plan).
\item Rédiger une introduction complète (accroche, définitions, problématique, annonce du plan).
\item Réaliser le développement structuré en deux grandes parties (I. Les luttes nationalistes : de la revendication à la confrontation armée ; II. Les négociations diplomatiques : de l'autonomie interne à la réunification fédérale), chaque partie comprenant deux paragraphes argumentés s'appuyant sur les documents et les connaissances.
\item Rédiger une conclusion synthétique et une ouverture sur la construction de l'État-nation (1961--1972).
\end{enumerate}
\end{exercice}

\begin{exercice}[Sujet type Probatoire : Commentaire de carte / Croquis]
\textbf{Sujet :} Le référendum du 11 février 1961 au Cameroun britannique : géographie du vote et conséquences territoriales.

\textbf{Travail demandé :}
\begin{enumerate}
\item Réaliser un croquis légendé et titré sur le fond de carte du Cameroun britannique fourni, faisant apparaître :
      \begin{itemize}
      \item La limite Nord/Sud (ligne pointillée).
      \item Le sens du vote par zone (hachures rouges Nord / vertes Sud).
      \item Les principales villes (Bamenda, Buea, Kumba, Garoua, Mamfe).
      \item La frontière internationale issue du vote (trait fort Nord/Nigeria, trait fort Sud/Rép. du Cameroun).
      \item Une flèche de rattachement du Sud vers Yaoundé.
      \end{itemize}
\item Rédiger un commentaire organisé (environ 20 lignes) structuré en trois parties :
      \begin{itemize}
      \item Le contexte et le déroulement du scrutin (question posée, rôle de l'ONU, participation).
      \item L'analyse géographique des résultats (opposition Nord/Sud, facteurs ethniques, religieux, économiques).
      \item Les conséquences territoriales et politiques immédiates (indépendance du Nord dans le Nigeria, réunification du Sud, Conférence de Foumban, fédéralisme).
      \end{itemize}
\end{enumerate}
\end{exercice}

\begin{exercice}[Sujet type Bac : Analyse de document + Mise en situation administrative]
\textbf{Partie A : Analyse de document (Lettre de Ruben Um Nyobé au Secrétaire général de l'ONU, 1954).}
\textbf{Document :} « Monsieur le Secrétaire Général, [...] Le peuple camerounais observe avec angoisse la transformation de la tutelle en colonisation déguisée. Le Code du travail indigène est maintenu, les libertés syndicales et politiques sont bafouées, l'administration française refuse tout dialogue avec l'UPC, seul parti de masse. Nous demandons l'envoi d'une mission de visite de l'ONU dotée de pouvoirs d'enquête réels et la fixation d'une date butoir pour l'indépendance. »

\begin{enumerate}
\item \textbf{Contexte :} Situer la date (1954) par rapport à la loi-cadre Defferre et à l'évolution de la situation au Cameroun français.
\item \textbf{Auteur :} Présenter Ruben Um Nyobé (origine, fonction, rôle à l'UPC).
\item \textbf{Thèses défendues :} Identifier les trois arguments principaux développés dans la lettre.
\item \textbf{Portée historique :} Expliquer pourquoi cette lettre marque une étape dans l'internationalisation du conflit.
\item \textbf{Limites du document :} Préciser ce que ce document ne dit pas sur la stratégie de l'UPC (lutte armée) et sur la position de la France.
\end{enumerate}

\textbf{Partie B : Mise en situation administrative (Savoir-faire Unité).}
\textbf{Consigne :} Vous êtes conseiller municipal dans une commune de l'Ouest (ex: Bafoussam) en 2024. La commune accueille d'importantes populations déplacées internes (PDI) venant des régions du Nord-Ouest et du Sud-Ouest (crise anglophone depuis 2016). Des tensions surgissent entre populations « autochtones » et « allogènes » (accès au foncier, commerces, services publics).
Proposer un \textbf{Plan d'action concret (3 mesures prioritaires)} pour renforcer le vivre-ensemble, en vous inspirant explicitement de l'esprit de la Réunification de 1961 (dialogue, fédéralisme originel, respect des particularismes, unité dans la diversité).
\textbf{Format :} Note administrative structurée (Objet, Contexte, 3 Mesures chiffrées/budgétisées si possible, Indicateurs de suivi, Signature).
\end{exercice}



\newpage

\coursec{Corrigés des exercices}

\coursec{Exercices corrigés — La marche vers l’indépendance (1946-1961)}

\begin{corrige}[Exercice 1 — QCM : La Tutelle et les institutions]
\begin{enumerate}
\item \textbf{1946}. Les accords de tutelle sont signés le 13 décembre 1946 (approuvés par l'Assemblée générale de l'ONU en 1947) : le Cameroun passe du mandat de la SDN à la tutelle de l'ONU.
\item \textbf{France}. Le Cameroun français (ex-mandat français de 1922) est administré par la France ; le Cameroun britannique l'est par la Grande-Bretagne.
\item \textbf{Defferre}. La loi-cadre du 23 juin 1956, portée par Gaston Defferre, ministre de la France d'Outre-mer, instaure le suffrage universel et l'autonomie interne.
\item \textbf{1948}. L'UPC est fondée en avril 1948 à Douala, avec Ruben Um Nyobé comme Secrétaire général.
\end{enumerate}
\textbf{Points de vigilance :} ne pas confondre la date de signature des accords de tutelle (1946) avec la date de la loi-cadre (1956) ; distinguer clairement les deux puissances administrantes (France à l'Est, Grande-Bretagne à l'Ouest).
\end{corrige}

\begin{corrige}[Exercice 2 — Vrai / Faux justifié : Le nationalisme et l'UPC]
\begin{enumerate}
\item \textbf{Vrai.} Les statuts de l'UPC (avril 1948) fixent deux objectifs : l'indépendance immédiate du Cameroun et la réunification des deux territoires sous tutelle (français et britannique).
\item \textbf{Vrai.} Ruben Um Nyobé se rend à New York et intervient devant la IV\textsuperscript{e} Commission de l'ONU en décembre 1952 puis en 1954 pour dénoncer la violation des accords de tutelle par la France.
\item \textbf{Faux.} L'ordre chronologique est inversé : l'insurrection éclate en mai 1955 (Sanaga-Maritime), \emph{avant} l'interdiction officielle de l'UPC prononcée le 13 juillet 1955 par le Haut-commissaire Roland Pré. L'interdiction est une conséquence des troubles, non leur cause.
\item \textbf{Vrai.} Après l'assassinat de Ruben Um Nyobé le 13 septembre 1958 dans la Sanaga-Maritime, Félix-Roland Moumié prend la direction de l'UPC en exil (Conakry, Accra) jusqu'à son propre assassinat le 3 novembre 1960.
\end{enumerate}
\textbf{Points de vigilance :} la justification doit toujours contenir un repère précis (date, acteur, lieu). Pour l'affirmation 3, c'est la maîtrise de la chronologie mai 1955 / juillet 1955 qui fait la différence.
\end{corrige}

\begin{corrige}[Exercice 3 — Définitions et repères chronologiques]
\begin{enumerate}
\item \cle{Tutelle de l'ONU} : statut international institué en 1946 (accords signés le 13 décembre 1946, approuvés en 1947) qui remplace le mandat de la SDN de 1922. La puissance administrante (France, Grande-Bretagne) a l'\emph{obligation} de conduire le territoire vers l'autonomie ou l'indépendance, sous le contrôle de l'ONU (rapports annuels, missions de visite, examen des pétitions). Le mandat de la SDN, lui, n'imposait qu'un rapport annuel, sans perspective d'indépendance ni contrôle effectif.
\item \cle{Loi-cadre Defferre} : loi française du 23 juin 1956 (Gaston Defferre, ministre de la France d'Outre-mer). Elle instaure le suffrage universel, une Assemblée territoriale élue et un Conseil de gouvernement responsable devant elle : c'est l'autonomie interne, étape vers l'indépendance. Au Cameroun, elle permet les élections du 23 décembre 1956 et l'avènement d'André-Marie Mbida puis d'Ahmadou Ahidjo.
\item \cle{« Guerre cachée » ou « maquis »} : conflit armé de 1955 à 1971 entre l'UPC (maquisards) et les forces françaises puis camerounaises, surtout en Sanaga-Maritime, au pays Bassa, au pays Bamiléké et dans le Mungo. Répression massive (regroupements de populations, bombardements), élimination des leaders : Um Nyobé (1958), Moumié (1960), Ouandié (exécuté en 1971).
\item \cle{Référendum du 11 février 1961} : scrutin organisé au Cameroun britannique sous supervision de l'ONU. La question porte sur le mode d'accession à l'indépendance : rejoindre le Nigeria ou rejoindre la République du Cameroun. Le Nord choisit le Nigeria (environ 60 \%), le Sud la réunification (environ 79 \%). Conséquences : rattachement du Nord au Nigeria le 1\textsuperscript{er} juin 1961 ; indépendance du Sud par réunification le 1\textsuperscript{er} octobre 1961.
\end{enumerate}
\textbf{Points de vigilance :} une définition au Bac doit toujours comporter trois éléments : la nature, la date, le contexte. Éviter les définitions circulaires (« la tutelle est le fait d'être sous tutelle »).
\end{corrige}

\begin{corrige}[Exercice 4 — Calculs directs et données chiffrées]
\begin{enumerate}
\item Taux de participation aux élections de 1956 :
\[ \text{Taux} = \frac{\text{Votants}}{\text{Inscrits}} \times 100 = \frac{720\,000}{1\,200\,000} \times 100 = 60 \%. \]
\item Nombre de voix pour la réunification au Sud :
\[ 79\,\% \times 230\,000 = 0,79 \times 230\,000 = \textbf{181\,700 voix}. \]
Vérification : les voix pour l'intégration au Nigeria au Sud valent $0,21 \times 230\,000 = 48\,300$ ; or $181\,700 + 48\,300 = 230\,000$ : le total des votants est bien retrouvé.
\item Pourcentage de voix pour l'intégration au Nigeria au Nord : \textbf{60 \%} (lecture directe du tableau), soit $0,60 \times 170\,000 = 102\,000$ voix.
\item Choix majoritaires et conséquences :
\begin{itemize}
\item \textbf{Nord camerounais} : majorité pour l'intégration au Nigeria (60 \%) $\rightarrow$ rattachement effectif au Nigeria le 1\textsuperscript{er} juin 1961.
\item \textbf{Sud camerounais} : majorité pour la réunification (79 \%) $\rightarrow$ indépendance par rattachement à la République du Cameroun le 1\textsuperscript{er} octobre 1961, donnant naissance à la République fédérale du Cameroun.
\end{itemize}
\end{enumerate}
\textbf{Points de vigilance :} bien appliquer le pourcentage aux \emph{votants} et non aux inscrits ; ne pas oublier de conclure par la conséquence territoriale datée, c'est elle qui est historiquement significative.
\end{corrige}

\begin{corrige}[Exercice 5 — Analyse de document : Discours d'Um Nyobé à l'ONU (1952)]
\begin{enumerate}
\item \textbf{Présentation du document :} auteur : Ruben Um Nyobé, Secrétaire général de l'UPC ; nature : pétition (discours officiel) ; date : décembre 1952 ; destinataire : la IV\textsuperscript{e} Commission de l'Assemblée générale de l'ONU, chargée des questions de tutelle, à New York.
\item \textbf{Deux griefs contre la puissance administrante :} (a) la France a violé les termes de l'accord de tutelle en refusant d'associer les véritables représentants du peuple à l'administration ; (b) la France maintient ses troupes sur le territoire et empêche l'expression libre du peuple (absence d'élections libres).
\item \textbf{« Élections libres sous contrôle de l'ONU » :} l'UPC demande que le scrutin soit organisé et surveillé par l'ONU, et non par l'administration française, afin de garantir la sincérité du vote et de permettre l'élection d'une Assemblée constituante véritablement représentative, chargée de préparer l'indépendance.
\item \textbf{Portée historique :} cette intervention internationalise la question camerounaise, fait de l'UPC l'interlocuteur reconnu du peuple devant l'ONU et contraint la France à justifier sa gestion devant les missions de visite ; elle marque une étape décisive de la lutte nationaliste pacifique avant la radicalisation de 1955.
\end{enumerate}
\textbf{Points de vigilance :} la présentation du document (question 1) exige quatre éléments : auteur, nature, date, destinataire. Pour la portée, rester concis et conclure sur l'effet historique, sans paraphraser le texte.
\end{corrige}

\begin{corrige}[Exercice 6 — Cartographie : Le référendum de 1961]
\textbf{Corrigé du croquis attendu} (description précise, à reproduire sur le fond de carte) :
\begin{itemize}
\item \textbf{Titre} (en haut, en capitales) : « Le référendum du 11 février 1961 au Cameroun britannique ».
\item \textbf{Limite Nord/Sud} : ligne pointillée horizontale au niveau de la cuvette de la Bénoué, séparant le Nord camerounais (région de Garoua) du Sud camerounais (régions de Bamenda, Buea).
\item \textbf{Hachures rouges} sur le Nord : zone ayant voté majoritairement « Intégration au Nigeria » (60 \%).
\item \textbf{Hachures vertes} sur le Sud : zone ayant voté majoritairement « Réunification avec le Cameroun » (79 \%).
\item \textbf{Villes} (points nommés) : Garoua (Nord), Bamenda, Buea, Kumba, Mamfe (Sud).
\item \textbf{Flèche} partant du Sud (Buea) vers Yaoundé, légendée « Réunification, 1\textsuperscript{er} octobre 1961 » ; flèche du Nord vers l'ouest (Nigeria), légendée « Rattachement au Nigeria, 1\textsuperscript{er} juin 1961 ».
\item \textbf{Éléments obligatoires} : flèche du nord (orientation), échelle indicative (par exemple 1 cm pour 100 km), légende organisée et numérotée dans un cartouche.
\end{itemize}

\begin{popfigure}
\centering
\begin{tikzpicture}[scale=0.8]
% Contour très simplifié du Cameroun britannique (forme trapézoïdale)
\draw[popdark, thick] (0,0) -- (8,0) -- (9,5) -- (7,7) -- (1,7) -- cycle;
% Limite Nord/Sud (cuvette de la Bénoué ~ y=3.5)
\draw[popdark, dashed, thick] (0.5,3.5) -- (8.5,3.5) node[midway, above, sloped, font=\small] {Limite Nord/Sud};
% Hachures Nord (Rouge)
\foreach \x in {0.7,1.0,...,8.3} \draw[poporange!70!popdark, thick] (\x,0.2) -- (\x+0.3,3.3);
% Hachures Sud (Vert)
\foreach \x in {0.7,1.0,...,8.3} \draw[popgreen!70!popdark, thick] (\x,3.7) -- (\x+0.3,6.7);
% Villes
\fill[popdark] (4.5,1.5) circle (2pt) node[below] {Garoua};
\fill[popdark] (2.5,5.0) circle (2pt) node[left] {Bamenda};
\fill[popdark] (5.5,5.5) circle (2pt) node[right] {Buea};
\fill[popdark] (6.0,4.5) circle (2pt) node[right] {Kumba};
\fill[popdark] (3.5,6.0) circle (2pt) node[above] {Mamfe};
% Yaoundé (hors carte, à l'Est)
\fill[popdark] (9.5,3.5) circle (2pt) node[right] {Yaoundé};
% Flèches
\draw[popblue, thick, ->] (5.5,5.5) -- (9.0,3.8) node[midway, above, sloped, font=\small, popblue] {Réunification 1er oct. 1961};
\draw[poporange, thick, ->] (4.5,1.5) -- (0.5,1.0) node[midway, below, sloped, font=\small, poporange] {Rattachement Nigeria 1er juin 1961};
% Nord (flèche orientation)
\draw[popdark, ->, thick] (0.5,6.5) -- (0.5,7.2) node[above] {N};
% Échelle
\draw[popdark, thick] (6,0.2) -- (7,0.2) node[midway, below] {100 km};
% Légende (simplifiée dans le code, à faire à la main sur la copie)
\node[draw, popdark, fill=popcream, rounded corners, font=\small, anchor=south east, align=center] at (9.5, -0.5){
\begin{tabular}{ll}
\textcolor{poporange!70!popdark}{\rule{1em}{1ex}} & Intégration au Nigeria (Nord, 60\%) \\
\textcolor{popgreen!70!popdark}{\rule{1em}{1ex}} & Réunification (Sud, 79\%) \\
-- -- -- & Limite Nord/Sud \\
$\bullet$ & Villes principales
\end{tabular}
};
\end{tikzpicture}
\legende{Croquis schématique du référendum du 11 février 1961 : résultats par zone, villes principales, flux territoriaux et éléments cartographiques obligatoires.}
\end{popfigure}

\textbf{Barème indicatif (sur 10) :} titre 1 pt ; orientation 1 pt ; échelle 1 pt ; limite Nord/Sud 1 pt ; hachures correctes 2 pts ; villes 2 pts ; flèches datées 1 pt ; légende organisée 1 pt.
\textbf{Points de vigilance :} une légende non organisée (symboles dispersés dans le cartouche) est pénalisée ; ne pas confondre les dates de rattachement du Nord (juin) et du Sud (octobre).
\end{corrige}

\begin{corrige}[Exercice 7 — Tableau comparatif : Les grandes figures du nationalisme]
Le tableau complété est le suivant :
\begin{center}
\begin{tabularx}{\linewidth}{|l|X|X|X|X|}
\hline
\rowcolor{popblueL}
\textbf{Figure} & \textbf{Formation / Débuts} & \textbf{Parti dirigé} & \textbf{Revendications} & \textbf{Fin de parcours} \\ \hline
Ruben Um Nyobé & Instituteur (École normale de Foulassi) ; syndicaliste CGT & UPC (fondateur, SG 1948) & Indépendance immédiate, réunification, réformes sociales & Assassiné par l'armée française le 13 septembre 1958 (Sanaga-Maritime) \\ \hline
Félix-Roland Moumié & Médecin (Dakar) ; militant UPC & UPC (SG en exil, 1958--1960) & Lutte armée, internationalisation de la cause, réunification & Empoisonné au thallium par les services secrets français, 3 novembre 1960, Genève \\ \hline
Ernest Ouandié & Instituteur ; vice-président de l'UPC & UPC (chef militaire du maquis, 1958--1970) & Indépendance réelle, réunification, justice sociale & Capturé en août 1970, exécuté le 15 janvier 1971 à Bafoussam \\ \hline
Ahmadou Ahidjo & École normale de Yaoundé ; fonctionnaire, député & UC (1958), puis UNC (1966) & Autonomie interne (1956), indépendance négociée (1960), unité nationale & 1\textsuperscript{er} Président de la République (1960--1982) ; mort en exil à Dakar en 1989 \\ \hline
John Ngu Foncha & Instituteur, études au Nigeria ; syndicaliste & KNDP (1955), Premier ministre du Cameroun du Sud (1959) & Réunification par le fédéralisme, autonomie régionale & Vice-président fédéral (1961--1966) ; mort en 2010 \\ \hline
\end{tabularx}
\end{center}
\textbf{Points de vigilance :} bien distinguer les deux voies nationalistes : la voie radicale (UPC : Um Nyobé, Moumié, Ouandié) et la voie négociée (Ahidjo, Foncha). Les dates de mort des leaders UPC sont des repères exigibles.
\end{corrige}

\begin{corrige}[Exercice 8 — Savoir-faire : Argumentation historique]
\textbf{Paragraphe 1 — L'indépendance du Cameroun français (1\textsuperscript{er} janvier 1960) : une victoire de la voie négociée.} L'indépendance de 1960 résulte d'abord du processus institutionnel ouvert par la loi-cadre Defferre du 23 juin 1956 : suffrage universel, Assemblée territoriale, autonomie interne. Ahmadou Ahidjo, à la tête de l'Union Camerounaise, remporte les élections (l'UPC, interdite depuis juillet 1955, appelle au boycott), devient Premier ministre en février 1958 et négocie avec la France le transfert des compétences. L'ONU entérine la fin de la tutelle (résolution de décembre 1959), et l'indépendance est proclamée le 1\textsuperscript{er} janvier 1960. La lutte armée de l'UPC (maquis de 1955, assassinat d'Um Nyobé en 1958) a certes pesé sur la décision française, mais elle est \emph{exclue} du processus : l'indépendance se fait sans elle et contre elle.

\textbf{Paragraphe 2 — La réunification (1961) : un acte diplomatique distinct de l'insurrection.} La réunification concerne le Cameroun britannique. Le référendum du 11 février 1961, organisé sous contrôle de l'ONU, voit le Sud choisir la réunification (79 \%) sous l'impulsion de John Ngu Foncha (KNDP), tandis que le Nord rejoint le Nigeria. La conférence constitutionnelle de Foumban (juillet 1961) entre Ahidjo et Foncha fixe le cadre fédéral, et le 1\textsuperscript{er} octobre 1961 naît la République fédérale du Cameroun. L'UPC (Moumié assassiné en 1960, Ouandié dans le maquis) rejette ce compromis qu'elle qualifie de « néocolonial » et poursuit la lutte armée jusqu'en 1971. La souveraineté camerounaise est donc bien le produit d'une \emph{conjonction} : la pression nationaliste a rendu le statu quo impossible, mais ce sont les négociations (Ahidjo--France, Ahidjo--Foncha, arbitrage de l'ONU) qui ont fixé les formes et les dates de l'indépendance et de la réunification.

\textbf{Points de vigilance :} exigence majeure du sujet : \emph{nuancer}, c'est-à-dire ne pas attribuer l'indépendance à la seule UPC ni à la seule négociation ; distinguer rigoureusement les deux dates (1960 / 1961) et les deux processus.
\end{corrige}

\begin{corrige}[Exercice 9 — Sujet type Baccalauréat : Dissertation structurée]
\textbf{1. Analyse du sujet.} Termes clés : « souveraineté internationale » (reconnaissance de l'État camerounais comme sujet de droit international : indépendance du 1\textsuperscript{er} janvier 1960, réunification du 1\textsuperscript{er} octobre 1961) ; « luttes nationalistes » (pétitions, syndicalisme, UPC, maquis) ; « négociations diplomatiques » (loi-cadre, accords franco-camerounais, référendum ONU, Foumban). Problème : l'indépendance du Cameroun est-elle le fruit des armes ou de la table des négociations ? Plan en deux parties imposé par le libellé.

\textbf{2. Introduction (modèle).} Le 1\textsuperscript{er} janvier 1960, le Cameroun français accède à l'indépendance ; le 1\textsuperscript{er} octobre 1961, le Sud du Cameroun britannique le rejoint pour former la République fédérale du Cameroun. Cette double souveraineté est l'aboutissement de quinze années de mobilisation. D'un côté, les luttes nationalistes : pétitions de l'UPC à l'ONU dès 1948 (document 1), insurrection et maquis à partir de 1955. De l'autre, les négociations diplomatiques : loi-cadre Defferre de 1956 (document 2), référendum de 1961 (document 3), conférence de Foumban (document 4). \textbf{Problématique :} comment la pression nationaliste et le compromis diplomatique se sont-ils articulés pour aboutir à la souveraineté du Cameroun ? \textbf{Annonce du plan :} nous étudierons d'abord les luttes nationalistes, de la revendication pacifique à la confrontation armée (I), puis les négociations diplomatiques, de l'autonomie interne à la réunification fédérale (II).

\textbf{3. Développement (plan détaillé attendu).}
\textbf{I. Les luttes nationalistes : de la revendication à la confrontation armée.}
\begin{itemize}
\item \textbf{I.1. La revendication pacifique et internationalisée (1945--1955) :} essor syndical (USCC, grèves de 1945 et 1953) ; création de l'UPC (avril 1948) ; pétitions et interventions d'Um Nyobé à l'ONU (1952, 1954 — document 1) : indépendance immédiate et réunification ; l'UPC devient le parti de masse.
\item \textbf{I.2. La confrontation armée (1955--1971) :} insurrection de mai 1955, interdiction de l'UPC (13 juillet 1955), maquis en Sanaga-Maritime, au pays Bassa et Bamiléké ; répression (« guerre cachée ») ; élimination des leaders (Um Nyobé 1958, Moumié 1960, Ouandié 1971). Bilan : la lutte armée échoue militairement mais déstabilise durablement l'ordre colonial.
\end{itemize}
\textbf{II. Les négociations diplomatiques : de l'autonomie interne à la réunification fédérale.}
\begin{itemize}
\item \textbf{II.1. De la loi-cadre à l'indépendance (1956--1960) :} la loi-cadre Defferre (document 2) instaure suffrage universel et autonomie interne ; élections de 1956 et 1957 ; Ahidjo Premier ministre (1958) ; négociations franco-camerounaises ; résolution de l'ONU (décembre 1959) ; indépendance le 1\textsuperscript{er} janvier 1960.
\item \textbf{II.2. Le référendum et la réunification (1961) :} référendum du 11 février 1961 sous contrôle de l'ONU (document 3) : le Nord choisit le Nigeria, le Sud la réunification (79 \%) ; conférence de Foumban (juillet 1961) et Constitution fédérale (document 4) : deux États fédérés ; naissance de la République fédérale le 1\textsuperscript{er} octobre 1961.
\end{itemize}

\textbf{4. Conclusion (modèle).} La souveraineté du Cameroun n'est ni le seul fruit des armes ni le seul produit des négociations : la lutte nationaliste a rendu le maintien de la tutelle intenable, tandis que la diplomatie (Ahidjo, Foncha, ONU, France) a fixé les formes et le calendrier de l'indépendance et de la réunification. Mais cette indépendance négociée sans l'UPC laisse un héritage conflictuel : la répression du maquis se poursuit jusqu'en 1971 et la mémoire des nationalistes reste longtemps occultée. \textbf{Ouverture :} la construction de l'État-nation se poursuit avec la République unie de 1972, dont les équilibres sont encore questionnés aujourd'hui.

\textbf{Barème indicatif (sur 20) :} présentation et introduction 4 pts ; analyse et mobilisation des documents 6 pts ; connaissances personnelles précises (dates, acteurs) 4 pts ; construction du développement en deux parties 4 pts ; conclusion et ouverture 2 pts.
\textbf{Points de vigilance :} citer explicitement chaque document dans le développement ; ne pas faire de I.2 un récit interminable du maquis au détriment de la partie II ; la problématique doit être une question, pas une affirmation.
\end{corrige}

\begin{corrige}[Exercice 10 — Sujet type Probatoire : Commentaire de carte / Croquis]
\textbf{1. Croquis attendu} : se reporter au croquis TikZ fourni à l'exercice 6 (mêmes exigences : titre, limite Nord/Sud, hachures, villes, flèches datées, échelle, nord, légende).

\textbf{2. Commentaire organisé (modèle).}
\textbf{Contexte et déroulement.} Le Cameroun britannique, sous tutelle de l'ONU depuis 1946, est administré par la Grande-Bretagne en union avec le Nigeria. Lorsque celui-ci devient indépendant (1\textsuperscript{er} octobre 1960), l'ONU organise un référendum le 11 février 1961 : au Nord, la question est de rejoindre le Nigeria ou la République du Cameroun ; au Sud, la même alternative se pose. La participation est forte et le scrutin se déroule sous supervision internationale, garant de sa régularité.
\textbf{Analyse géographique des résultats.} Le vote dessine deux ensembles opposés. Au Nord, environ 60 \% des suffrages se portent sur l'intégration au Nigeria : cette zone, majoritairement musulmane et peule, entretient des liens religieux, politiques et commerciaux anciens avec les émirats du Nord nigérian. Au Sud, environ 79 \% des votants choisissent la réunification : populations en partie christianisées, élites anglophones hostiles à la domination du Nigeria, liens historiques et linguistiques avec le Cameroun de l'Est, et rôle moteur du KNDP de John Ngu Foncha. Le facteur ethnique et religieux apparaît donc déterminant dans la géographie du vote.
\textbf{Conséquences territoriales et politiques.} Le Nord camerounais est rattaché au Nigeria le 1\textsuperscript{er} juin 1961 (région du Sardauna). Le Sud accède à l'indépendance par réunification avec la République du Cameroun le 1\textsuperscript{er} octobre 1961, après la conférence de Foumban (juillet 1961) qui adopte une Constitution fédérale : deux États fédérés (Est et Ouest), un président (Ahidjo), un vice-président (Foncha). La réunification réalise ainsi l'un des deux objectifs historiques de l'UPC, mais selon des modalités fédérales que celle-ci conteste.

\textbf{Barème indicatif (sur 20) :} croquis 10 pts (titre 1, orientation 1, échelle 1, limite Nord/Sud 1, hachures 2, villes 2, flèches 1, légende 1) ; commentaire 10 pts (contexte 3, analyse géographique 4, conséquences 3).
\textbf{Points de vigilance :} le commentaire doit suivre le plan en trois parties annoncé ; chaque affirmation géographique doit être appuyée par un chiffre du scrutin ; ne pas confondre « indépendance » (Sud, par réunification) et « rattachement » (Nord, sans indépendance propre).
\end{corrige}

\begin{corrige}[Exercice 11 — Sujet type Bac : Analyse de document + Mise en situation administrative]
\textbf{Partie A — Analyse de la lettre d'Um Nyobé (1954).}
\begin{enumerate}
\item \textbf{Contexte :} la lettre date de 1954, soit deux ans \emph{avant} la loi-cadre Defferre (23 juin 1956). Le Cameroun français vit alors sous le statu quo colonial : le Code du travail indigène est encore en vigueur, les grèves syndicales de 1953 sont réprimées, et l'UPC, parti de masse légal, est marginalisée par l'administration qui favorise des formations modérées. C'est aussi l'année de la deuxième intervention d'Um Nyobé à l'ONU.
\item \textbf{Auteur :} Ruben Um Nyobé, né en 1913 au pays Bassa, instituteur formé à l'École normale de Foulassi, syndicaliste CGT, Secrétaire général fondateur de l'UPC (1948). Surnommé « Mpodol » (le porte-parole), il porte la cause camerounaise devant l'ONU à trois reprises (1952, 1954, 1955).
\item \textbf{Thèses défendues :} (a) la tutelle est détournée en « colonisation déguisée » (maintien du Code indigène, libertés bafouées) ; (b) l'administration française refuse le dialogue avec l'UPC, seul parti de masse, ce qui délégitime sa gestion ; (c) il faut une mission de visite de l'ONU dotée de pouvoirs d'enquête réels et la fixation d'une date butoir pour l'indépendance.
\item \textbf{Portée historique :} la lettre saisit directement le Secrétaire général de l'ONU : elle internationalise le conflit, contribue à l'envoi de la mission de visite de 1955 et accroît la pression sur la France, qui répondra par le durcissement (interdiction de l'UPC le 13 juillet 1955).
\item \textbf{Limites du document :} texte de propagande diplomatique, il ne dit rien de la préparation de la lutte armée (le maquis de 1955), ni de la répression déjà en cours, ni des débats internes à l'UPC entre partisans de la voie légale et de la radicalisation. Il présente aussi la position française uniquement de manière dénonciatrice.
\end{enumerate}

\textbf{Partie B — Note administrative (modèle attendu).}
\begin{quote}
\textbf{NOTE ADMINISTRATIVE --- Commune de Bafoussam, 2024}\\
\textbf{OBJET :} Plan d'action « Vivre-Ensemble » pour l'intégration des personnes déplacées internes (PDI) du Nord-Ouest et du Sud-Ouest.\\
\textbf{CONTEXTE :} depuis 2016, la crise anglophone provoque l'arrivée de nombreuses PDI. Des tensions (foncier, commerces, services publics) opposent populations autochtones et allogènes. En cohérence avec l'esprit de la Réunification de 1961 (dialogue de Foumban, respect des particularismes, unité dans la diversité), la commune propose trois mesures.\\
\textbf{MESURE 1 --- Dialogue institutionnalisé :} création d'un Conseil communal de la réconciliation, paritaire (autochtones, PDI, chefs traditionnels et religieux). Budget : 2\,000\,000 FCFA/an. Indicateur : 4 sessions par an, 80 \% des conflits réglés à l'amiable.\\
\textbf{MESURE 2 --- Accès équitable aux ressources :} programme « Foncier solidaire » : mise à disposition de parcelles communales en bail de 10 ans pour les PDI, avec appui agricole (semences, microcrédits). Budget : 15\,000\,000 FCFA (partenariat MINADER/ONG). Indicateur : 200 ha attribués par an.\\
\textbf{MESURE 3 --- Éducation à l'unité dans la diversité :} modules « Histoire partagée » (Réunification, fédéralisme de 1961) dans les écoles de la commune ; journée annuelle du vivre-ensemble le 11 février. Budget : 3\,000\,000 FCFA. Indicateur : 100 \% des écoles couvertes.\\
\textbf{SUIVI :} tableau de bord trimestriel (Maire, Préfet, MINAT). \textbf{SIGNATURE :} Le Maire.
\end{quote}
\textbf{Barème indicatif (sur 20) :} Partie A : 12 pts (contexte 3, auteur 2, thèses 3, portée 2, limites 2). Partie B : 8 pts (format de note 2, trois mesures concrètes et chiffrées 4, lien explicite avec 1961 2).
\textbf{Points de vigilance :} en Partie A, la question des limites est souvent négligée : tout document a un point de vue, il faut dire ce qu'il tait. En Partie B, le lien avec l'esprit de 1961 (dialogue, fédéralisme originel, unité dans la diversité) doit être \emph{explicite}, c'est le cœur du savoir-faire évalué.
\end{corrige}

\newpage

\coursec{Fiche bilan}

\begin{popfigure}
\begin{tikzpicture}[
    mindmap,
    concept color=popblue,
    level 1/.style={level distance=4.5cm, sibling angle=360/5},
    level 2/.style={level distance=3cm, sibling angle=360/6},
    every node/.style={concept, font=\small\bfseries, text=white},
    grow cyclic,
    text width=2.8cm,
    align=center
]
\node[concept color=popdark, align=center]{Marche vers\\ l'Indépendance}
    child[concept color=popblue] { node[align=center]{1. Tutelle ONU\\ (1946--1960)}
        child { node {Statut fiduciaire} }
        child { node {Double admin. FR/UK} }
        child { node {Rapports ONU} }
    }
    child[concept color=poporange] { node[align=center]{2. Nationalisme\\ (1945--1955)}
        child { node {Syndicats (USCC)} }
        child { node {Partis (UPC, ANC)} }
        child { node {Revendications} }
    }
    child[concept color=popgreen] { node[align=center]{3. Figures\\ emblématiques}
        child { node {R. Moumié} }
        child { node {U. Ouandié} }
        child { node {F. Eboué (héritage)} }
        child { node {A. Ahidjo} }
        child { node {J. Foncha} }
    }
    child[concept color=poppurple] { node[align=center]{4. Vers la\\ Souveraineté}
        child { node {Insurrection 1955} }
        child { node {Loi-cadre 1956} }
        child { node {Référendum 1960} }
        child { node {Indép. 1er janv 1960} }
        child { node {Réunif. 1961} }
    }
    child[concept color=poppink] { node[align=center]{5. Construction\\ État-Nation}
        child { node {Ahmadou Ahidjo} }
        child { node {Référendum 1972} }
        child { node {Rép. Unie} }
        child { node {Mémoire/Justice} }
    };
\end{tikzpicture}
\legende{Carte mentale de synthèse : La marche vers l'indépendance du Cameroun (1946--1972).}
\end{popfigure}

\begin{bilan}

\courssub{La Tutelle de l'ONU (1946--1960) : Double héritage}

\begin{definition}[Tutelle / Territoire sous tutelle]
Statut juridique international (Chap. XII Charte ONU) succédant au mandat SDN. L'administrateur (France ou Royaume-Uni) doit préparer l'autonomie politique, économique, sociale et éducative du territoire sous contrôle de l'ONU (Conseil de tutelle, missions de visite, rapports annuels).
\end{definition}

\begin{itemize}
    \item \textbf{Cameroun français :} Loi du 13 déc. 1946 $\rightarrow$ Territoire d'Outre-Mer (TOM). Assemblée territoriale (ARTCAM, puis ATCAM). Évolution : Loi-cadre Defferre (1956) $\rightarrow$ Autonomie interne (1957) $\rightarrow$ Gouvernement Ahidjo (1958).
    \item \textbf{Cameroun britannique :} Administré comme partie du Nigeria (Lagos). Deux zones : Nord (musulman, émirats) et Sud (bantou, élites éduquées). Représentation à l'Assemblée législative nigériane puis fédérale.
    \item \textbf{Contrôle ONU :} Missions de visite (1949, 1952, 1955, 1958), rapports critiques sur répression UPC, lenteur réformes, question réunification.
\end{itemize}

\courssub{Éveil et structuration du nationalisme (1945--1955)}

\begin{propriete}[Facteurs du réveil nationaliste]
\begin{itemize}
    \item \textbf{Contexte international :} Charte ONU (droit des peuples), Guerre froide, décolonisation asiatique (Inde 1947), Conférence de Bandung (1955).
    \item \textbf{Facteurs endogènes :} Exploitation coloniale (travail forcé, code de l'indigénat), retour des tirailleurs (conscience politique), essor syndical (USCC, CGT), presse (\emph{L'Écho camerounais}, \emph{La Voix du Cameroun}).
    \item \textbf{Acteurs clés :} Syndicats (Um Nyobé, Moumié), Élites traditionnelles (chefs Bamileke, Bassa), Évolués (fonctionnaires, enseignants).
\end{itemize}
\end{propriete}

\begin{itemize}
    \item \textbf{Partis majeurs :}
    \begin{itemize}
        \item \cle{UPC} (Union des Populations du Cameroun, 1948) : Um Nyobé, Moumié, Ouandié. Idéologie : nationalisme radical, réunification, indépendance immédiate, marxisme. Interdite 1955 $\rightarrow$ maquis.
        \item \cle{ANC} (Action Nationale Camerounaise, 1955) : Ahmadou Ahidjo. Modérés, légalistes, fédéralisme, coopération France.
        \item \cle{KNDP} (Kamerun National Democratic Party, 1955) : John Foncha (Sud britannique). Réunification avec Cameroun français.
        \item \cle{ONE} (Organisation Nord-Étudiants) / \cle{CU} (Cameroon United) : Nord britannique, pro-intégration Nigeria.
    \end{itemize}
\end{itemize}

\courssub{Dossier 3 : Grandes figures (Portraits croisés)}

\begin{center}
\begin{tabularx}{\linewidth}{|l|X|X|X|}
\hline
\rowcolor{popblueL}
\textbf{Figure} & \textbf{Rôle / Trajectoire} & \textbf{Idées forces} & \textbf{Sort / Héritage} \\
\hline
\textbf{Ruben Um Nyobé} (1913--1958) & Secrétaire général UPC, « \cle{Mpodol} ». Député ATCAM. Maquis 1955. Tué par armée française. & Indépendance immédiate, réunification, justice sociale, non-violence initiale. & Martyr national. Réhabilité 1991. Symbole unité. \\
\hline
\textbf{Félix-Roland Moumié} (1925--1960) & Président UPC (après 1958). Médecin. Exil (Ghana, Guinée, Maroc). Empoisonné (Genève) par SDECE. & Internationalisation lutte (ONU, OUA, Bloc de l'Est), lutte armée. & Figure panafricaine. Assassinat impuni. \\
\hline
\textbf{Ernest Ouandié} (1924--1971) & Vice-président UPC. Chef militaire maquis (Bamileke). Arrêté 1970. Exécuté public 1971. & Fidélité à Nyobé/Moumié, résistance armée jusqu'au bout. & Dernier grand chef maquis. Symbole répression post-indép. \\
\hline
\textbf{Ahmadou Ahidjo} (1924--1989) & Député ANC, VP Assemblée, PM 1958, 1er Président (1960--1982). & Pragmatisme, unité nationale, fédéralisme puis unitarisme, alliance France. & Architecte État moderne. Démission 1982. \\
\hline
\textbf{John Ngu Foncha} (1916--1999) & Premier ministre Sud-Britannique (1959), VP République Fédérale (1961--1970). & Réunification, fédéralisme, droits minorités anglophones. & Démission 1970 (conteste unitarisme). Père anglophone. \\
\hline
\end{tabularx}
\end{center}

\courssub{De la revendication à la souveraineté : Étapes clés (1955--1961)}

\begin{methode}[Chronologie décisive]
\begin{enumerate}
    \item \textbf{Mai 1955 :} Insurrection UPC (Douala, Bassa, Bamileke). Répression sanglante (dizaines/centaines de morts). Interdiction UPC (13 juil.). Début guerre de libération (maquis).
    \item \textbf{23 juin 1956 :} \cle{Loi-cadre Defferre}. Autonomie interne, suffrage universel, conseil de gouvernement.
    \item \textbf{1957 :} 1er gouvernement africain (André-Marie Mbida, puis Ahidjo fév. 1958). Demande indépendance + réunification à l'ONU.
    \item \textbf{12 juin 1958 :} Ahidjo à l'ONU : « Indépendance le 1er janv 1960 ».
    \item \textbf{13 mars 1959 :} Résolution ONU 1349 (XIV) : Fin tutelle 1er janv 1960.
    \item \textbf{1er janv 1960 :} \cle{Indépendance} République du Cameroun (Ahidjo Président).
    \item \textbf{11 fév 1961 :} \cle{Référendum} ONU Cameroun britannique. Nord $\rightarrow$ Nigeria (60\%). Sud $\rightarrow$ Réunification (70\%).
    \item \textbf{1er oct 1961 :} \cle{Réunification}. République Fédérale du Cameroun (2 États : Est / Ouest). Foumban (juil. 1961) : Constitution fédérale.
\end{enumerate}
\end{methode}

\courssub{Construire l'État-Nation : Unité, Réconciliation, Mémoire (1961--Aujourd'hui)}

\begin{attention}[Pièges historiographiques]
\begin{itemize}
    \item Ne pas confondre \emph{Indépendance} (1er janv 1960, Cameroun français) et \emph{Réunification} (1er oct 1961, Sud britannique).
    \item La « guerre cachée » (maquis UPC vs armée française puis camerounaise) se poursuit après 1960 (jusqu'à mort Ouandié 1971).
    \item Le fédéralisme (1961) est un compromis pour la réunification ; l'unitarisme (1972) règle la question institutionnelle mais rouvre la « question anglophone ».
\end{itemize}
\end{attention}

\begin{itemize}
    \item \textbf{1972 :} Référendum 20 mai $\rightarrow$ \cle{République Unie du Cameroun}. Fin fédéralisme. Centralisation (Ahidjo).
    \item \textbf{1982--1984 :} Transition Biya. Tentative coup d'État 1984 (Garde républicaine). Renommage \cle{République du Cameroun} (1984).
    \item \textbf{Mémoire / Justice :} Réhabilitation Um Nyobé (1991), Moumié, Ouandié. Commission Vérité/Réconciliation (absente). Crise anglophone (2016--) : revendications fédéralisme / sécession (Ambazonie), racines dans gestion post-1972.
\end{itemize}

\end{bilan}

\begin{aretenir}[Checklist avant le Bac]
\begin{itemize}
    \item $\square$ Définir le statut de \cle{Tutelle ONU} (différence Mandat/Tutelle, rôle Conseil de tutelle).
    \item $\square$ Contraster l'administration du Cameroun français (TOM, Assemblée) et britannique (intégré Nigeria, indirect rule).
    \item $\square$ Citer 3 facteurs du nationalisme (contexte int., exploitation, syndicats, tirailleurs).
    \item $\square$ Identifier les 4 grands partis : \cle{UPC}, \cle{ANC}, \cle{KNDP}, \cle{ONE/CU} + leaders + idéologies.
    \item $\square$ Résumer le parcours des 3 leaders UPC : \cle{Um Nyobé} (Mpodol, ONU, maquis, mort 1958), \cle{Moumié} (exil, diplomatie, assassinat 1960), \cle{Ouandié} (chef maquis, exécution 1971).
    \item $\square$ Ordonner les 5 dates clés : 1955 (insurrection), 1956 (Loi-cadre), 1958 (Ahidjo ONU), 1960 (Indépendance), 1961 (Référendum + Réunification).
    \item $\square$ Expliquer le résultat du référendum 1961 (Nord $\rightarrow$ Nigeria, Sud $\rightarrow$ Réunification) et la Conférence de Foumban (Fédéralisme).
    \item $\square$ Analyser le passage Fédéralisme (1961) $\rightarrow$ Unitarisme (1972) : référendum 20 mai, motivations, conséquences.
    \item $\square$ Lier crise anglophone actuelle aux non-dits de la réunification (marginalisation, common law/éducation, fédéralisme).
    \item $\square$ Rédiger une intro/conclusion type : « De la tutelle à la souveraineté, le Cameroun construit son unité par la réunification (1961), mais l'héritage de la répression nationaliste (UPC) et le choix unitaire (1972) pèsent sur la cohésion nationale. »
    \item $\square$ Savoir analyser un document : discours Ahidjo à l'ONU (1958), rapport mission ONU, tract UPC, Constitution 1961/1972, photo référendum.
\end{itemize}
\end{aretenir}

\findecours

\end{document}
